% Statistiques à deux variables — Mathématiques Tle C — Cours signé OnBuch+
% Programme officiel MINESEC. Compiler avec : tectonic M10-statistiques-a-deux-variables.tex
\newcommand{\DOCMATIERE}{Mathématiques}
\newcommand{\DOCNIVEAU}{Tle C}
\newcommand{\DOCMODULE}{Organisation et gestion des données}
\newcommand{\DOCLECON}{10}
\newcommand{\DOCTITRE}{Statistiques à deux variables}
% OnBuch+ — préambule « cours signés » ENRICHI (compilé avec tectonic / XeLaTeX)
% Système visuel « Pop » d'OnBuch+ : fond crème (jamais blanc), bordures encre
% épaisses, ombres « dures » décalées, titres Archivo Black en capitales.
% Version MATHÉMATIQUES : graphes (pgfplots/tikz), boîtes pédagogiques
% (définition, propriété, méthode, attention, le savais-tu, exercice résolu,
% exercice, TP, bilan), page de garde et signature OnBuch+ ancrée en bas.
%
% Macros attendues dans main.tex AVANT \input{preamble} :
%   \DOCMATIERE  \DOCNIVEAU  \DOCMODULE  \DOCLECON (numéro)  \DOCTITRE
\documentclass[11pt]{article}

\usepackage{fontspec}
\usepackage[french]{babel}
\usepackage[a4paper,margin=2cm,top=3.4cm,bottom=2.8cm,headheight=1.4cm,footskip=1.3cm]{geometry}
\usepackage{amsmath,amssymb,amsfonts}
\usepackage{mathtools} % \xmapsto, \coloneqq... souvent utilisés par les modèles
\usepackage{cancel} % simplifications barrées
% \abs{z} (module, valeur absolue) et \norm{u} (norme), utilisés par les modèles
\providecommand{\abs}{}\let\abs\relax\DeclarePairedDelimiter{\abs}{\lvert}{\rvert}
\providecommand{\norm}{}\let\norm\relax\DeclarePairedDelimiter{\norm}{\lVert}{\rVert}
\usepackage[table]{xcolor} % [table] : fournit aussi \rowcolors (alternance de lignes)
\usepackage{pagecolor}
\usepackage{tikz}
\usetikzlibrary{arrows.meta,positioning,calc,shapes.geometric,shapes.misc,shapes.arrows,decorations.pathmorphing,decorations.markings,patterns,shadows,mindmap,trees,fit,backgrounds,matrix,intersections}
\usepackage{pgfplots}
\usepackage{pgf-pie} % diagrammes circulaires \pie (statistiques)
\pgfplotsset{compat=1.18}
\usepgfplotslibrary{fillbetween,groupplots} % aires entre courbes (calcul intégral)
\usepackage{enumitem}
\usepackage{fancyhdr}
% [most] charge toutes les bibliothèques tcolorbox (listings, minted, raster,
% documentation...) et gonfle inutilement la mémoire TeX sur un document long
% (~300 boîtes) : on ne charge que ce qui est réellement utilisé.
\usepackage[skins,breakable]{tcolorbox}
\usepackage{graphicx}
\usepackage{colortbl}
\usepackage{booktabs}
\usepackage{tabularx}
\usepackage{diagbox} % case d'en-tête diagonale des tableaux à double entrée
% Colonnes extensibles centrée / gauche / droite (C, L, R), souvent utilisées par les modèles
\newcolumntype{C}{>{\centering\arraybackslash}X}
\newcolumntype{L}{>{\raggedright\arraybackslash}X}
\newcolumntype{R}{>{\raggedleft\arraybackslash}X}
\usepackage{array}
\usepackage{multicol}
\usepackage{siunitx}
% parse-numbers=false : accepte aussi \SI{2,5 \times 10^{-3}}{...}
\sisetup{locale=FR,output-decimal-marker={,},group-minimum-digits=4,parse-numbers=false}
\DeclareSIUnit\ohm{\ensuremath{\symup{\Omega}}} % U+2126 absent de la police
\usepackage{qrcode}
% \degree : commande fréquemment utilisée par les modèles pour un angle ou une
% température en dehors de \SI{}{} (non fournie par les paquets chargés).
\providecommand{\degree}{^{\circ}}
% \qed (fin de démonstration), utilisé par les modèles sans amsthm
\providecommand{\qed}{\hfill$\square$}
% \barre : double barre des valeurs interdites dans un tableau de variations (tkz-tab)
\providecommand{\barre}{\ensuremath{\|}}
% environnement proof (amsthm non chargé) utilisé par les modèles
\ifdefined\proof\else\newenvironment{proof}[1][Démonstration]{\par\noindent\textit{#1.}\ }{\qed\par}\fi
\usepackage{lastpage}
\usepackage{hyperref}
\hypersetup{hidelinks}

% ============================================================
% Palette « Pop » OnBuch+ — mêmes hex que la classe P (plus_ui.dart)
% ============================================================
\definecolor{poporange}{HTML}{F59321}
\definecolor{popdark}{HTML}{E07A0C}
\definecolor{popcream}{HTML}{FFF6E8}
\definecolor{popink}{HTML}{1C1714}
\definecolor{poppurple}{HTML}{7A5AE0}
\definecolor{popgreen}{HTML}{2E9E6A}
\definecolor{poppink}{HTML}{E0457B}
\definecolor{popblue}{HTML}{2F7DE1}
\definecolor{popgold}{HTML}{C98A12}
\definecolor{popmuted}{HTML}{8A7F76}
\definecolor{popline}{HTML}{EADFD2}
\definecolor{popgray}{HTML}{8A7F76}
\colorlet{popgrey}{popgray}
% Alias tolérants (le modèle nomme parfois les couleurs différemment)
\colorlet{popwhite}{white}
\colorlet{popblack}{popink}
\colorlet{popred}{poppink}
\colorlet{popyellow}{popgold}
% Teintes claires pour fonds de tableaux / zones
\colorlet{popblueL}{popblue!10!white}
\colorlet{poporangeL}{poporange!14!white}
\colorlet{popgreenL}{popgreen!12!white}
\colorlet{poppurpleL}{poppurple!12!white}
\colorlet{poppinkL}{poppink!12!white}
\colorlet{popgoldL}{popgold!14!white}

\pagecolor{popcream}

% ============================================================
% Typographie
% ============================================================
\defaultfontfeatures{Ligatures=TeX}
\setmainfont{PlusJakartaSans-Medium.ttf}[Path=../../../fonts/,BoldFont=PlusJakartaSans-Bold.ttf]
% Maths en sans-serif (Fira Math) avec chiffres et lettres droites de Plus
% Jakarta, pour que les formules se fondent dans le texte.
\usepackage[math-style=ISO,bold-style=ISO]{unicode-math}
\setmathfont{FiraMath-Regular.otf}[Path=../../../fonts/]
\setmathfont{PlusJakartaSans-Medium.ttf}[Path=../../../fonts/,range={up/num,up/{Latin,latin},\mathup}]
\usepackage{icomma} % virgule décimale sans espace en mode math
% Caractères mathématiques Unicode absents des polices de texte (Plus Jakarta,
% Archivo Black) : rendus par leur équivalent mathématique, en texte comme en
% formule — sinon ils s'affichent en carré vide (ex. « Arithmétique dans ℤ »).
\usepackage{newunicodechar}
\newunicodechar{ℝ}{\ensuremath{\mathbb{R}}}
\newunicodechar{ℤ}{\ensuremath{\mathbb{Z}}}
\newunicodechar{ℂ}{\ensuremath{\mathbb{C}}}
\newunicodechar{ℕ}{\ensuremath{\mathbb{N}}}
\newunicodechar{ℚ}{\ensuremath{\mathbb{Q}}}
\newunicodechar{𝔻}{\ensuremath{\mathbb{D}}}
\newunicodechar{α}{\ensuremath{\alpha}}
\newunicodechar{β}{\ensuremath{\beta}}
\newunicodechar{γ}{\ensuremath{\gamma}}
\newunicodechar{δ}{\ensuremath{\delta}}
\newunicodechar{ε}{\ensuremath{\varepsilon}}
\newunicodechar{θ}{\ensuremath{\theta}}
\newunicodechar{λ}{\ensuremath{\lambda}}
\newunicodechar{μ}{\ensuremath{\mu}}
\newunicodechar{σ}{\ensuremath{\sigma}}
\newunicodechar{τ}{\ensuremath{\tau}}
\newunicodechar{φ}{\ensuremath{\varphi}}
\newunicodechar{ψ}{\ensuremath{\psi}}
\newunicodechar{ω}{\ensuremath{\omega}}
\newunicodechar{Σ}{\ensuremath{\Sigma}}
% majuscules produites par \MakeUppercase dans les titres (α-aminés -> Α)
\newunicodechar{Α}{\ensuremath{\alpha}}
\newunicodechar{Β}{\ensuremath{\beta}}
\newunicodechar{Γ}{\ensuremath{\Gamma}}
\newunicodechar{Δ}{\ensuremath{\Delta}}
\newunicodechar{Ε}{\ensuremath{\varepsilon}}
\newunicodechar{Θ}{\ensuremath{\Theta}}
\newunicodechar{Λ}{\ensuremath{\Lambda}}
\newunicodechar{Μ}{\ensuremath{\mu}}
\newunicodechar{Τ}{\ensuremath{\tau}}
\newunicodechar{Φ}{\ensuremath{\Phi}}
\newunicodechar{Ψ}{\ensuremath{\Psi}}
\newunicodechar{Ω}{\ensuremath{\Omega}}
\newunicodechar{↦}{\ensuremath{\mapsto}}
\newunicodechar{∈}{\ensuremath{\in}}
\newunicodechar{∉}{\ensuremath{\notin}}
\newunicodechar{∧}{\ensuremath{\wedge}}
\newunicodechar{⇔}{\ensuremath{\Leftrightarrow}}
\newunicodechar{⟺}{\ensuremath{\Longleftrightarrow}}
\newunicodechar{⇒}{\ensuremath{\Rightarrow}}
\newunicodechar{⟹}{\ensuremath{\Longrightarrow}}
\newunicodechar{∘}{\ensuremath{\circ}}
\newunicodechar{∪}{\ensuremath{\cup}}
\newunicodechar{∩}{\ensuremath{\cap}}
\newunicodechar{≡}{\ensuremath{\equiv}}
\newunicodechar{⊂}{\ensuremath{\subset}}
\newunicodechar{⊥}{\ensuremath{\perp}}
\newunicodechar{∃}{\ensuremath{\exists}}
\newunicodechar{∀}{\ensuremath{\forall}}
\newunicodechar{∛}{\ensuremath{\sqrt[3]{\,}}}
\newunicodechar{∜}{\ensuremath{\sqrt[4]{\,}}}
\newunicodechar{⃗}{\ensuremath{{}^{\to}}}
\newunicodechar{ⁿ}{\ifmmode^{n}\else\textsuperscript{\ensuremath{n}}\fi}
\newunicodechar{ˣ}{\ifmmode^{x}\else\textsuperscript{\ensuremath{x}}\fi}
\newunicodechar{ᵇ}{\ifmmode^{b}\else\textsuperscript{\ensuremath{b}}\fi}
\newunicodechar{ʳ}{\ifmmode^{r}\else\textsuperscript{\ensuremath{r}}\fi}
\newunicodechar{ᵘ}{\ifmmode^{u}\else\textsuperscript{\ensuremath{u}}\fi}
\newunicodechar{ʸ}{\ifmmode^{y}\else\textsuperscript{\ensuremath{y}}\fi}
\newunicodechar{ᵃ}{\ifmmode^{a}\else\textsuperscript{\ensuremath{a}}\fi}
\newunicodechar{ᵉ}{\ifmmode^{e}\else\textsuperscript{\ensuremath{e}}\fi}
\newunicodechar{ᵏ}{\ifmmode^{k}\else\textsuperscript{\ensuremath{k}}\fi}
\newunicodechar{ᵖ}{\ifmmode^{p}\else\textsuperscript{\ensuremath{p}}\fi}
\newunicodechar{ᴺ}{\ifmmode^{N}\else\textsuperscript{\ensuremath{N}}\fi}
\newunicodechar{ᶜ}{\ifmmode^{c}\else\textsuperscript{\ensuremath{c}}\fi}
\newunicodechar{ʰ}{\ifmmode^{h}\else\textsuperscript{\ensuremath{h}}\fi}
\newunicodechar{ᵠ}{\ifmmode^{q}\else\textsuperscript{\ensuremath{q}}\fi}
\newunicodechar{ₙ}{\ifmmode_{n}\else\textsubscript{\ensuremath{n}}\fi}
\newunicodechar{ₐ}{\ifmmode_{a}\else\textsubscript{\ensuremath{a}}\fi}
\newunicodechar{ᵢ}{\ifmmode_{i}\else\textsubscript{\ensuremath{i}}\fi}
\newunicodechar{ᵣ}{\ifmmode_{r}\else\textsubscript{\ensuremath{r}}\fi}
\newunicodechar{ₖ}{\ifmmode_{k}\else\textsubscript{\ensuremath{k}}\fi}
\newunicodechar{ᵦ}{\ifmmode_{\beta}\else\textsubscript{\ensuremath{\beta}}\fi}
\newunicodechar{ₓ}{\ifmmode_{x}\else\textsubscript{\ensuremath{x}}\fi}
\newfontfamily\popdisplay{ArchivoBlack-Regular.ttf}[Path=../../../fonts/]
\newfontfamily\poplogo{SpaceGrotesk-Bold.ttf}[Path=../../../fonts/]
\newfontfamily\popsemi{PlusJakartaSans-SemiBold.ttf}[Path=../../../fonts/]
\newfontfamily\popxbold{PlusJakartaSans-ExtraBold.ttf}[Path=../../../fonts/]
\newfontfamily\popxboldls{PlusJakartaSans-ExtraBold.ttf}[Path=../../../fonts/,LetterSpace=3.0]

\setlist[enumerate]{leftmargin=1.7em,itemsep=5pt,topsep=4pt}
\setlist[itemize]{leftmargin=1.7em,itemsep=4pt,topsep=4pt,label={\color{poporange}\textbullet}}
\setlength{\parindent}{0pt}
\setlength{\parskip}{5pt}
\renewcommand{\arraystretch}{1.25}

% pgfplots : style Pop par défaut
\pgfplotsset{
  every axis/.append style={axis line style={popink,line width=1pt},tick style={popink},
    grid style={popline,line width=0.5pt},label style={font=\small},tick label style={font=\footnotesize}},
  popcurve/.style={poporange,line width=1.8pt,no markers,smooth},
}
% Même style disponible dans un \draw TikZ simple (hors pgfplots)
\tikzset{popcurve/.style={poporange,line width=1.8pt}}

% ============================================================
% Pill / squiggle
% ============================================================
\newcommand{\poppill}[3]{%
  \tikz[baseline=-0.55ex]\node[draw=popink,line width=1.2pt,rounded corners=2.6pt,%
    fill=#2,inner xsep=6.5pt,inner ysep=3pt]{%
    \popxboldls\fontsize{7}{7}\selectfont\color{#3}\MakeUppercase{#1}};%
}
\newcommand{\popsquiggle}[1]{%
  \tikz[baseline=-0.4ex]{
    \draw[#1,line width=2.6pt,line cap=round]
      plot [smooth] coordinates {(0,0.09) (0.34,0.01) (0.68,0.21) (1.02,0.01) (1.36,0.10)};
  }%
}
% Mot-clé surligné (marqueur orange sous le texte)
\newcommand{\cle}[1]{{\popxbold\color{popdark}#1}}

% ============================================================
% En-tête / pied
% ============================================================
\pagestyle{fancy}
\fancyhf{}
\renewcommand{\headrulewidth}{0pt}
\renewcommand{\footrulewidth}{0pt}
\lhead{\raisebox{-0.15ex}{\poplogo\fontsize{13}{13}\selectfont\color{popink}OnBuch\color{poporange}+}}
\rhead{\poppill{\DOCMATIERE\ \textbullet\ \DOCNIVEAU\ \textbullet\ Leçon \DOCLECON}{white}{popink}}
\lfoot{\popxbold\fontsize{7.6}{7.6}\selectfont\color{popmuted}\MakeUppercase{Cours signé OnBuch+}\ \textbullet\ {\popsemi\DOCTITRE}}
\rfoot{\tikz[baseline=-0.6ex]\node[rounded corners=8.5pt,fill=popink,minimum height=17pt,minimum width=17pt,inner xsep=5pt,inner sep=0]{\popxbold\fontsize{8}{8}\selectfont\color{popcream}\thepage\,/\,\pageref*{LastPage}};}

% ============================================================
% Titres
% ============================================================
\newcommand{\chaptitre}[2]{%
  \begin{tcolorbox}[enhanced,colback=poporange,colframe=popink,boxrule=2.6pt,arc=11pt,
      drop shadow=popink,left=15pt,right=15pt,top=12pt,bottom=11pt,before skip=3pt,after skip=18pt]
    \poppill{#2}{popink}{popcream}\par\vspace{9pt}
    {\raggedright\hyphenpenalty=10000\exhyphenpenalty=10000\popdisplay\fontsize{22}{24}\selectfont\color{popcream}\MakeUppercase{#1}\par}
  \end{tcolorbox}
}
% Section numérotée : pastille orange avec le numéro + titre Archivo
\newcounter{popsec}
\newcommand{\coursec}[1]{%
  \stepcounter{popsec}\setcounter{popsub}{0}%
  \par\vspace{16pt}\noindent
  \begin{minipage}{\linewidth}
  \tikz[baseline=-0.7ex]\node[draw=popink,line width=1.6pt,fill=poporange,rounded corners=4pt,minimum size=22pt,inner sep=0]{\popdisplay\fontsize{12}{12}\selectfont\color{popcream}\thepopsec};%
  \hspace{8pt}{\popdisplay\fontsize{15}{17}\selectfont\color{popink}\MakeUppercase{#1}}\par
  \vspace{4pt}\noindent{\color{poporange}\rule{36pt}{3.6pt}}
  \end{minipage}\par\nopagebreak\vspace{8pt}
}
\newcounter{popsub}
\newcommand{\courssub}[1]{%
  \stepcounter{popsub}%
  \par\vspace{9pt}\noindent{\popxbold\fontsize{11.5}{13}\selectfont\color{popink}{\color{poporange}\thepopsec.\thepopsub}\hspace{6pt}#1}\par\nopagebreak\vspace{4pt}
}

% ============================================================
% Boîtes pédagogiques — même recette : fond blanc, bordure épaisse colorée,
% ombre dure encre, badge plein. Titre optionnel [#1].
% ============================================================
\tcbset{popbase/.style={breakable,enhanced,colback=white,boxrule=2.3pt,arc=9pt,
  shadow={3pt}{-3pt}{0pt}{popink},left=14pt,right=14pt,top=11pt,bottom=11pt,before skip=11pt,after skip=15pt,
  fontupper=\popsemi\fontsize{10.6}{14.5}\selectfont\color{popink},
  fontlower=\popsemi\fontsize{10.6}{14.5}\selectfont\color{popink}}}
% Coche dessinée (le glyphe ✓ n'existe pas dans les polices du document)
\newcommand{\popcheck}[1]{\tikz[baseline=-0.5ex]\draw[#1,line width=1.6pt,line cap=round,line join=round](-0.13,0.0)--(-0.04,-0.09)--(0.13,0.1);}
\newcommand{\popboxhead}[3]{\poppill{#1}{#2}{white}\ifx\relax#3\relax\else\hspace{6pt}{\popxbold\fontsize{10.6}{12}\selectfont\color{popink}#3}\fi\par\vspace{7pt}}

\newtcolorbox{aretenir}[1][]{popbase,colframe=popgreen,before upper={\popboxhead{À retenir}{popgreen}{#1}}}
\newtcolorbox{exemplebox}[1][]{popbase,colframe=poporange,before upper={\popboxhead{Exemple}{poporange}{#1}}}
\newtcolorbox{definition}[1][]{popbase,colframe=popblue,before upper={\popboxhead{Définition}{popblue}{#1}}}
\newtcolorbox{propriete}[1][]{popbase,colframe=poppurple,before upper={\popboxhead{Propriété}{poppurple}{#1}}}
\newtcolorbox{methode}[1][]{popbase,colframe=popgold,before upper={\popboxhead{Méthode}{popgold}{#1}}}
\newtcolorbox{attention}[1][]{popbase,colframe=poppink,before upper={\popboxhead{Attention}{poppink}{#1}}}
\newtcolorbox{experience}[1][]{popbase,colframe=popblue,colback=popblueL,before upper={\popboxhead{Expérience}{popblue}{#1}}}
% « Le savais-tu ? » : carte violette pleine (contexte, culture, Cameroun)
\newtcolorbox{savaistu}[1][]{popbase,colback=poppurple,colframe=popink,
  fontupper=\popsemi\fontsize{10.4}{14.2}\selectfont\color{white},
  before upper={\poppill{Le savais-tu\,?}{white}{poppurple}\ifx\relax#1\relax\else\hspace{6pt}{\popxbold\color{white}#1}\fi\par\vspace{7pt}}}
% Exercice résolu : énoncé en haut, \tcblower puis la solution sur fond crème
\newcounter{popexo}\newcounter{popexr}
\newtcolorbox{exoresolu}[1][]{popbase,colframe=poporange,colbacklower=popcream,
  segmentation style={popink,line width=1.2pt,dashed},
  before upper={\stepcounter{popexr}\popboxhead{Exercice résolu \thepopexr}{poporange}{#1}},
  before lower={\poppill{Solution}{popink}{popcream}\par\vspace{6pt}}}
% Exercice (non résolu en place) — la correction est dans la partie Corrigés
\newtcolorbox{exercice}[1][]{popbase,colframe=popink,
  before upper={\stepcounter{popexo}\popboxhead{Exercice \thepopexo}{popink}{#1}}}
% Corrigé
\newtcolorbox{corrige}[1][]{popbase,colframe=popgreen,colback=popgreenL,
  before upper={\popboxhead{Corrigé}{popgreen}{#1}}}
% Objectifs / prérequis / bilan
\newtcolorbox{objectifs}{popbase,colframe=popink,colback=poporangeL,
  before upper={\popboxhead{Objectifs de la leçon}{popink}{}}}
\newtcolorbox{prerequis}{popbase,colframe=popink,colback=popblueL,
  before upper={\popboxhead{Prérequis}{popblue}{}}}
\newtcolorbox{bilan}{popbase,colframe=popink,colback=white,boxrule=2.8pt,
  before upper={\popboxhead{Fiche bilan}{poporange}{L'essentiel à connaître pour l'examen}}}
% Figure encadrée avec légende
\newtcolorbox{popfigure}[1][]{enhanced,breakable=false,colback=white,colframe=popline,boxrule=1.4pt,arc=8pt,
  left=10pt,right=10pt,top=10pt,bottom=8pt,before skip=10pt,after skip=12pt,halign=center,
  lower separated=false,halign lower=center,
  fontlower=\popsemi\fontsize{9}{11.5}\selectfont\color{popmuted},
  #1}
\newcommand{\legende}[1]{\par\vspace{6pt}{\popsemi\fontsize{9}{11.5}\selectfont\color{popmuted}#1}}

% ============================================================
% Signature OnBuch+
% ============================================================
\newcommand{\poplogoinline}[1]{{\poplogo\fontsize{#1}{#1}\selectfont\color{popink}OnBuch\color{poporange}+}}
\newcommand{\signatureonbuch}{%
  \begin{tcolorbox}[enhanced,colback=popink,colframe=popink,boxrule=2.2pt,arc=10pt,
      drop shadow=poporange,left=14pt,right=14pt,top=10pt,bottom=10pt,before skip=0pt,after skip=0pt]
    \tikz[baseline=-0.7ex]\node[circle,fill=popgreen,draw=popcream,line width=1.3pt,minimum size=22pt,inner sep=0]{\popcheck{white}};%
    \hspace{9pt}\begin{minipage}[c]{0.62\linewidth}
      {\popxbold\fontsize{12}{13}\selectfont\color{popcream}Signé\ \textbullet\ {\poplogo OnBuch\color{poporange}+}}\par\vspace{2pt}
      {\raggedright\popsemi\fontsize{8}{10.5}\selectfont\color{popline}\MakeUppercase{Équipe pédagogique OnBuch+}\\\MakeUppercase{Conforme au programme officiel MINESEC}\par}
    \end{minipage}\hfill
    {\poplogo\fontsize{22}{22}\selectfont\color{popcream}OnBuch\color{poporange}+}
  \end{tcolorbox}%
}

% ============================================================
% Page de garde
% #1 titre · #2 matière · #3 niveau · #4 module · #5 n° leçon · #6 durée ·
% #7 description / objectifs (texte ou itemize)
% ============================================================
\newcommand{\pagegarde}[7]{%
  \thispagestyle{empty}
  \newgeometry{a4paper,margin=1.7cm,top=1.6cm,bottom=1.4cm}%
  \begin{tikzpicture}[remember picture,overlay]
    \fill[poporange!18] ([xshift=-3.2cm,yshift=-3.6cm]current page.north east) circle (4.6cm);
    \fill[poppurple!14] ([xshift=2.4cm,yshift=7.6cm]current page.south west) circle (3.8cm);
  \end{tikzpicture}%
  {\poplogo\fontsize{30}{30}\selectfont\color{popink}OnBuch\color{poporange}+}\hfill
  \raisebox{0.6ex}{\poppill{Cours signé \textbullet\ Premium}{popink}{popcream}}\par
  \vspace{1pt}\popsquiggle{poppurple}\par
  \vspace{8pt}
  \begin{tcolorbox}[enhanced,colback=poporange,colframe=popink,boxrule=2.8pt,arc=13pt,
      drop shadow=popink,left=18pt,right=18pt,top=12pt,bottom=13pt,after skip=10pt]
    \poppill{#2 \textbullet\ #3}{popink}{popcream}\hspace{5pt}\poppill{#4}{white}{popink}\par\vspace{8pt}
    {\popdisplay\fontsize{14}{14}\selectfont\color{popink}LEÇON #5}\par\vspace{4pt}
    {\raggedright\hyphenpenalty=10000\exhyphenpenalty=10000\popdisplay\fontsize{25}{27}\selectfont\color{popcream}\MakeUppercase{#1}\par}
    \vspace{8pt}
    {\popxbold\fontsize{9.5}{9.5}\selectfont\color{popink}\MakeUppercase{Durée conseillée : #6}}
  \end{tcolorbox}
  \begin{tcolorbox}[enhanced,colback=white,colframe=popink,boxrule=2.2pt,arc=10pt,
      drop shadow=popink,left=16pt,right=16pt,top=10pt,bottom=10pt,after skip=8pt]
    \poppill{Dans cette leçon}{poppurple}{white}\par\vspace{5pt}
    \popsemi\fontsize{9.3}{12.6}\selectfont\color{popink}#7
  \end{tcolorbox}
  \noindent\begin{minipage}[t]{0.487\linewidth}
    \begin{tcolorbox}[enhanced,colback=popgreen,colframe=popink,boxrule=2.2pt,arc=10pt,
        drop shadow=popink,left=11pt,right=11pt,top=10pt,bottom=10pt,height=3.1cm,valign=center]
      \tcbox[colback=white,colframe=popink,boxrule=1.3pt,arc=5pt,boxsep=0pt,nobeforeafter,
        left=3pt,right=3pt,top=3pt,bottom=3pt,valign=center]{\qrcode[height=1.9cm]{https://play.google.com/store/apps/details?id=com.luvvix.onbuch&hl=fr}}%
      \hspace{8pt}\begin{minipage}[c]{0.5\linewidth}
        {\popdisplay\fontsize{11}{12.5}\selectfont\color{white}\MakeUppercase{Télécharge OnBuch}}\par\vspace{4pt}
        {\raggedright\popsemi\fontsize{8.4}{11}\selectfont\color{white}Gratuit \textbullet\ Google Play\\Tuteur IA, annales, concours, résultats\par}
      \end{minipage}
    \end{tcolorbox}
  \end{minipage}\hfill
  \begin{minipage}[t]{0.487\linewidth}
    \begin{tcolorbox}[enhanced,colback=poppurple,colframe=popink,boxrule=2.2pt,arc=10pt,
        drop shadow=popink,left=11pt,right=11pt,top=10pt,bottom=10pt,height=3.1cm,valign=center]
      \tikz[baseline=-0.7ex]\node[circle,fill=white,draw=popink,line width=1.3pt,minimum size=1.6cm,inner sep=0]{\popdisplay\fontsize{24}{24}\selectfont\color{poppurple}+};%
      \hspace{8pt}\begin{minipage}[c]{0.6\linewidth}
        {\popdisplay\fontsize{11}{12.5}\selectfont\color{white}\MakeUppercase{Passe à OnBuch+}}\par\vspace{4pt}
        {\raggedright\popsemi\fontsize{8.4}{11}\selectfont\color{white}Tous les cours signés, Tuteur IA illimité, fiches PDF — onglet OnBuch+ de l'app\par}
      \end{minipage}
    \end{tcolorbox}
  \end{minipage}
  \vspace{8pt}
  \popcontenu
  \vfill
  \signatureonbuch
  \restoregeometry
  \newpage
}

% Rangée « Ce cours contient » (page de garde)
\newcommand{\poptile}[3]{% #1 couleur #2 gros texte #3 légende
  \begin{tcolorbox}[enhanced,colback=#1,colframe=popink,boxrule=2pt,arc=9pt,shadow={2.5pt}{-2.5pt}{0pt}{popink},
      left=6pt,right=6pt,top=9pt,bottom=9pt,halign=center,valign=center,height=1.75cm,width=\linewidth,nobeforeafter]
    {\popdisplay\fontsize{12.5}{13}\selectfont\color{white}\MakeUppercase{#2}}\par\vspace{3pt}
    {\centering\popsemi\fontsize{7.8}{9.5}\selectfont\color{white}#3\par}
  \end{tcolorbox}}
\newcommand{\popcontenu}{%
  \par\noindent{\popxboldls\fontsize{7.5}{7.5}\selectfont\color{popmuted}CE COURS SIGNÉ CONTIENT}\par\vspace{7pt}
  \noindent\begin{minipage}[t]{0.235\linewidth}\poptile{popblue}{Cours}{complet, illustré, pas à pas}\end{minipage}\hfill
  \begin{minipage}[t]{0.235\linewidth}\poptile{poporange}{Méthodes}{et exercices résolus}\end{minipage}\hfill
  \begin{minipage}[t]{0.235\linewidth}\poptile{poppink}{Exercices}{type examen + corrigés}\end{minipage}\hfill
  \begin{minipage}[t]{0.235\linewidth}\poptile{popgreen}{Fiche bilan}{l'essentiel à retenir}\end{minipage}\par}

% Sommaire
\newtcolorbox{sommaire}{enhanced,colback=white,colframe=popink,boxrule=2.2pt,arc=10pt,
  drop shadow=popink,left=18pt,right=18pt,top=13pt,bottom=13pt,before skip=6pt,after skip=20pt,
  before upper={\poppill{Sommaire}{poppurple}{white}\par\vspace{9pt}\popsemi\fontsize{10.6}{14.5}\selectfont\color{popink}}}

% Signature de fin de cours — poussée tout en bas de la dernière page
\newcommand{\findecours}{%
  \par\vfill\nopagebreak
  \begin{center}\popsquiggle{poporange}\end{center}\vspace{4pt}
  \signatureonbuch
}
\AtBeginDocument{\def\llbracket{\mathopen{[\mkern-3.5mu[}}\def\rrbracket{\mathclose{]\mkern-3.5mu]}}} % intervalles d'entiers (glyphe ⟦ absent de la police)
% Glyphes absents de Fira Math : redéfinis à partir de glyphes disponibles
\AtBeginDocument{%
  \def\setminus{\mathbin{\backslash}}\let\smallsetminus\setminus
  \def\vdots{\mathord{\vbox{\baselineskip3.2pt\lineskiplimit0pt\kern4pt\hbox{.}\hbox{.}\hbox{.}}}}%
  \def\ddots{\mathinner{\mkern1mu\raise6.5pt\hbox{.}\mkern2mu\raise3.6pt\hbox{.}\mkern2mu\raise0.7pt\hbox{.}\mkern1mu}}%
  \def\triangle{\mathord{\scalebox{0.9}{$\symup{\Delta}$}}}%
  \def\nabla{\mathord{\raisebox{\depth}{\scalebox{1}[-1]{$\symup{\Delta}$}}}}%
  \def\checkmark{\popcheck{popdark}}%
  \def\star{\mathbin{\ast}}\def\bigstar{\mathord{\ast}}%
  \def\diamond{\mathbin{\scalebox{0.6}{$\Diamond$}}}%
}

%% FIGFIX-BEGIN v5 — correctifs globaux des figures (docs/onbuch-generation/tools/figfix)
\usepackage{adjustbox}\usepackage{etoolbox}\tcbuselibrary{raster}
% toute figure TikZ/pgfplots plus large que la ligne est réduite à la largeur disponible
\BeforeBeginEnvironment{tikzpicture}{\begin{adjustbox}{max width=\linewidth}}
\AfterEndEnvironment{tikzpicture}{\end{adjustbox}}
% légendes pgfplots lisibles par-dessus les courbes
\pgfplotsset{every axis/.append style={legend style={fill=white,fill opacity=0.92,text opacity=1,draw=black!15,inner sep=3pt}}}
% étiquettes d'arêtes (syntaxe edge["..."]) avec fond blanc
\tikzset{every edge quotes/.append style={fill=white,inner sep=1.2pt}}
%% FIGFIX-END
\begin{document}

\pagegarde{Statistiques à deux variables}{Mathématiques}{Tle C}{Organisation et gestion des données}{10}{9 h}{Cette leçon introduit l'étude simultanée de deux caractères statistiques : représentation par nuage de points, ajustement affine (méthodes de Mayer et des moindres carrés), mesure de la corrélation et prévisions. Elle constitue un chapitre central de l'épreuve de mathématiques au Baccalauréat C, D, E et TI.
\vspace{4pt}
\begin{multicols}{2}\small
\begin{itemize}
\item À la fin de cette leçon, je sais construire et exploiter un tableau à double entrée et en déduire les séries marginales
\item À la fin de cette leçon, je sais construire un nuage de points et placer le point moyen G
\item À la fin de cette leçon, je sais déterminer une droite d'ajustement affine par la méthode de Mayer
\item À la fin de cette leçon, je sais calculer la covariance de deux variables statistiques
\item À la fin de cette leçon, je sais déterminer les droites de régression de y en x et de x en y par les moindres carrés
\item À la fin de cette leçon, je sais calculer et interpréter le coefficient de corrélation linéaire
\end{itemize}\end{multicols}}

\begin{sommaire}
\begin{multicols}{2}
\begin{enumerate}
\item Séries statistiques à deux caractères
\item Nuage de points et point moyen
\item Ajustement affine par la méthode de Mayer
\item Covariance et moindres carrés
\item Coefficient de corrélation linéaire
\item Prévisions et synthèse méthodologique
\item Activité d'intégration
\item Méthodes et exercices résolus
\item Exercices
\item Corrigés des exercices
\item Fiche bilan
\end{enumerate}
\end{multicols}
\end{sommaire}

\begin{objectifs}
\begin{itemize}
\item À la fin de cette leçon, je sais construire et exploiter un tableau à double entrée et en déduire les séries marginales
\item À la fin de cette leçon, je sais construire un nuage de points et placer le point moyen G
\item À la fin de cette leçon, je sais déterminer une droite d'ajustement affine par la méthode de Mayer
\item À la fin de cette leçon, je sais calculer la covariance de deux variables statistiques
\item À la fin de cette leçon, je sais déterminer les droites de régression de y en x et de x en y par les moindres carrés
\item À la fin de cette leçon, je sais calculer et interpréter le coefficient de corrélation linéaire
\item À la fin de cette leçon, je sais juger de la qualité d'un ajustement affine
\item À la fin de cette leçon, je sais utiliser un ajustement affine pour effectuer une prévision et en discuter la fiabilité
\end{itemize}
\end{objectifs}

\begin{prerequis}
\begin{itemize}
\item Statistiques à une variable (Première) : moyenne, variance, écart-type, formules avec effectifs
\item Équations de droites dans le plan : pente (coefficient directeur), point d'appartenance
\item Repérage dans le plan et lecture graphique
\item Résolution de systèmes de deux équations à deux inconnues
\item Calculs avec le symbole somme Σ et manipulation de moyennes pondérées
\end{itemize}
\end{prerequis}

\newpage

\coursec{Séries statistiques à deux caractères}

\textbf{Situation d'accroche.} Mme Etoundi gère une station-service sur l'axe lourd Yaoundé--Soa. Chaque soir, elle note dans son cahier deux nombres : le nombre de véhicules comptés devant sa station dans la journée, et le volume de carburant vendu (en litres). Au bout de huit semaines, elle observe une chose frappante : les jours de fort trafic, les ventes grimpent ; les jours calmes, elles s'effondrent. Deux questions la taraudent alors : \emph{peut-on prévoir les ventes de la semaine prochaine à partir du trafic annoncé ?} Et surtout : \emph{comment mesurer la force du lien entre ces deux quantités ?} Dire « ça augmente ensemble » ne suffit pas à un banquier pour accorder un crédit de stock : il faut un nombre, un coefficient, une méthode. C'est précisément l'objet de cette leçon : la \cle{statistique à deux variables}. Avant de prévoir (sections 3 à 6), il faut d'abord apprendre à \emph{organiser} ces données doubles : c'est l'objet de cette première section.

\courssub{Qu'est-ce qu'une série statistique à deux caractères ?}

\textbf{Intuition.} En classe de Première, tu as étudié des séries statistiques à \emph{un} caractère : la taille des élèves d'une classe, les notes à un devoir. Mais dans la vie réelle, on s'intéresse souvent à \emph{deux} grandeurs mesurées \emph{sur les mêmes individus} : le poids et la taille d'un nouveau-né à l'hôpital de Bafoussam, le revenu et la dépense téléphonique d'un ménage, l'âge d'un véhicule et sa consommation de carburant. La question n'est plus « comment se répartit telle grandeur ? » mais « \emph{ces deux grandeurs sont-elles liées, et comment ?} ».

\begin{definition}[Série statistique à deux caractères]
Soit une population de $N$ individus. On observe sur \emph{chaque} individu \emph{deux} caractères quantitatifs $X$ et $Y$. L'ensemble des couples $(x_i\,;\,y_i)$ obtenus s'appelle une \cle{série statistique à deux caractères} (ou série double, ou série bivariée). L'entier $N$ est l'\cle{effectif total} de la série.
\end{definition}

Deux présentations des données sont possibles, selon la taille de la population.

\textbf{Première présentation : la liste de couples.} Quand l'effectif est modeste (comme les huit semaines de Mme Etoundi), on donne simplement la liste des couples $(x_i\,;\,y_i)$, souvent sous forme d'un tableau à deux lignes :

\begin{center}
\begin{tabularx}{\linewidth}{|l|X|X|X|X|X|X|X|X|}
\hline
\rowcolor{popblueL} Semaine $i$ & 1 & 2 & 3 & 4 & 5 & 6 & 7 & 8 \\
\hline
Trafic $x_i$ (centaines de véhicules) & 12 & 15 & 9 & 18 & 14 & 20 & 11 & 16 \\
\hline
Ventes $y_i$ (centaines de litres) & 30 & 38 & 25 & 45 & 36 & 50 & 28 & 41 \\
\hline
\end{tabularx}
\end{center}

Chaque colonne est un couple $(x_i\,;\,y_i)$ : par exemple, la semaine 4 correspond au couple $(18\,;\,45)$.

\textbf{Deuxième présentation : le tableau à double entrée.} Quand l'effectif est grand et que les valeurs se répètent (ou sont regroupées en classes), on croise les valeurs de $X$ (en lignes) et celles de $Y$ (en colonnes) dans un tableau.

\begin{definition}[Tableau à double entrée]
On note $x_1, x_2, \dots, x_p$ les valeurs (ou classes) prises par $X$ et $y_1, y_2, \dots, y_q$ celles prises par $Y$. Le \cle{tableau à double entrée} donne, à l'intersection de la ligne $i$ et de la colonne $j$, l'\cle{effectif conjoint} $n_{ij}$ : le nombre d'individus qui présentent \emph{à la fois} la valeur $x_i$ de $X$ \emph{et} la valeur $y_j$ de $Y$.
\end{definition}

\begin{attention}[Effectif conjoint $\neq$ effectifs marginaux]
L'effectif $n_{ij}$ compte les individus vérifiant \emph{simultanément} les deux conditions ($X = x_i$ ET $Y = y_j$). Ne le confonds jamais avec $n_{i\bullet}$ (individus vérifiant $X = x_i$, quelle que soit la valeur de $Y$) ni avec $n_{\bullet j}$ (individus vérifiant $Y = y_j$, quelle que soit la valeur de $X$). Au Baccalauréat, cette confusion est la première cause de perte de points dans les tableaux à double entrée : vérifie toujours que la somme de tous les $n_{ij}$ redonne l'effectif total $N$.
\end{attention}

\courssub{Un exemple complet : téléphonie mobile et revenu à Douala}

\begin{exemplebox}[Dépenses de téléphonie mobile selon le revenu]
Une enquête porte sur $N = 50$ ménages d'un quartier de Douala. On relève pour chaque ménage son revenu mensuel $X$ (en dizaines de milliers de FCFA, regroupé en classes) et sa dépense mensuelle de téléphonie mobile $Y$ (en milliers de FCFA, regroupée en classes). Les résultats sont consignés dans le tableau à double entrée suivant :
\begin{center}
\begin{tabularx}{\linewidth}{|l|X|X|X|X|}
\hline
\rowcolor{popblueL} $X$ \textbackslash{} $Y$ & $[2\,;\,5[$ & $[5\,;\,8[$ & $[8\,;\,11[$ & Total ligne \\
\hline
$[10\,;\,20[$ & 8 & 4 & 1 & 13 \\
\hline
$[20\,;\,30[$ & 5 & 9 & 3 & 17 \\
\hline
$[30\,;\,40[$ & 2 & 6 & 5 & 13 \\
\hline
$[40\,;\,50[$ & 0 & 3 & 4 & 7 \\
\hline
Total colonne & 15 & 22 & 13 & 50 \\
\hline
\end{tabularx}
\end{center}
\end{exemplebox}

\textbf{Lecture du tableau.} La case située à l'intersection de la ligne $[20\,;\,30[$ et de la colonne $[5\,;\,8[$ contient $n_{22} = 9$ : cela signifie que \emph{9 ménages} ont un revenu compris entre \num{200000} et \num{300000} FCFA \emph{et} une dépense téléphonique comprise entre \num{5000} et \num{8000} FCFA. De même, $n_{43} = 4$ : quatre ménages gagnent entre \num{400000} et \num{500000} FCFA et dépensent entre \num{8000} et \num{11000} FCFA en téléphonie.

\textbf{Vérification de l'effectif total.} Additionnons tous les effectifs conjoints :
\[ \sum_{i=1}^{4}\sum_{j=1}^{3} n_{ij} = (8+4+1) + (5+9+3) + (2+6+5) + (0+3+4) = 13+17+13+7 = 50. \]
On retrouve bien $N = 50$ : le tableau est cohérent. \emph{Cette vérification doit devenir un réflexe systématique.}

\courssub{Séries marginales et moyennes marginales}

\textbf{Intuition.} Le tableau à double entrée contient \emph{trop} d'information si l'on veut étudier un seul caractère à la fois. Pour retrouver la série « classique » du caractère $X$ seul, il suffit d'« écraser » le tableau horizontalement : pour chaque ligne, on additionne tous les effectifs de la ligne. On obtient les \emph{effectifs marginaux}, placés dans les marges du tableau — d'où leur nom.

\begin{definition}[Effectifs marginaux et séries marginales]
\begin{itemize}
\item L'\cle{effectif marginal} de la valeur $x_i$ est la somme des effectifs de la ligne $i$ :
\[ n_{i\bullet} = \sum_{j=1}^{q} n_{ij} = n_{i1} + n_{i2} + \dots + n_{iq}. \]
La série $(x_i\,;\,n_{i\bullet})$ est la \cle{série marginale de $X$}.
\item L'\cle{effectif marginal} de la valeur $y_j$ est la somme des effectifs de la colonne $j$ :
\[ n_{\bullet j} = \sum_{i=1}^{p} n_{ij} = n_{1j} + n_{2j} + \dots + n_{pj}. \]
La série $(y_j\,;\,n_{\bullet j})$ est la \cle{série marginale de $Y$}.
\end{itemize}
On a toujours : $\displaystyle \sum_{i=1}^{p} n_{i\bullet} = \sum_{j=1}^{q} n_{\bullet j} = \sum_{i,j} n_{ij} = N$.
\end{definition}

\begin{definition}[Moyennes marginales]
Les \cle{moyennes marginales} sont les moyennes des séries marginales :
\[ \bar{x} = \frac{1}{N}\sum_{i=1}^{p} n_{i\bullet}\, x_i \qquad \text{et} \qquad \bar{y} = \frac{1}{N}\sum_{j=1}^{q} n_{\bullet j}\, y_j. \]
\end{definition}

\begin{methode}[Extraire les séries marginales d'un tableau à double entrée]
\begin{enumerate}
\item \textbf{Calculer les totaux de lignes} : pour chaque ligne $i$, additionner les $n_{ij}$ ; écrire le résultat $n_{i\bullet}$ dans une colonne « Total » à droite.
\item \textbf{Calculer les totaux de colonnes} : pour chaque colonne $j$, additionner les $n_{ij}$ ; écrire le résultat $n_{\bullet j}$ dans une ligne « Total » en bas.
\item \textbf{Contrôler} : la somme des totaux de lignes et la somme des totaux de colonnes doivent toutes deux valoir $N$.
\item \textbf{Écrire les séries marginales} sous forme de tableaux à une variable : $(x_i\,;\,n_{i\bullet})$ pour $X$ et $(y_j\,;\,n_{\bullet j})$ pour $Y$.
\item \textbf{Calculer les moyennes marginales} $\bar{x}$ et $\bar{y}$ avec les formules habituelles de la moyenne pondérée (en utilisant les \emph{centres de classes} si les données sont regroupées).
\end{enumerate}
\end{methode}

\begin{exoresolu}[Séries marginales de l'enquête de Douala]
À partir du tableau des 50 ménages de Douala : (1) déterminer les deux séries marginales ; (2) calculer les moyennes marginales $\bar{x}$ et $\bar{y}$.
\tcblower
\textbf{1. Séries marginales.} Les totaux de lignes donnent la série marginale de $X$ :
\begin{center}
\begin{tabularx}{0.8\linewidth}{|l|X|X|X|X|}
\hline
\rowcolor{popblueL} Classe de $X$ & $[10;20[$ & $[20;30[$ & $[30;40[$ & $[40;50[$ \\
\hline
Effectif $n_{i\bullet}$ & 13 & 17 & 13 & 7 \\
\hline
\end{tabularx}
\end{center}
Les totaux de colonnes donnent la série marginale de $Y$ :
\begin{center}
\begin{tabularx}{0.7\linewidth}{|l|X|X|X|}
\hline
\rowcolor{popblueL} Classe de $Y$ & $[2;5[$ & $[5;8[$ & $[8;11[$ \\
\hline
Effectif $n_{\bullet j}$ & 15 & 22 & 13 \\
\hline
\end{tabularx}
\end{center}
Contrôle : $13+17+13+7 = 50$ et $15+22+13 = 50$. \checkmark

\textbf{2. Moyennes marginales.} Les données étant regroupées en classes, on utilise les \emph{centres de classes} : 15, 25, 35, 45 pour $X$ et 3,5 ; 6,5 ; 9,5 pour $Y$.
\begin{align*}
\bar{x} &= \frac{13\times 15 + 17\times 25 + 13\times 35 + 7\times 45}{50} = \frac{195 + 425 + 455 + 315}{50} = \frac{1390}{50} = 27{,}8 ; \\
\bar{y} &= \frac{15\times 3{,}5 + 22\times 6{,}5 + 13\times 9{,}5}{50} = \frac{52{,}5 + 143 + 123{,}5}{50} = \frac{319}{50} = 6{,}38.
\end{align*}
\textbf{Interprétation.} Le revenu mensuel moyen des ménages enquêtés est d'environ \num{278000} FCFA, et leur dépense téléphonique moyenne d'environ \num{6380} FCFA par mois.
\end{exoresolu}

\courssub{Le tableau à double entrée annoté}

La figure ci-dessous résume l'architecture d'un tableau à double entrée : repère bien la place des effectifs conjoints $n_{ij}$ (au cœur du tableau), des effectifs marginaux $n_{i\bullet}$ (marge droite) et $n_{\bullet j}$ (marge inférieure), et de l'effectif total $N$ (coin inférieur droit).

\begin{popfigure}
\begin{center}
\begin{tikzpicture}[x=1.55cm,y=0.62cm]
% Grille 4x4 + marges
\foreach \i in {0,...,5} \draw[popline] (\i,0) -- (\i,5);
\foreach \j in {0,...,5} \draw[popline] (0,\j) -- (5,\j);
% En-têtes
\node at (0.5,4.5) {\small $X\backslash Y$};
\node at (1.5,4.5) {\small $y_1$};
\node at (2.5,4.5) {\small $y_2$};
\node at (3.5,4.5) {\small $y_3$};
\node[poporange] at (4.5,4.5) {\small\textbf{Total}};
\node at (0.5,3.5) {\small $x_1$};
\node at (0.5,2.5) {\small $x_2$};
\node at (0.5,1.5) {\small $x_3$};
\node[poporange] at (0.5,0.5) {\small\textbf{Total}};
% Effectifs conjoints
\node[popblue] at (1.5,3.5) {$n_{11}$};
\node[popblue] at (2.5,3.5) {$n_{12}$};
\node[popblue] at (3.5,3.5) {$n_{13}$};
\node[popblue] at (1.5,2.5) {$n_{21}$};
\node[popblue] at (2.5,2.5) {$n_{22}$};
\node[popblue] at (3.5,2.5) {$n_{23}$};
\node[popblue] at (1.5,1.5) {$n_{31}$};
\node[popblue] at (2.5,1.5) {$n_{32}$};
\node[popblue] at (3.5,1.5) {$n_{33}$};
% Effectifs marginaux lignes
\node[poporange] at (4.5,3.5) {$n_{1\bullet}$};
\node[poporange] at (4.5,2.5) {$n_{2\bullet}$};
\node[poporange] at (4.5,1.5) {$n_{3\bullet}$};
% Effectifs marginaux colonnes
\node[popgreen] at (1.5,0.5) {$n_{\bullet 1}$};
\node[popgreen] at (2.5,0.5) {$n_{\bullet 2}$};
\node[popgreen] at (3.5,0.5) {$n_{\bullet 3}$};
% Effectif total
\node[poppurple] at (4.5,0.5) {$N$};
% Annotations
\draw[-{Stealth},popblue,thick] (2.5,2.9) -- (2.5,2.15);
\node[popblue,align=center] at (2.5,3.15) {\footnotesize effectifs conjoints};
\draw[-{Stealth},poporange,thick] (5.15,2.5) -- (4.85,2.5);
\node[poporange,align=left,anchor=west] at (5.2,2.5) {\footnotesize effectifs marginaux de $X$};
\draw[-{Stealth},popgreen,thick] (2.5,-0.15) -- (2.5,0.15);
\node[popgreen,align=center] at (2.5,-0.55) {\footnotesize effectifs marginaux de $Y$};
\end{tikzpicture}
\end{center}
\legende{Structure d'un tableau à double entrée : en bleu les effectifs conjoints $n_{ij}$, en orange les totaux de lignes $n_{i\bullet}$, en vert les totaux de colonnes $n_{\bullet j}$, en violet l'effectif total $N$.}
\end{popfigure}

\courssub{Données regroupées en classes et retour à la liste de couples}

\textbf{Centres de classes.} Lorsque les valeurs de $X$ ou de $Y$ sont regroupées en classes (comme dans l'enquête de Douala), on ne connaît plus les valeurs exactes : on remplace chaque classe par son \cle{centre}. Pour une classe $[a\,;\,b[$, le centre est $c = \dfrac{a+b}{2}$. Ainsi la classe $[20\,;\,30[$ est représentée par le centre $25$. Tous les calculs ultérieurs (moyennes, covariance, coefficient de corrélation) se feront sur les centres : c'est une approximation, mais c'est la convention officielle du programme.

\textbf{Du tableau à la liste de couples.} Pour les calculs des sections suivantes (point moyen, covariance, droites d'ajustement), il est souvent plus pratique de revenir à une liste de couples. Le principe : chaque case $(i\,;\,j)$ du tableau fournit le couple $(x_i\,;\,y_j)$ \emph{compté $n_{ij}$ fois}. Par exemple, dans l'enquête de Douala, la case $n_{22} = 9$ fournit neuf fois le couple $(25\,;\,6{,}5)$. On obtient ainsi une série de 50 couples (avec répétitions), sur laquelle toutes les formules des séries à une variable pondérées s'appliquent.

\begin{aretenir}[L'essentiel de la section]
\begin{itemize}
\item Une \cle{série statistique à deux caractères} associe à chaque individu d'une population un couple $(x_i\,;\,y_i)$.
\item Le \cle{tableau à double entrée} croise les valeurs de $X$ et de $Y$ ; $n_{ij}$ est l'effectif conjoint.
\item Les \cle{effectifs marginaux} s'obtiennent en sommant lignes ($n_{i\bullet}$) et colonnes ($n_{\bullet j}$) ; leur somme commune est $N$.
\item Les \cle{moyennes marginales} $\bar{x}$ et $\bar{y}$ sont les moyennes des séries marginales (avec les centres de classes si les données sont groupées).
\item Réflexe de contrôle : $\sum n_{ij} = \sum n_{i\bullet} = \sum n_{\bullet j} = N$.
\end{itemize}
\end{aretenir}

\begin{exercice}[Consommation selon l'âge du conducteur]
Un assureur de Bafoussam étudie, sur 40 conducteurs, l'âge $X$ (en années) et la consommation mensuelle de carburant $Y$ (en litres). Le tableau à double entrée donne : $n_{11} = 6$, $n_{12} = 2$, $n_{13} = 1$ (ligne $[20\,;\,30[$) ; $n_{21} = 5$, $n_{22} = 8$, $n_{23} = 3$ (ligne $[30\,;\,40[$) ; $n_{31} = 2$, $n_{32} = 7$, $n_{33} = 6$ (ligne $[40\,;\,50[$), les colonnes étant $[40\,;\,70[$, $[70\,;\,100[$ et $[100\,;\,130[$.
\begin{enumerate}
\item Construire le tableau complet avec les effectifs marginaux et vérifier que $N = 40$.
\item Déterminer les séries marginales de $X$ et de $Y$.
\item Calculer $\bar{x}$ et $\bar{y}$ à l'aide des centres de classes.
\end{enumerate}
\end{exercice}

\begin{corrige}[Exercice 1]
\textbf{1.} Totaux de lignes : $n_{1\bullet} = 6+2+1 = 9$ ; $n_{2\bullet} = 5+8+3 = 16$ ; $n_{3\bullet} = 2+7+6 = 15$. Totaux de colonnes : $n_{\bullet 1} = 6+5+2 = 13$ ; $n_{\bullet 2} = 2+8+7 = 17$ ; $n_{\bullet 3} = 1+3+6 = 10$. Contrôle : $9+16+15 = 40$ et $13+17+10 = 40$. \checkmark

\textbf{2.} Série marginale de $X$ : $([20;30[\,;\,9)$, $([30;40[\,;\,16)$, $([40;50[\,;\,15)$. Série marginale de $Y$ : $([40;70[\,;\,13)$, $([70;100[\,;\,17)$, $([100;130[\,;\,10)$.

\textbf{3.} Centres de $X$ : 25, 35, 45 ; centres de $Y$ : 55, 85, 115.
\[ \bar{x} = \frac{9\times 25 + 16\times 35 + 15\times 45}{40} = \frac{225+560+675}{40} = \frac{1460}{40} = 36{,}5 \text{ ans} ; \]
\[ \bar{y} = \frac{13\times 55 + 17\times 85 + 10\times 115}{40} = \frac{715+1445+1150}{40} = \frac{3310}{40} = 82{,}75 \text{ litres}. \]
\end{corrige}

\begin{savaistu}[Pourquoi « marginal » ?]
Le mot \emph{marginal} vient tout simplement de la \emph{marge} du tableau : les totaux $n_{i\bullet}$ et $n_{\bullet j}$ s'écrivent traditionnellement dans les bords droit et inférieur. La notation avec le point ($n_{i\bullet}$) signifie « on a sommé sur l'indice remplacé par le point » : $n_{i\bullet}$ est la somme sur toutes les valeurs de $j$. Ces séries marginales sont exactement les séries à une variable que tu as étudiées en Première : la statistique à deux variables les englobe et les relie. Dans la prochaine section, nous donnerons une image géométrique de ces données : le \emph{nuage de points}, première étape vers la prévision des ventes de Mme Etoundi.
\end{savaistu}

\coursec{Nuage de points et point moyen}

Dans la section précédente, nous avons appris à \emph{organiser} les données d'une série statistique à deux caractères : tableau à double entrée, effectifs marginaux, lectures croisées. Mais un tableau de nombres, aussi bien présenté soit-il, ne parle pas immédiatement à l'œil. Or, en statistiques comme en géométrie, \textbf{voir} c'est déjà comprendre à moitié. Lorsqu'un gérant de station-service à Bonabéri observe ses ventes d'essence semaine après semaine, il ne veut pas seulement des chiffres : il veut \emph{deviner une tendance}. C'est exactement le rôle du \cle{nuage de points} : transformer un tableau de couples $(x_i\,;\,y_i)$ en une image, puis résumer cette image par un point unique, le \cle{point moyen}, qui en sera le « centre de gravité ».

\courssub{Le nuage de points : donner une image aux données}

\textbf{Intuition.} Chaque couple $(x_i\,;\,y_i)$ de la série statistique est formé de deux nombres. Or deux nombres, c'est exactement ce qu'il faut pour placer un point dans un repère du plan : l'un en abscisse, l'autre en ordonnée. En plaçant ainsi tous les couples observés, on obtient un ensemble de points que l'on appelle, très naturellement, un \emph{nuage}.

\begin{definition}[Nuage de points]
Soit une série statistique à deux variables constituée de $n$ couples $(x_i\,;\,y_i)$, $i$ allant de $1$ à $n$. Dans un repère orthogonal $(O\,;\,\vec{\imath},\vec{\jmath})$ du plan, on appelle \cle{nuage de points} associé à cette série l'ensemble des points $M_i$ de coordonnées $(x_i\,;\,y_i)$ :
\[ \mathcal{N} = \left\{\, M_i\,(x_i\,;\,y_i)\ \middle|\ i \in \{1, 2, \ldots, n\} \,\right\}. \]
\end{definition}

\begin{attention}[Deux variables, deux rôles]
Par convention, la variable $x$ (dite \emph{variable explicative}) est portée en abscisse et la variable $y$ (dite \emph{variable expliquée}) en ordonnée. Au Baccalauréat, l'énoncé précise toujours lequel des deux caractères joue le rôle de $x$. Intervertir les axes change complètement l'allure du nuage et fausse toute la suite de l'exercice.
\end{attention}

\courssub{Construire un nuage lisible : le choix des échelles}

Un nuage de points n'est utile que s'il est \emph{lisible}. Un mauvais choix d'unités peut écraser le nuage sur une ligne ou le diluer dans une feuille entière. Voici la démarche professionnelle.

\begin{methode}[Construire un nuage de points lisible]
\begin{enumerate}
\item \textbf{Repérer les plages de valeurs} : noter les valeurs minimales et maximales de $x$ et de $y$ dans le tableau de données.
\item \textbf{Choisir les unités graphiques} de sorte que le nuage occupe une grande partie de la feuille (typiquement 10 à \SI{15}{cm} dans chaque direction), sans être coupé. Les unités peuvent être \emph{différentes} sur les deux axes : le repère est orthogonal, pas nécessairement orthonormé.
\item \textbf{Décaler l'origine si nécessaire} : si les valeurs de $x$ sont toutes comprises entre 40 et 50, il est inutile de graduer l'axe depuis 0. On fait commencer l'axe à 40 (on dit qu'on effectue un \emph{zoom sur la zone utile}), et on le signale proprement sur la copie.
\item \textbf{Graduer régulièrement} chaque axe et indiquer clairement les unités choisies (par exemple : « 1 cm représente 2 unités en abscisse »).
\item \textbf{Placer les points} $M_i(x_i\,;\,y_i)$ avec soin, en les marquant par des croix ou des petits points bien visibles. On ne relie \textbf{jamais} les points entre eux : un nuage n'est pas une courbe.
\end{enumerate}
\end{methode}

\begin{attention}[Pièges de construction]
\begin{itemize}
\item Ne jamais relier les points du nuage : il n'y a aucune raison que la « courbe » passant par les points ait un sens, les observations étant dispersées.
\item Ne pas oublier d'écrire les unités choisies sur la copie : au Baccalauréat, l'absence d'unités graphiques coûte des points de présentation.
\item Un décalage d'origine doit être \emph{visible} : on peut marquer une cassure sur l'axe ou écrire explicitement « l'axe des abscisses commence à 40 ».
\end{itemize}
\end{attention}

\courssub{Le point moyen : le centre de gravité du nuage}

\textbf{Intuition.} Imagine que chaque point du nuage soit une petite bille de masse identique posée sur une plaque. La plaque resterait en équilibre si on la posait sur une pointe placée exactement sous le \emph{centre de gravité} des billes. Ce centre de gravité, en statistiques, s'appelle le \cle{point moyen} : ses coordonnées sont tout simplement les moyennes des coordonnées des points du nuage.

\begin{definition}[Point moyen]
Soit un nuage de $n$ points $M_i(x_i\,;\,y_i)$. On appelle \cle{point moyen} du nuage le point $G$ de coordonnées $(\bar{x}\,;\,\bar{y})$, où :
\[ \bar{x} = \frac{1}{n}\sum_{i=1}^{n} x_i \qquad \text{et} \qquad \bar{y} = \frac{1}{n}\sum_{i=1}^{n} y_i. \]
Autrement dit : l'abscisse de $G$ est la moyenne des $x_i$, et son ordonnée est la moyenne des $y_i$.
\end{definition}

\begin{propriete}[Point moyen d'une série avec effectifs]
Si la série est donnée par un tableau d'effectifs (chaque couple $(x_i\,;\,y_i)$ étant observé $n_i$ fois, avec $n = \sum n_i$), alors :
\[ \bar{x} = \frac{1}{n}\sum_{i} n_i\, x_i \qquad \text{et} \qquad \bar{y} = \frac{1}{n}\sum_{i} n_i\, y_i. \]
Le point moyen est alors le barycentre des points $M_i$ affectés des coefficients $n_i$.
\end{propriete}

\begin{aretenir}[Ce qu'il faut retenir du point moyen]
\begin{itemize}
\item $G$ se calcule \textbf{séparément} sur chaque variable : moyenne des $x$, moyenne des $y$. C'est un calcul de moyenne de Première, appliqué deux fois.
\item $G$ résume la \emph{position} du nuage, mais pas sa \emph{forme}.
\item Propriété fondamentale pour la suite de la leçon : la droite des moindres carrés (section 4) \textbf{passe toujours par le point moyen} $G$. C'est pourquoi on le calcule presque systématiquement.
\end{itemize}
\end{aretenir}

\begin{attention}[Le point moyen n'est généralement pas un point du nuage]
C'est l'erreur classique du candidat pressé : chercher $G$ \emph{parmi} les points $M_i$. Le point moyen est un point \emph{calculé}, pas un point \emph{observé}. Il peut tomber au milieu du nuage, dans une zone vide où aucune observation ne se trouve. De même, la moyenne des notes d'une classe n'est la note d'aucun élève en particulier. Sur la figure, on marque $G$ d'un symbole différent (couleur, croix épaisse) pour bien le distinguer des points du nuage.
\end{attention}

\courssub{Exemple chiffré : les ventes de la station-service}

Reprenons la situation qui nous servira de fil rouge dans toute la leçon. Le gérant d'une station-service de Douala a relevé, pendant 8 semaines, le prix de vente du litre d'essence (en centaines de FCFA au-dessus d'un prix de référence) et la quantité vendue (en milliers de litres).

\begin{center}
\begin{tabularx}{\linewidth}{|l|X|X|X|X|X|X|X|X|}
\hline
\rowcolor{popblueL} Semaine $i$ & 1 & 2 & 3 & 4 & 5 & 6 & 7 & 8 \\
\hline
$x_i$ (prix, centaines F) & 1 & 2 & 2 & 3 & 4 & 5 & 5 & 6 \\
\hline
$y_i$ (ventes, milliers L) & 42 & 40 & 38 & 35 & 33 & 30 & 28 & 26 \\
\hline
\end{tabularx}
\end{center}

\begin{exemplebox}[Calcul du point moyen]
Calculons les coordonnées de $G$ :
\begin{align*}
\bar{x} &= \frac{1+2+2+3+4+5+5+6}{8} = \frac{28}{8} = 3{,}5 \\
\bar{y} &= \frac{42+40+38+35+33+30+28+26}{8} = \frac{272}{8} = 34
\end{align*}
Le point moyen est donc $G\,(3{,}5\,;\,34)$. On remarque qu'aucune semaine ne correspond au couple $(3{,}5\,;\,34)$ : $G$ n'est pas un point du nuage, conformément à la mise en garde précédente.
\end{exemplebox}

\begin{popfigure}
\begin{center}
\begin{tikzpicture}
\begin{axis}[
  width=0.85\linewidth, height=7cm,
  xlabel={Prix $x$ (centaines de FCFA)},
  ylabel={Ventes $y$ (milliers de litres)},
  xmin=0, xmax=7, ymin=20, ymax=46,
  xtick={0,1,2,3,4,5,6,7},
  ytick={20,25,30,35,40,45},
  grid=both, grid style={popline!40},
  axis lines=left,
]
\addplot[only marks, mark=*, mark size=2.5pt, popblue] coordinates {
  (1,42) (2,40) (2,38) (3,35) (4,33) (5,30) (5,28) (6,26)
};
\addplot[only marks, mark=*, mark size=4pt, popink] coordinates {(3.5,34)};
\node[popink, anchor=south west] at (axis cs:3.5,34) {$G\,(3{,}5\,;\,34)$};
\end{axis}
\end{tikzpicture}
\end{center}
\legende{Nuage de points des ventes de la station-service (8 semaines). Le point moyen $G$, marqué en rouge, n'est pas un point du nuage.}
\end{popfigure}

Sur cette figure, le nuage est visiblement \emph{allongé} autour d'une direction descendante : quand le prix augmente, les ventes diminuent, et ce de façon assez régulière. C'est exactement le type de lecture que nous allons maintenant apprendre à formaliser.

\courssub{Lecture qualitative d'un nuage : pressentir la corrélation}

Avant tout calcul, un bon statisticien \emph{regarde} son nuage. Trois questions suffisent :

\begin{enumerate}
\item \textbf{Le nuage est-il allongé ?} Si les points se regroupent autour d'une direction privilégiée, les deux variables semblent liées. Si le nuage est rond ou informe, les variables semblent indépendantes.
\item \textbf{Quelle est l'orientation ?} Si le nuage « monte » de gauche à droite, $y$ a tendance à augmenter quand $x$ augmente : on pressent une \cle{corrélation positive}. S'il « descend », on pressent une \cle{corrélation négative}.
\item \textbf{Quelle est la forme de l'allongement ?} Si le nuage s'allonge autour d'une \emph{droite}, un ajustement affine (sections 3 et 4) est pertinent. S'il s'allonge autour d'une \emph{courbe} (parabole, exponentielle...), il faudra autre chose.
\end{enumerate}

\begin{popfigure}
\begin{center}
\begin{tikzpicture}
\begin{axis}[
  name=plot1, width=0.34\linewidth, height=4.6cm,
  title={Corrélation positive forte},
  xmin=0, xmax=10, ymin=0, ymax=10,
  xtick=\empty, ytick=\empty, axis lines=left,
]
\addplot[only marks, mark=*, mark size=1.6pt, popblue] coordinates {
  (1,1.4) (2,2.3) (3,2.8) (4,4.1) (5,4.9) (6,6.2) (7,6.8) (8,8.1) (9,8.8)
};
\end{axis}
\begin{axis}[
  name=plot2, at={($(plot1.east)+(1.1cm,0)$)}, anchor=west,
  width=0.34\linewidth, height=4.6cm,
  title={Corrélation négative},
  xmin=0, xmax=10, ymin=0, ymax=10,
  xtick=\empty, ytick=\empty, axis lines=left,
]
\addplot[only marks, mark=*, mark size=1.6pt, poporange] coordinates {
  (1,8.6) (2,7.9) (3,7.2) (4,6.4) (5,5.1) (6,4.6) (7,3.4) (8,2.6) (9,1.5)
};
\end{axis}
\begin{axis}[
  name=plot3, at={($(plot2.east)+(1.1cm,0)$)}, anchor=west,
  width=0.34\linewidth, height=4.6cm,
  title={Absence de corrélation},
  xmin=0, xmax=10, ymin=0, ymax=10,
  xtick=\empty, ytick=\empty, axis lines=left,
]
\addplot[only marks, mark=*, mark size=1.6pt, popgreen] coordinates {
  (1,4) (2,7) (3,2.5) (4,6) (5,3.5) (6,8) (7,2) (8,6.5) (9,4.5) (3,7.5) (7,6)
};
\end{axis}
\end{tikzpicture}
\end{center}
\legende{Trois formes typiques de nuages : à gauche, un nuage allongé montant (corrélation positive forte) ; au centre, un nuage allongé descendant (corrélation négative) ; à droite, un nuage sans direction privilégiée (pas de corrélation linéaire).}
\end{popfigure}

\begin{aretenir}[Vocabulaire de la lecture d'un nuage]
\begin{itemize}
\item Nuage \textbf{allongé autour d'une droite montante} : corrélation linéaire \emph{positive} pressentie.
\item Nuage \textbf{allongé autour d'une droite descendante} : corrélation linéaire \emph{négative} pressentie.
\item Nuage \textbf{arrondi, sans direction} : pas de corrélation linéaire ; un ajustement affine n'a pas de sens.
\item Nuage \textbf{allongé autour d'une courbe} : il existe une liaison, mais \emph{non linéaire} ; l'ajustement affine direct est inadapté.
\end{itemize}
Cette lecture sera \emph{confirmée et mesurée} par le coefficient de corrélation (section 5) : l'œil propose, le calcul dispose.
\end{aretenir}

\courssub{Et si le nuage n'est pas rectiligne ? (complément utile au Bac)}

Il arrive — et le Baccalauréat aime cette situation — que le nuage soit nettement allongé mais \emph{courbé} : croissance de plus en plus rapide (allure exponentielle), ou au contraire qui s'aplanit (allure logarithmique). Dans ce cas, chercher une droite d'ajustement sur $(x\,;\,y)$ serait une erreur.

\begin{methode}[Changement de variable (signalé comme complément au programme)]
Lorsque le nuage des points $(x_i\,;\,y_i)$ n'est pas rectiligne, l'énoncé du Baccalauréat propose un \cle{changement de variable} pour se ramener à un nuage rectiligne :
\begin{enumerate}
\item On pose une nouvelle variable, par exemple $z_i = \ln(y_i)$, ou $z_i = \sqrt{y_i}$, ou $z_i = \dfrac{1}{x_i}$ (c'est \textbf{toujours l'énoncé qui le précise}).
\item On dresse le nouveau tableau des couples $(x_i\,;\,z_i)$ et on construit le nouveau nuage.
\item Si ce nouveau nuage est rectiligne, on effectue l'ajustement affine de $z$ en $x$ : $z = ax + b$.
\item On \textbf{revient à la variable initiale} : par exemple, si $\ln y = ax + b$, alors $y = \mathrm{e}^{ax+b} = \mathrm{e}^{b}\,\mathrm{e}^{ax}$ : l'ajustement de $y$ en $x$ est exponentiel.
\end{enumerate}
\end{methode}

\begin{attention}[Ne pas oublier le retour à la variable initiale]
L'erreur fatale : donner comme réponse finale $z = ax + b$ alors que la question porte sur $y$. Le changement de variable n'est qu'un \emph{détour de calcul} ; la réponse doit toujours être exprimée avec les variables de départ, celles du contexte concret (prix, ventes, effectifs...).
\end{attention}

\begin{exoresolu}[Nuage, point moyen et lecture qualitative]
Une agence MTN de Yaoundé étudie le lien entre le nombre $x$ de spots publicitaires diffusés dans la semaine et le nombre $y$ de nouvelles SIM activées (en centaines) :
\begin{center}
\begin{tabular}{|l|c|c|c|c|c|c|}
\hline
\rowcolor{popblueL} $x_i$ (spots) & 2 & 4 & 5 & 7 & 8 & 10 \\
\hline
$y_i$ (centaines de SIM) & 3 & 5 & 6 & 8 & 9 & 11 \\
\hline
\end{tabular}
\end{center}
\begin{enumerate}
\item Calculer les coordonnées du point moyen $G$.
\item Le point $G$ est-il un point du nuage ? Justifier.
\item Décrire qualitativement le nuage et dire si un ajustement affine semble pertinent.
\end{enumerate}
\tcblower
\textbf{Solution détaillée.}
\begin{enumerate}
\item On calcule séparément les deux moyennes :
\[ \bar{x} = \frac{2+4+5+7+8+10}{6} = \frac{36}{6} = 6 \qquad ; \qquad \bar{y} = \frac{3+5+6+8+9+11}{6} = \frac{42}{6} = 7. \]
Donc $G\,(6\,;\,7)$.
\item Le couple $(6\,;\,7)$ ne figure pas dans le tableau : aucune semaine n'a donné 6 spots pour 700 SIM. Le point moyen n'est donc \textbf{pas} un point du nuage : c'est un point calculé, centre de gravité du nuage.
\item En plaçant les points, on observe qu'ils sont presque alignés autour d'une direction \emph{montante} : quand le nombre de spots augmente, les activations augmentent de façon régulière. Le nuage est allongé et rectiligne, à corrélation positive forte : un \textbf{ajustement affine est tout à fait pertinent}. C'est ce que nous construirons dans les sections 3 et 4.
\end{enumerate}
\end{exoresolu}

\begin{savaistu}[D'où vient le mot « régression » ?]
Le terme \emph{régression}, que nous rencontrerons dans la section 4 avec les « droites de régression », a une origine surprenante. En 1885, le savant britannique Francis Galton (cousin de Darwin) étudia la transmission de la taille entre pères et fils. Il constata que les fils des pères très grands étaient grands, mais \emph{moins} que leurs pères, et que les fils des pères très petits étaient petits, mais \emph{moins} que leurs pères : les tailles semblaient « revenir » vers la moyenne. Galton parla de \emph{régression vers la médiocrité}. Le mot est resté, même si aujourd'hui il désigne simplement l'ensemble des méthodes d'ajustement. C'est en traçant des nuages de points de tailles (père, fils) et en cherchant leur point moyen que Galton a fondé la statistique à deux variables... exactement ce que tu viens d'apprendre à faire.
\end{savaistu}

\begin{exercice}[À toi de jouer]
Un coopérative agricole de l'Ouest a relevé, pour 6 exploitations, la surface $x$ plantée en maïs (en hectares) et la récolte $y$ (en tonnes) : $(2\,;\,5)$, $(3\,;\,7)$, $(4\,;\,9)$, $(5\,;\,10)$, $(7\,;\,15)$, $(9\,;\,18)$.
\begin{enumerate}
\item Construire le nuage de points (on choisira des unités adaptées et on les précisera).
\item Calculer les coordonnées du point moyen $G$ et le placer sur la figure.
\item Un ajustement affine semble-t-il justifié ? Justifier par une lecture qualitative du nuage.
\end{enumerate}
\end{exercice}

\begin{corrige}[Exercice]
\begin{enumerate}
\item Les valeurs de $x$ vont de 2 à 9, celles de $y$ de 5 à 18 : on peut prendre 1 cm pour 1 hectare en abscisse et 1 cm pour 2 tonnes en ordonnée, sans décalage d'origine. On place les six points, sans les relier.
\item $\bar{x} = \dfrac{2+3+4+5+7+9}{6} = \dfrac{30}{6} = 5$ et $\bar{y} = \dfrac{5+7+9+10+15+18}{6} = \dfrac{64}{6} \approx 10{,}67$. Donc $G\,(5\,;\,\frac{32}{3})$, soit environ $G\,(5\,;\,10{,}67)$. On le marque d'un symbole distinct.
\item Le nuage est nettement allongé autour d'une direction montante : la récolte croît régulièrement avec la surface. La corrélation linéaire positive est forte : un ajustement affine est justifié.
\end{enumerate}
\end{corrige}

Nous savons désormais \emph{voir} une série à deux variables et la résumer par son point moyen. Mais regarder ne suffit pas : il faut maintenant \emph{tracer la droite} qui résume le mieux le nuage. La section suivante présente une première méthode, simple et rapide, fondée précisément sur le point moyen : la \textbf{méthode de Mayer}.

\coursec{Ajustement affine par la méthode de Mayer}

Dans la section précédente, nous avons appris à représenter une série statistique à deux variables par un \cle{nuage de points} et à calculer son \cle{point moyen} $G$. Nous avons aussi observé que, lorsque le nuage s'allonge autour d'une direction privilégiée, il est naturel de chercher une droite qui le « résume ». Mais comment choisir cette droite de manière objective, sans la tracer « au jugé » ? La \cle{méthode de Mayer} est la première réponse rigoureuse à cette question : c'est une méthode simple, rapide, qui ne demande que des calculs de moyennes. Elle porte le nom de l'astronome allemand Tobias Mayer, qui l'utilisait au XVIII\textsuperscript{e} siècle pour ajuster ses observations de la Lune.

\courssub{Principe de la méthode : partager le nuage en deux}

\textbf{Une idée toute simple.}\par

Imagine un commerçant de Douala qui note chaque mois le prix du sac d'oignons et la quantité vendue. Les points obtenus s'alignent à peu près. Pour tracer une droite « au milieu » du nuage, l'idée de Mayer est la suivante : au lieu de chercher une droite qui passe « au milieu de tous les points » (problème difficile), on partage le nuage en \textbf{deux groupes} de même effectif, on calcule le point moyen de chaque groupe, et l'on trace la droite qui relie ces deux points moyens. Chaque point moyen « représente » fidèlement son groupe : la droite qui les joint passe donc naturellement au cœur du nuage.

\begin{definition}[Partage du nuage selon Mayer]
Soit un nuage de $n$ points $M_i(x_i\,;y_i)$ d'une série statistique à deux variables. On \textbf{ordonne les points par abscisses croissantes}, puis on partage le nuage en deux sous-nuages de même effectif (ou d'effectifs différant d'une unité) :
\begin{itemize}
\item le premier sous-nuage contient les points de plus petites abscisses ;
\item le second sous-nuage contient les points de plus grandes abscisses.
\end{itemize}
On note $G_1$ le point moyen du premier sous-nuage et $G_2$ celui du second.
\end{definition}

\begin{definition}[Droite de Mayer]
On appelle \cle{droite de Mayer} (ou droite d'ajustement affine par la méthode de Mayer) la droite $(G_1G_2)$ passant par les points moyens des deux sous-nuages.
\end{definition}

\begin{attention}[Classer avant de partager !]
Le partage se fait \textbf{impérativement après avoir rangé les points par abscisses croissantes}, et non dans l'ordre du tableau de données. Si le tableau n'est pas trié, les deux groupes ne correspondent plus aux « petites » et « grandes » valeurs de $x$, et la droite obtenue peut être complètement fausse. C'est l'erreur la plus fréquente au Baccalauréat : vérifie toujours l'ordre des $x_i$ avant de couper le nuage en deux.
\end{attention}

\courssub{Calcul des points moyens $G_1$ et $G_2$}

Rappelons la formule du point moyen vue dans la section précédente : le point moyen d'un nuage de points $(x_i\,;y_i)$ a pour coordonnées les moyennes arithmétiques des coordonnées. On l'applique ici séparément à chaque sous-nuage.

\begin{propriete}[Coordonnées de $G_1$ et $G_2$]
Si le premier sous-nuage contient les points $(x_1\,;y_1),\dots,(x_p\,;y_p)$ et le second les points $(x_{p+1}\,;y_{p+1}),\dots,(x_n\,;y_n)$, alors :
\[
G_1\left(\bar{x}_1 = \frac{x_1+\dots+x_p}{p}\ ;\ \bar{y}_1 = \frac{y_1+\dots+y_p}{p}\right)
\quad\text{et}\quad
G_2\left(\bar{x}_2 = \frac{x_{p+1}+\dots+x_n}{n-p}\ ;\ \bar{y}_2 = \frac{y_{p+1}+\dots+y_n}{n-p}\right).
\]
\end{propriete}

\begin{propriete}[La droite de Mayer équilibre les deux sous-nuages (admise)]
La droite de Mayer $(G_1G_2)$ passe par deux points qui sont les « centres de gravité » respectifs de la moitié gauche et de la moitié droite du nuage. Elle réalise ainsi un \textbf{équilibre} : les écarts des points du premier groupe se compensent autour de $G_1$, ceux du second groupe autour de $G_2$.
\end{propriete}

\begin{experience}[Justification intuitive du partage]
Pourquoi ce partage en deux donne-t-il une droite raisonnable ? Raisonnons pas à pas.
\begin{enumerate}
\item Le point moyen d'un nuage est son \emph{barycentre} : la somme des écarts $\sum (y_i - \bar{y})$ y est nulle. Autrement dit, $G_1$ est le point qui « résume » le mieux le premier groupe, et $G_2$ celui qui résume le mieux le second.
\item Une droite est entièrement déterminée par deux points. En imposant à la droite de passer par $G_1$ et $G_2$, on la force à passer par les deux « centres » du nuage.
\item Comme $G_1$ est à gauche (petites valeurs de $x$) et $G_2$ à droite (grandes valeurs de $x$), la droite $(G_1G_2)$ traverse le nuage dans sa longueur, en suivant sa direction naturelle.
\item De plus, on peut montrer que le point moyen global $G$ du nuage complet est le barycentre de $G_1$ et $G_2$ affectés des effectifs des deux groupes ; il appartient donc à la droite $(G_1G_2)$ : \textbf{la droite de Mayer passe toujours par le point moyen $G$} (lorsque les deux groupes ont le même effectif, $G$ est même le milieu de $[G_1G_2]$).
\end{enumerate}
Cette justification n'est pas exigible au Baccalauréat, mais elle éclaire pourquoi la méthode fonctionne.
\end{experience}

\begin{aretenir}[Point moyen et droite de Mayer]
La droite de Mayer passe par le \cle{point moyen} $G$ du nuage complet. C'est un excellent moyen de vérification : après avoir tracé la droite, contrôle qu'elle passe bien par $G(\bar{x}\,;\bar{y})$.
\end{aretenir}

\courssub{Équation réduite de la droite de Mayer}

Une fois $G_1(x_{G_1}\,;y_{G_1})$ et $G_2(x_{G_2}\,;y_{G_2})$ calculés, il reste à déterminer l'équation réduite $y = ax + b$ de la droite $(G_1G_2)$. C'est un calcul classique de Première.

\begin{propriete}[Équation de la droite $(G_1G_2)$]
Comme $x_{G_1} \neq x_{G_2}$ (les deux groupes portent sur des plages distinctes de valeurs de $x$), la droite $(G_1G_2)$ n'est pas verticale et admet une équation réduite $y = ax + b$ avec :
\[
a = \frac{y_{G_2}-y_{G_1}}{x_{G_2}-x_{G_1}}
\qquad\text{et}\qquad
b = y_{G_1} - a\,x_{G_1}.
\]
\end{propriete}

\begin{methode}[Déterminer une droite d'ajustement par la méthode de Mayer]
\begin{enumerate}
\item \textbf{Ordonner} les points du nuage par abscisses croissantes.
\item \textbf{Partager} le nuage en deux sous-nuages de même effectif (ou quasi si $n$ est impair).
\item \textbf{Calculer} les coordonnées des points moyens $G_1$ et $G_2$ de chaque sous-nuage (moyennes des abscisses, moyennes des ordonnées).
\item \textbf{Calculer} le coefficient directeur $a = \dfrac{y_{G_2}-y_{G_1}}{x_{G_2}-x_{G_1}}$.
\item \textbf{Calculer} l'ordonnée à l'origine $b = y_{G_1} - a\,x_{G_1}$ (on peut aussi utiliser $G_2$ : le résultat doit être identique, c'est une vérification).
\item \textbf{Écrire} l'équation $y = ax + b$ et \textbf{tracer} la droite sur le nuage, en vérifiant qu'elle passe par le point moyen $G$.
\end{enumerate}
\end{methode}

\courssub{Exemple complet : les ventes de carburant}

\begin{exemplebox}[Ventes de carburant à la station d'Ebolowa]
Une station-service d'Ebolowa relève, pendant 8 semaines, le prix de vente du litre d'essence $x_i$ (en centaines de FCFA au-dessus d'un prix de référence) et le volume hebdomadaire vendu $y_i$ (en milliers de litres). Les données, déjà rangées par abscisses croissantes, sont :
\begin{center}
\begin{tabularx}{\linewidth}{|l|*{8}{>{\centering\arraybackslash}X|}}
\hline
\rowcolor{popblueL} Semaine & 1 & 2 & 3 & 4 & 5 & 6 & 7 & 8 \\
\hline
$x_i$ & 1 & 2 & 3 & 4 & 5 & 6 & 7 & 8 \\
\hline
$y_i$ & 2 & 3 & 4 & 4,6 & 5 & 6,4 & 7 & 8 \\
\hline
\end{tabularx}
\end{center}
\textbf{Étape 1 et 2 : partage.} L'effectif est $n = 8$, pair : on partage en deux groupes de 4 points.
\[
\text{Groupe 1 : } (1\,;2),\ (2\,;3),\ (3\,;4),\ (4\,;4,6)
\qquad
\text{Groupe 2 : } (5\,;5),\ (6\,;6,4),\ (7\,;7),\ (8\,;8).
\]
\textbf{Étape 3 : points moyens.}
\begin{align*}
G_1 : \quad & \bar{x}_1 = \frac{1+2+3+4}{4} = \frac{10}{4} = 2{,}5
& \bar{y}_1 &= \frac{2+3+4+4{,}6}{4} = \frac{13{,}6}{4} = 3{,}4 \\
G_2 : \quad & \bar{x}_2 = \frac{5+6+7+8}{4} = \frac{26}{4} = 6{,}5
& \bar{y}_2 &= \frac{5+6{,}4+7+8}{4} = \frac{26{,}4}{4} = 6{,}6
\end{align*}
Donc $G_1(2{,}5\,;3{,}4)$ et $G_2(6{,}5\,;6{,}6)$.

\textbf{Étapes 4 et 5 : équation.}
\[
a = \frac{6{,}6-3{,}4}{6{,}5-2{,}5} = \frac{3{,}2}{4} = 0{,}8
\qquad
b = 3{,}4 - 0{,}8 \times 2{,}5 = 3{,}4 - 2 = 1{,}4.
\]
Vérification avec $G_2$ : $b = 6{,}6 - 0{,}8 \times 6{,}5 = 6{,}6 - 5{,}2 = 1{,}4$. Cohérent.

\textbf{Étape 6 : conclusion.} La droite de Mayer a pour équation :
\[
\boxed{y = 0{,}8x + 1{,}4}
\]
Vérification par le point moyen global : $\bar{x} = \dfrac{1+2+\dots+8}{8} = 4{,}5$ et $\bar{y} = \dfrac{40}{8} = 5$. Or $0{,}8 \times 4{,}5 + 1{,}4 = 3{,}6 + 1{,}4 = 5 = \bar{y}$ : le point $G(4{,}5\,;5)$ est bien sur la droite.
\end{exemplebox}

\begin{popfigure}
\begin{tikzpicture}
\begin{axis}[
width=\linewidth, height=7.5cm,
xmin=0, xmax=9, ymin=0, ymax=9.5,
xlabel={$x$ : prix (centaines de FCFA)}, ylabel={$y$ : volume (milliers de litres)},
grid=both, grid style={popline!40},
axis lines=left, ticklabel style={font=\small},
legend style={at={(0.02,0.98)}, anchor=north west, font=\small}
]
\addplot[popcurve, domain=0:9, samples=2] {0.8*x + 1.4};
\addlegendentry{Droite de Mayer $y=0{,}8x+1{,}4$}
\addplot[only marks, mark=*, mark size=2.2pt, popblue] coordinates {(1,2) (2,3) (3,4) (4,4.6)};
\addlegendentry{Sous-nuage 1}
\addplot[only marks, mark=*, mark size=2.2pt, poporange] coordinates {(5,5) (6,6.4) (7,7) (8,8)};
\addlegendentry{Sous-nuage 2}
\addplot[only marks, mark=square*, mark size=2.8pt, poppurple] coordinates {(2.5,3.4) (6.5,6.6)};
\addlegendentry{$G_1$ et $G_2$}
\addplot[only marks, mark=diamond*, mark size=3pt, popgreen] coordinates {(4.5,5)};
\addlegendentry{Point moyen $G$}
\node[font=\small, poppurple, anchor=south east] at (axis cs:2.5,3.4) {$G_1$};
\node[font=\small, poppurple, anchor=south east] at (axis cs:6.5,6.6) {$G_2$};
\node[font=\small, popgreen, anchor=north west] at (axis cs:4.5,5) {$G$};
\end{axis}
\end{tikzpicture}
\legende{Nuage des ventes de carburant partagé en deux sous-nuages (bleu et orange), points moyens $G_1$ et $G_2$, et droite de Mayer passant par le point moyen global $G$.}
\end{popfigure}

\textbf{Lecture graphique.}\par

Sur la figure, on observe que la droite de Mayer traverse le nuage dans sa longueur : les points bleus (petites valeurs de $x$) gravitent autour de $G_1$, les points orange autour de $G_2$, et la droite passe bien par le point moyen $G$. C'est la \textbf{vérification graphique} à effectuer systématiquement : si ta droite semble « à côté » du nuage, il y a une erreur de calcul ou de partage.

\courssub{Cas d'un nombre impair de points}

Lorsque $n$ est impair, on ne peut pas partager le nuage en deux groupes de même effectif. Deux conventions sont admises ; l'énoncé du Baccalauréat précise toujours laquelle utiliser.

\begin{aretenir}[Conventions pour $n$ impair]
Si le nuage comporte $n = 2p+1$ points, le point médian (celui de rang $p+1$) peut être :
\begin{itemize}
\item \textbf{exclu} du calcul : les deux groupes ont alors chacun $p$ points ;
\item \textbf{affecté aux deux groupes} : chaque groupe contient alors $p+1$ points, le point médian comptant deux fois.
\end{itemize}
Les deux conventions donnent des droites légèrement différentes : applique strictement celle indiquée dans l'énoncé.
\end{aretenir}

\begin{exoresolu}[Nuage de 7 points : dépenses de connexion Internet]
Une cybercafétéria de Bafoussam relève pendant 7 jours le nombre de clients $x_i$ et la recette $y_i$ (en milliers de FCFA). Les données sont déjà triées :
\begin{center}
\begin{tabularx}{\linewidth}{|l|*{7}{>{\centering\arraybackslash}X|}}
\hline
\rowcolor{popblueL} $x_i$ & 10 & 15 & 20 & 25 & 30 & 35 & 40 \\
\hline
$y_i$ & 5 & 7 & 9 & 10 & 13 & 14 & 17 \\
\hline
\end{tabularx}
\end{center}
Déterminer l'équation de la droite de Mayer en \textbf{excluant le point médian}.
\tcblower
\textbf{Solution.} L'effectif est $n = 7$ : le point médian est le 4\textsuperscript{e} point $(25\,;10)$, que l'on exclut. Les deux groupes sont :
\[
\text{Groupe 1 : } (10\,;5),\ (15\,;7),\ (20\,;9)
\qquad
\text{Groupe 2 : } (30\,;13),\ (35\,;14),\ (40\,;17).
\]
Points moyens :
\[
G_1\left(\frac{10+15+20}{3}\,;\,\frac{5+7+9}{3}\right) = G_1(15\,;7),
\qquad
G_2\left(\frac{30+35+40}{3}\,;\,\frac{13+14+17}{3}\right) = G_2\left(35\,;\frac{44}{3}\right).
\]
Coefficient directeur :
\[
a = \frac{\frac{44}{3}-7}{35-15} = \frac{\frac{44-21}{3}}{20} = \frac{23}{60} \approx 0{,}383.
\]
Ordonnée à l'origine : $b = 7 - \frac{23}{60}\times 15 = 7 - \frac{23}{4} = \frac{5}{4} = 1{,}25$.
La droite de Mayer a donc pour équation $y = \frac{23}{60}x + \frac{5}{4}$, soit approximativement $y \approx 0{,}38x + 1{,}25$.
\end{exoresolu}

\courssub{Limites de la méthode de Mayer}

La méthode de Mayer est séduisante par sa simplicité : quelques moyennes et une équation de droite suffisent. Mais il faut connaître ses faiblesses pour l'utiliser à bon escient.

\begin{attention}[Limites de la méthode de Mayer]
\begin{itemize}
\item \textbf{Dépendance au découpage.} Le résultat dépend de la façon dont on partage le nuage (surtout pour $n$ impair). Deux découpages différents donnent deux droites différentes : la méthode ne fournit pas une droite « unique et optimale ».
\item \textbf{Perte d'information.} Seules les moyennes des deux groupes interviennent : deux nuages très différents peuvent avoir les mêmes $G_1$ et $G_2$, donc la même droite de Mayer.
\item \textbf{Sensibilité aux points extrêmes.} Un point aberrant (erreur de saisie, événement exceptionnel) déplace fortement le point moyen de son groupe, donc la droite entière.
\item \textbf{Moins précise que les moindres carrés.} La méthode des moindres carrés, que nous étudierons dans la section suivante, utilise \emph{tous} les points et minimise globalement les écarts : elle fournit un ajustement mathématiquement optimal, au prix de calculs plus lourds.
\end{itemize}
En résumé : la méthode de Mayer est \textbf{rapide et suffisante pour une première estimation}, mais la droite des moindres carrés lui est généralement préférée dès qu'on dispose d'une calculatrice.
\end{attention}

\begin{savaistu}[Tobias Mayer, astronome cartographe]
Tobias Mayer (1723--1762) était astronome à l'observatoire de Göttingen. Pour cartographier la surface de la Lune avec une précision inédite, il devait combiner de nombreuses observations entachées d'erreurs : c'est ainsi qu'il eut l'idée de regrouper ses mesures et de moyenner les groupes, ancêtre de notre méthode. Ses cartes lunaires sont restées des références pendant des décennies. Aujourd'hui encore, au Cameroun, les techniciens de CAMTEL ou les météorologues qui ajustent des relevés dispersés utilisent, sans le savoir, l'héritage de cette idée : résumer des groupes de données par leurs moyennes avant d'ajuster un modèle.
\end{savaistu}

\begin{exercice}[À toi de jouer : livraisons d'une brasserie]
Une brasserie de Yaoundé note, sur 6 livraisons, la distance parcourue $x_i$ (en dizaines de km) et la consommation de gasoil $y_i$ (en litres) :
\begin{center}
\begin{tabularx}{\linewidth}{|l|*{6}{>{\centering\arraybackslash}X|}}
\hline
\rowcolor{popblueL} $x_i$ & 1 & 2 & 3 & 4 & 5 & 6 \\
\hline
$y_i$ & 9 & 12 & 14 & 17 & 20 & 22 \\
\hline
\end{tabularx}
\end{center}
\begin{enumerate}
\item Déterminer les points moyens $G_1$ et $G_2$ des deux sous-nuages.
\item En déduire l'équation de la droite de Mayer.
\item Vérifier qu'elle passe par le point moyen $G$ du nuage.
\end{enumerate}
\end{exercice}

\begin{corrige}[Exercice : livraisons d'une brasserie]
\begin{enumerate}
\item Groupe 1 : $(1\,;9),(2\,;12),(3\,;14)$, donc $G_1(2\,;\frac{35}{3})$ soit environ $G_1(2\,;11{,}67)$. Groupe 2 : $(4\,;17),(5\,;20),(6\,;22)$, donc $G_2(5\,;\frac{59}{3})$ soit environ $G_2(5\,;19{,}67)$.
\item $a = \dfrac{\frac{59}{3}-\frac{35}{3}}{5-2} = \dfrac{8}{3} \approx 2{,}67$ et $b = \frac{35}{3} - \frac{8}{3}\times 2 = \frac{19}{3} \approx 6{,}33$. D'où $y = \frac{8}{3}x + \frac{19}{3}$.
\item $\bar{x} = 3{,}5$ et $\bar{y} = \frac{94}{6} = \frac{47}{3}$. Or $\frac{8}{3}\times 3{,}5 + \frac{19}{3} = \frac{28}{3}+\frac{19}{3} = \frac{47}{3} = \bar{y}$ : $G$ est bien sur la droite.
\end{enumerate}
\end{corrige}

\coursec{Covariance et moindres carrés}

\courssub{Pourquoi dépasser la méthode de Mayer ?}

Dans la section précédente, tu as appris à ajuster un nuage de points par la \cle{méthode de Mayer} : on partage le nuage en deux sous-nuages, on calcule leurs points moyens $G_1$ et $G_2$, et l'on trace la droite $(G_1G_2)$. Cette méthode est rapide et donne souvent un résultat correct. Mais elle souffre de deux défauts majeurs :

\begin{itemize}
\item le résultat dépend du \emph{découpage} choisi : deux élèves qui regroupent différemment les points obtiennent deux droites différentes ;
\item aucun critère ne garantit que la droite obtenue est la « meilleure possible » : elle passe par deux points moyens, certes, mais rien ne dit qu'elle minimise l'erreur commise en remplaçant le nuage par une droite.
\end{itemize}

Les mathématiciens Legendre et Gauss, au tout début du XIX\textsuperscript{e} siècle, ont posé la bonne question : \emph{quelle droite rend la plus petite possible l'erreur globale d'ajustement ?} Leur réponse, la méthode des \cle{moindres carrés}, repose sur un outil nouveau qui mesure la façon dont deux caractères varient ensemble : la \cle{covariance}. C'est l'objet de cette section, au cœur du programme de Terminale et très fréquente au Baccalauréat.

\courssub{La covariance : mesurer la variation conjointe de deux caractères}

\textbf{Intuition.} Reprenons notre station-service de la route Yaoundé--Douala. On note $X$ le nombre de véhicules passant devant la station chaque jour et $Y$ le volume de carburant vendu (en hectolitres). Les jours où le trafic $X$ est \emph{au-dessus} de sa moyenne, les ventes $Y$ sont-elles plutôt \emph{au-dessus} ou \emph{en dessous} de leur moyenne ? Pour le savoir, l'idée naturelle est de regarder, pour chaque jour $i$, le produit des écarts aux moyennes :
\[
(x_i-\bar{x})(y_i-\bar{y}).
\]
Si $x_i$ et $y_i$ sont tous deux au-dessus (ou tous deux en dessous) de leur moyenne, ce produit est \textbf{positif}. Si l'un est au-dessus et l'autre en dessous, il est \textbf{négatif}. En faisant la moyenne de tous ces produits, on obtient un nombre qui résume le sens de la variation conjointe des deux caractères : c'est la covariance.

\begin{definition}[Covariance de deux variables statistiques]
Soit une série statistique double $(x_i\,;\,y_i)_{1\le i\le n}$ d'effectif total $N$, de moyennes $\bar{x}$ et $\bar{y}$. On appelle \cle{covariance} du couple $(X\,;\,Y)$ le nombre réel noté $\mathrm{cov}(X,Y)$ défini par :
\[
\mathrm{cov}(X,Y)=\frac{1}{N}\sum_{i=1}^{n}n_i\,(x_i-\bar{x})(y_i-\bar{y}).
\]
En développant le produit et en utilisant $\frac{1}{N}\sum n_i x_i=\bar{x}$ et $\frac{1}{N}\sum n_i y_i=\bar{y}$, on obtient la \textbf{forme développée}, beaucoup plus pratique pour les calculs :
\[
\boxed{\ \mathrm{cov}(X,Y)=\frac{1}{N}\sum_{i=1}^{n}n_i\,x_i y_i-\bar{x}\,\bar{y}\ }
\]
Autrement dit : \emph{moyenne des produits moins produit des moyennes}.
\end{definition}

\begin{aretenir}[Signe de la covariance]
\begin{itemize}
\item Si $\mathrm{cov}(X,Y)>0$ : les deux caractères ont tendance à varier \textbf{dans le même sens} (quand l'un augmente, l'autre augmente globalement). Le nuage « monte ».
\item Si $\mathrm{cov}(X,Y)<0$ : les deux caractères ont tendance à varier \textbf{en sens contraires}. Le nuage « descend ».
\item Si $\mathrm{cov}(X,Y)\approx 0$ : il n'y a pas de tendance linéaire globale entre les deux caractères.
\end{itemize}
\end{aretenir}

\textbf{Lien avec la variance.} Tu connais déjà la variance d'une série simple : $V(X)=\frac{1}{N}\sum n_i(x_i-\bar{x})^2$. Or $(x_i-\bar{x})^2=(x_i-\bar{x})(x_i-\bar{x})$, donc la variance n'est rien d'autre qu'une covariance particulière :
\[
V(X)=\mathrm{cov}(X,X)\qquad\text{et}\qquad V(Y)=\mathrm{cov}(Y,Y).
\]
La covariance généralise donc la variance : au lieu de mesurer comment un caractère varie par rapport à lui-même, elle mesure comment deux caractères varient \emph{l'un par rapport à l'autre}.

\begin{attention}[La covariance dépend des unités]
Contrairement au coefficient de corrélation que tu verras dans la section suivante, la covariance \textbf{dépend des unités} choisies : exprimer les ventes en litres plutôt qu'en hectolitres multiplie la covariance par $100\%$. On ne peut donc pas comparer directement deux covariances issues de séries différentes, ni conclure à une corrélation « forte » ou « faible » à partir de sa seule valeur. Son \textbf{signe}, en revanche, est toujours interprétable.
\end{attention}

\courssub{Le principe des moindres carrés}

\textbf{Intuition.} Cherchons une droite d'équation $y=ax+b$ qui « colle » au mieux au nuage. Pour chaque point $M_i(x_i\,;\,y_i)$, la droite prédit la valeur $ax_i+b$, alors que la valeur réellement observée est $y_i$. L'\textbf{écart vertical} (ou \cle{résidu}) est :
\[
e_i=y_i-(ax_i+b).
\]
On veut rendre tous ces écarts petits \emph{simultanément}. On ne peut pas simplement additionner les $e_i$, car les écarts positifs et négatifs se compenseraient. L'idée de Gauss : élever chaque écart au carré (ce qui les rend tous positifs et pénalise fortement les gros écarts), puis chercher $a$ et $b$ qui \textbf{minimisent} la somme :
\[
S(a,b)=\sum_{i=1}^{n}n_i\bigl(y_i-ax_i-b\bigr)^2.
\]
D'où le nom : méthode des \emph{moindres carrés}.

\begin{popfigure}
\begin{center}
\begin{tikzpicture}[scale=0.85]
\draw[->,popmuted] (-0.3,0) -- (8.5,0) node[below right]{$x$};
\draw[->,popmuted] (0,-0.3) -- (0,6) node[left]{$y$};
% droite de régression y = 0.6x + 1
\draw[popblue,very thick,domain=0:8] plot (\x,{0.6*\x+1}) node[right,popblue]{$D_{y/x}$};
% points et résidus
\foreach \x/\y in {1/2.1, 2/1.6, 3/3.4, 4/3.0, 5/4.6, 6/3.9, 7/5.6}{
  \pgfmathsetmacro{\yline}{0.6*\x+1}
  \draw[poporange,thick] (\x,\y) -- (\x,\yline);
  \fill[popink] (\x,\y) circle (2.2pt);
  \fill[popblue!60] (\x,\yline) circle (1.4pt);
}
\node[poporange] at (7.6,4.2) {$e_i$};
\fill[popgreen] (4,3.4) circle (3pt) node[below right,popgreen]{$G$};
\end{tikzpicture}
\end{center}
\legende{La droite des moindres carrés minimise la somme des carrés des segments verticaux (résidus $e_i$) figurés en orange. Elle passe par le point moyen $G$.}
\end{popfigure}

\begin{propriete}[Droite de régression de $y$ en $x$ — démonstration exigible]
Soit une série double de covariance $\mathrm{cov}(X,Y)$, avec $V(X)\neq 0$. Il existe une unique droite qui minimise $S(a,b)$ : c'est la \cle{droite de régression de $y$ en $x$}, notée $D_{y/x}$, d'équation :
\[
\boxed{\ y-\bar{y}=\frac{\mathrm{cov}(X,Y)}{V(X)}\,(x-\bar{x})\ }
\]
Son coefficient directeur $a=\dfrac{\mathrm{cov}(X,Y)}{V(X)}$ a donc le \textbf{signe de la covariance}.
\end{propriete}

\begin{experience}[Démonstration guidée du théorème des moindres carrés]
On cherche $a$ et $b$ minimisant $S(a,b)=\sum n_i(y_i-ax_i-b)^2$. On procède en deux étapes, en étudiant le trinôme du second degré obtenu (ou en annulant la dérivée).

\textbf{Étape 1 : choix optimal de $b$, à $a$ fixé.} Écrivons $y_i-ax_i-b=\bigl[(y_i-ax_i)-b\bigr]$ et posons $z_i=y_i-ax_i$. Alors $S=\sum n_i(z_i-b)^2$. Développons :
\[
S=\sum n_i z_i^2-2b\sum n_i z_i+Nb^2.
\]
À $a$ fixé, c'est un trinôme du second degré en $b$, de coefficient dominant $N>0$ : il est minimal lorsque sa dérivée s'annule, c'est-à-dire $-2\sum n_i z_i+2Nb=0$, soit :
\[
b=\frac{1}{N}\sum n_i z_i=\frac{1}{N}\sum n_i(y_i-ax_i)=\bar{y}-a\bar{x}.
\]
Ainsi, quel que soit $a$, le meilleur $b$ vérifie $\bar{y}=a\bar{x}+b$ : \emph{la droite optimale passe forcément par le point moyen $G(\bar{x}\,;\,\bar{y})$}.

\textbf{Étape 2 : choix optimal de $a$.} Reportons $b=\bar{y}-a\bar{x}$ dans $S$ :
\[
S(a)=\sum n_i\bigl[(y_i-\bar{y})-a(x_i-\bar{x})\bigr]^2.
\]
En développant (les sommes portent sur $i$, pondérées par $n_i$, et l'on divise par $N$ pour faire apparaître les quantités connues) :
\[
\frac{S(a)}{N}=V(Y)-2a\,\mathrm{cov}(X,Y)+a^2 V(X).
\]
C'est un trinôme du second degré en $a$, de coefficient dominant $V(X)>0$ : il atteint son minimum en
\[
a=-\frac{-2\,\mathrm{cov}(X,Y)}{2V(X)}=\frac{\mathrm{cov}(X,Y)}{V(X)}.
\]
\textbf{Conclusion.} La droite minimisant la somme des carrés des écarts a pour coefficient directeur $a=\dfrac{\mathrm{cov}(X,Y)}{V(X)}$ et passe par $G$ : son équation est bien $y-\bar{y}=\dfrac{\mathrm{cov}(X,Y)}{V(X)}(x-\bar{x})$. $\blacksquare$
\end{experience}

\begin{propriete}[Propriété fondamentale : passage par le point moyen]
La droite de régression de $y$ en $x$ passe par le point moyen $G(\bar{x}\,;\,\bar{y})$. C'est une conséquence directe de la forme de son équation : pour $x=\bar{x}$, on obtient $y=\bar{y}$. On verra ci-dessous que la droite de régression de $x$ en $y$ passe elle aussi par $G$.
\end{propriete}

\courssub{La droite de régression de $x$ en $y$}

Jusqu'ici, nous avons cherché à \emph{prédire $y$ à partir de $x$} en minimisant les écarts \textbf{verticaux}. Mais on peut vouloir faire l'inverse : prédire $x$ à partir de $y$ (par exemple, estimer le trafic à partir des ventes observées). On minimise alors la somme des carrés des écarts \textbf{horizontaux} $\sum n_i\bigl(x_i-a'y_i-b'\bigr)^2$. Le calcul est rigoureusement le même, en échangeant les rôles de $x$ et $y$ :

\begin{propriete}[Droite de régression de $x$ en $y$]
La \cle{droite de régression de $x$ en $y$}, notée $D_{x/y}$, a pour équation :
\[
\boxed{\ x-\bar{x}=\frac{\mathrm{cov}(X,Y)}{V(Y)}\,(y-\bar{y})\ }
\]
Elle passe également par le point moyen $G$.
\end{propriete}

\textbf{Comparaison des deux droites.} Les deux droites $D_{y/x}$ et $D_{x/y}$ passent toutes deux par $G$ : elles sont donc \textbf{sécantes en $G$}, sauf cas exceptionnel. Sont-elles confondues ? Rapportées au même repère, leurs coefficients directeurs respectifs sont :
\[
a=\frac{\mathrm{cov}(X,Y)}{V(X)}\qquad\text{et}\qquad \frac{1}{a'}=\frac{V(Y)}{\mathrm{cov}(X,Y)}.
\]
Elles sont confondues si et seulement si $a=\dfrac{1}{a'}$, c'est-à-dire $\bigl[\mathrm{cov}(X,Y)\bigr]^2=V(X)\,V(Y)$. On démontrera dans la section suivante que cela n'arrive que lorsque les points sont \textbf{parfaitement alignés} (corrélation parfaite). Plus le nuage est resserré autour d'une droite, plus les deux droites de régression sont proches l'une de l'autre ; plus le nuage est dispersé, plus l'angle entre elles est grand.

\begin{attention}[Les deux droites ne sont pas réciproques l'une de l'autre]
Erreur classique au Baccalauréat : croire que $D_{x/y}$ s'obtient en « isolant $x$ » dans l'équation de $D_{y/x}$. C'est \textbf{faux} ! Si l'on isole $x$ dans $y-\bar{y}=a(x-\bar{x})$, on obtient $x-\bar{x}=\dfrac{1}{a}(y-\bar{y})$ avec $\dfrac{1}{a}=\dfrac{V(X)}{\mathrm{cov}(X,Y)}$, alors que le vrai coefficient de $D_{x/y}$ est $\dfrac{\mathrm{cov}(X,Y)}{V(Y)}$. Ces deux coefficients ne sont égaux que si la corrélation est parfaite. Chaque droite répond à une question différente : $D_{y/x}$ sert à prédire $y$ connaissant $x$, $D_{x/y}$ sert à prédire $x$ connaissant $y$.
\end{attention}

\begin{popfigure}
\begin{center}
\begin{tikzpicture}[scale=0.85]
\draw[->,popmuted] (-0.3,0) -- (8.5,0) node[below right]{$x$};
\draw[->,popmuted] (0,-0.3) -- (0,6) node[left]{$y$};
\foreach \x/\y in {1/2.1, 2/1.6, 3/3.4, 4/3.0, 5/4.6, 6/3.9, 7/5.6}{
  \fill[popink] (\x,\y) circle (2.2pt);
}
% D_y/x : y = 0.6x + 1
\draw[popblue,very thick,domain=0:8] plot (\x,{0.6*\x+1});
\node[popblue,rotate=31] at (1.6,2.4) {$D_{y/x}$};
% D_x/y : x = 0.8 y + 1  =>  y = (x-1)/0.8
\draw[poporange,very thick,domain=1:5.8] plot (\x,{(\x-1)/0.8});
\node[poporange] at (6.9,5.8) {$D_{x/y}$};
\fill[popgreen] (4,3.4) circle (3.2pt);
\node[popgreen,below left] at (4,3.4) {$G(\bar{x}\,;\,\bar{y})$};
\end{tikzpicture}
\end{center}
\legende{Les deux droites de régression sont sécantes au point moyen $G$. Elles ne sont confondues que si les points sont parfaitement alignés.}
\end{popfigure}

\courssub{Organisation des calculs et exemple complet}

Les formules font intervenir $\bar{x}$, $\bar{y}$, $\sum n_i x_i^2$, $\sum n_i y_i^2$ et $\sum n_i x_i y_i$. Pour éviter toute erreur, on rassemble tout dans un tableau à cinq colonnes.

\begin{methode}[Organiser les calculs en tableau]
\begin{enumerate}
\item Construire un tableau à colonnes $x_i$, $y_i$, $x_i^2$, $y_i^2$, $x_i y_i$ (avec pondération par les effectifs $n_i$ si la série est groupée), puis une ligne « Totaux » contenant $\sum n_i x_i$, $\sum n_i y_i$, $\sum n_i x_i^2$, $\sum n_i y_i^2$, $\sum n_i x_i y_i$.
\item Calculer les moyennes : $\bar{x}=\dfrac{\sum n_i x_i}{N}$ et $\bar{y}=\dfrac{\sum n_i y_i}{N}$.
\item Calculer les variances par la formule de König-Huygens : $V(X)=\dfrac{\sum n_i x_i^2}{N}-\bar{x}^2$ et $V(Y)=\dfrac{\sum n_i y_i^2}{N}-\bar{y}^2$.
\item Calculer la covariance : $\mathrm{cov}(X,Y)=\dfrac{\sum n_i x_i y_i}{N}-\bar{x}\,\bar{y}$.
\item En déduire les équations : $y-\bar{y}=\dfrac{\mathrm{cov}(X,Y)}{V(X)}(x-\bar{x})$ et $x-\bar{x}=\dfrac{\mathrm{cov}(X,Y)}{V(Y)}(y-\bar{y})$.
\item Vérification rapide : le point $G(\bar{x}\,;\,\bar{y})$ doit appartenir aux deux droites obtenues.
\end{enumerate}
\end{methode}

\begin{exoresolu}[Ventes de carburant à la station Total de Mbankomo]
Le gérant d'une station-service relève, pendant 6 jours, le nombre $x_i$ de véhicules (en dizaines) passant devant la station et le volume $y_i$ de carburant vendu (en hectolitres) :
\begin{center}
\begin{tabular}{|l|*{6}{c|}}
\hline
\rowcolor{popblueL} Jour & 1 & 2 & 3 & 4 & 5 & 6 \\
\hline
$x_i$ (dizaines de véhicules) & 8 & 10 & 12 & 14 & 16 & 18 \\
\hline
$y_i$ (hL) & 12 & 15 & 16 & 20 & 21 & 26 \\
\hline
\end{tabular}
\end{center}
\begin{enumerate}
\item Calculer $\mathrm{cov}(X,Y)$ et interpréter son signe.
\item Déterminer l'équation de la droite de régression de $y$ en $x$.
\item Déterminer l'équation de la droite de régression de $x$ en $y$.
\end{enumerate}
\tcblower
\textbf{Solution détaillée.} On complète le tableau de calcul :
\begin{center}
\begin{tabular}{|c|c|c|c|c|}
\hline
\rowcolor{popblueL} $x_i$ & $y_i$ & $x_i^2$ & $y_i^2$ & $x_i y_i$ \\
\hline
8 & 12 & 64 & 144 & 96 \\
10 & 15 & 100 & 225 & 150 \\
12 & 16 & 144 & 256 & 192 \\
14 & 20 & 196 & 400 & 280 \\
16 & 21 & 256 & 441 & 336 \\
18 & 26 & 324 & 676 & 468 \\
\hline
\rowcolor{popcream} \textbf{Totaux ($\sum x_i=78$)} & $\sum y_i=110$ & $\sum x_i^2=1084$ & $\sum y_i^2=2142$ & $\sum x_i y_i=1522$ \\
\hline
\end{tabular}
\end{center}
avec $N=6$.

\textbf{1. Moyennes, variances et covariance.}
\[
\bar{x}=\frac{78}{6}=13 \qquad ; \qquad \bar{y}=\frac{110}{6}=\frac{55}{3}\approx 18{,}33.
\]
\[
V(X)=\frac{1084}{6}-13^2=\frac{1084}{6}-169=\frac{1084-1014}{6}=\frac{70}{6}=\frac{35}{3}\approx 11{,}67.
\]
\[
V(Y)=\frac{2142}{6}-\left(\frac{55}{3}\right)^2=357-\frac{3025}{9}=\frac{3213-3025}{9}=\frac{188}{9}\approx 20{,}89.
\]
\[
\mathrm{cov}(X,Y)=\frac{1522}{6}-13\times\frac{55}{3}=\frac{1522}{6}-\frac{715}{3}=\frac{1522-1430}{6}=\frac{92}{6}=\frac{46}{3}\approx 15{,}33.
\]
La covariance est \textbf{positive} : trafic et ventes varient dans le même sens, ce qui confirme l'intuition du gérant — plus il passe de véhicules, plus on vend de carburant.

\textbf{2. Droite de régression de $y$ en $x$.}
\[
a=\frac{\mathrm{cov}(X,Y)}{V(X)}=\frac{46/3}{35/3}=\frac{46}{35}\approx 1{,}314.
\]
L'équation est $y-\dfrac{55}{3}=\dfrac{46}{35}(x-13)$, soit en forme réduite :
\[
y=\frac{46}{35}x+\frac{55}{3}-\frac{46\times 13}{35}=\frac{46}{35}x+\frac{1925-1794}{105}=\frac{46}{35}x+\frac{131}{105}\approx 1{,}31x+1{,}25.
\]
Vérification : pour $x=\bar{x}=13$, $y=\frac{46}{35}\times 13+\frac{131}{105}=\frac{1794}{105}+\frac{131}{105}=\frac{1925}{105}=\frac{55}{3}=\bar{y}$. La droite passe bien par $G$.

\textbf{3. Droite de régression de $x$ en $y$.}
\[
a'=\frac{\mathrm{cov}(X,Y)}{V(Y)}=\frac{46/3}{188/9}=\frac{46}{3}\times\frac{9}{188}=\frac{414}{564}=\frac{69}{94}\approx 0{,}734.
\]
L'équation est $x-13=\dfrac{69}{94}\left(y-\dfrac{55}{3}\right)$, soit $x=\dfrac{69}{94}y+\dfrac{21}{94}\approx 0{,}73\,y+0{,}22$.
On remarque que $\dfrac{1}{a}=\dfrac{35}{46}\approx 0{,}76\neq a'\approx 0{,}73$ : les deux droites sont bien distinctes, comme prévu.
\end{exoresolu}

\begin{exercice}[À toi de jouer]
Une coopérative de planteurs de la région de l'Ouest relève, pour 5 exploitations, la surface $x_i$ plantée en cacao (en ha) et la récolte $y_i$ (en tonnes) : $(2\,;\,1{,}4)$, $(3\,;\,2{,}1)$, $(5\,;\,3{,}0)$, $(6\,;\,3{,}8)$, $(9\,;\,5{,}7)$.
\begin{enumerate}
\item Calculer $\bar{x}$, $\bar{y}$, $V(X)$, $V(Y)$ et $\mathrm{cov}(X,Y)$.
\item Déterminer l'équation réduite de la droite de régression de $y$ en $x$.
\item Vérifier que $G$ appartient à cette droite.
\end{enumerate}
\end{exercice}

\begin{corrige}[Exercice : coopérative de planteurs]
\textbf{1.} Tableau de calculs :
\begin{center}
\begin{tabular}{|c|c|c|c|c|}
\hline
\rowcolor{popblueL} $x_i$ & $y_i$ & $x_i^2$ & $y_i^2$ & $x_i y_i$ \\
\hline
2 & 1,4 & 4 & 1,96 & 2,8 \\
3 & 2,1 & 9 & 4,41 & 6,3 \\
5 & 3,0 & 25 & 9,00 & 15,0 \\
6 & 3,8 & 36 & 14,44 & 22,8 \\
9 & 5,7 & 81 & 32,49 & 51,3 \\
\hline
\rowcolor{popcream} \textbf{Totaux} & \textbf{16,0} & \textbf{155} & \textbf{62,30} & \textbf{98,2} \\
\hline
\end{tabular}
\end{center}
$\sum x_i = 25$, $N=5$.
\[
\bar{x}=\frac{25}{5}=5 \quad ; \quad \bar{y}=\frac{16,0}{5}=3{,}2
\]
\[
V(X)=\frac{155}{5}-5^2=31-25=6
\]
\[
V(Y)=\frac{62,30}{5}-3,2^2=12,46-10,24=2,22
\]
\[
\mathrm{cov}(X,Y)=\frac{98,2}{5}-5\times 3,2=19,64-16=3,64
\]

\textbf{2.} $a=\dfrac{\mathrm{cov}(X,Y)}{V(X)}=\dfrac{3,64}{6}=\dfrac{91}{150}\approx 0{,}6067$.
$b=\bar{y}-a\bar{x}=3,2-\dfrac{91}{150}\times 5=3,2-\dfrac{91}{30}=\dfrac{96-91}{30}=\dfrac{5}{30}=\dfrac{1}{6}\approx 0{,}1667$.
D'où $D_{y/x}:\ y=\dfrac{91}{150}x+\dfrac{1}{6}\approx 0{,}607x+0{,}167$.

\textbf{3.} Pour $x=5$ : $y=\dfrac{91}{150}\times 5+\dfrac{1}{6}=\dfrac{91}{30}+\dfrac{5}{30}=\dfrac{96}{30}=3{,}2=\bar{y}$. Le point $G(5\,;\,3{,}2)$ appartient bien à la droite.
\end{corrige}

\begin{savaistu}[Gauss, Legendre et l'astéroïde Cérès]
La méthode des moindres carrés est née d'une querelle de priorité entre deux géants : Adrien-Marie Legendre la publie en 1805, mais Carl Friedrich Gauss affirme l'utiliser depuis 1795. C'est Gauss qui lui doit sa consécration spectaculaire : en 1801, l'astéroïde Cérès, aperçu puis « perdu » par les astronomes, fut retrouvé exactement là où Gauss l'avait prédit grâce à un ajustement par moindres carrés des rares observations disponibles. Aujourd'hui, cette méthode est partout : prévisions de consommation d'électricité à ENEO, estimation du trafic par les opérateurs MTN et Orange, modèles agricoles de l'IRAD... Chaque fois qu'on prédit une grandeur à partir d'une autre, Gauss est dans la place.
\end{savaistu}

\begin{aretenir}[L'essentiel de la section]
\begin{itemize}
\item $\mathrm{cov}(X,Y)=\dfrac{1}{N}\sum n_i x_i y_i-\bar{x}\bar{y}$ ; son \textbf{signe} indique le sens de la variation conjointe.
\item $V(X)=\mathrm{cov}(X,X)$ et $V(Y)=\mathrm{cov}(Y,Y)$.
\item Droite de régression de $y$ en $x$ : $y-\bar{y}=\dfrac{\mathrm{cov}(X,Y)}{V(X)}(x-\bar{x})$.
\item Droite de régression de $x$ en $y$ : $x-\bar{x}=\dfrac{\mathrm{cov}(X,Y)}{V(Y)}(y-\bar{y})$.
\item Les deux droites passent par $G(\bar{x}\,;\,\bar{y})$ et ne sont confondues que si les points sont parfaitement alignés.
\item Les calculs s'organisent toujours dans un tableau à colonnes $x_i$, $y_i$, $x_i^2$, $y_i^2$, $x_i y_i$.
\end{itemize}
\end{aretenir}

Reste une question en suspens : nous savons désormais \emph{construire} la meilleure droite, mais comment \emph{juger} si cet ajustement est de bonne qualité, et s'il autorise des prévisions fiables ? C'est précisément le rôle du coefficient de corrélation linéaire, que nous étudions dans la section suivante.

\coursec{Coefficient de corrélation linéaire}

Dans la section précédente, nous avons appris à construire les droites de régression par la méthode des moindres carrés : la droite de régression de $y$ en $x$, d'équation $y = ax + b$ avec $a = \dfrac{\mathrm{cov}(X,Y)}{V(X)}$, et la droite de régression de $x$ en $y$, d'équation $x = a'y + b'$ avec $a' = \dfrac{\mathrm{cov}(X,Y)}{V(Y)}$. Mais une question fondamentale reste en suspens : \textbf{ces droites ont-elles un sens ?} Après tout, on peut calculer une droite des moindres carrés pour n'importe quel nuage de points, même un nuage complètement dispersé qui ne ressemble à rien ! Il nous faut donc un \cle{indicateur numérique} capable de mesurer la \emph{qualité} de l'alignement des points, avant de se lancer dans un ajustement et des prévisions. Cet indicateur, c'est le \cle{coefficient de corrélation linéaire}, noté $r$. C'est lui qui dira : « oui, un ajustement affine est justifié » ou « non, méfie-toi, ces points ne suivent pas une tendance linéaire ».

\courssub{Définition et premières propriétés}

\textbf{Intuition.}\par
La covariance $\mathrm{cov}(X,Y)$ mesure déjà le sens de la relation entre $X$ et $Y$ : positive si les deux variables varient dans le même sens, négative sinon. Mais elle possède un défaut majeur : elle dépend des \emph{unités}. Une covariance de $12$ entre une taille en centimètres et une masse en kilogrammes ne veut rien dire en soi ; la même covariance vaudrait $0{,}12$ si les tailles étaient exprimées en mètres. Pour comparer des situations différentes, il faut « normaliser » la covariance en la divisant par les écarts types des deux variables. On obtient alors un nombre sans unité, toujours compris entre $-1$ et $1$ : c'est le coefficient de corrélation linéaire.

\begin{definition}[Coefficient de corrélation linéaire]
Soit $(X, Y)$ une série statistique double dont les variances $V(X)$ et $V(Y)$ sont non nulles. Le \cle{coefficient de corrélation linéaire} de $X$ et $Y$ est le nombre réel noté $r$ défini par :
\[ r = \frac{\mathrm{cov}(X,Y)}{\sqrt{V(X)\,V(Y)}} = \frac{\mathrm{cov}(X,Y)}{\sigma_X \cdot \sigma_Y} \]
où $\sigma_X = \sqrt{V(X)}$ et $\sigma_Y = \sqrt{V(Y)}$ sont les écarts types de $X$ et de $Y$.
\end{definition}

\begin{propriete}[Encadrement de $r$ — admis]
Pour toute série statistique double $(X,Y)$ de variances non nulles :
\[ -1 \leq r \leq 1. \]
Ce résultat est une conséquence de l'\cle{inégalité de Cauchy-Schwarz}, qui affirme que $\bigl(\mathrm{cov}(X,Y)\bigr)^2 \leq V(X)\,V(Y)$. Cette inégalité est \emph{admise} en Terminale.

De plus, on a les cas limites suivants :
\begin{itemize}
\item $r = 1$ si et seulement si tous les points du nuage sont alignés sur une droite de coefficient directeur \textbf{strictement positif} ;
\item $r = -1$ si et seulement si tous les points du nuage sont alignés sur une droite de coefficient directeur \textbf{strictement négatif}.
\end{itemize}
\end{propriete}

\begin{experience}[Pourquoi $r = \pm 1$ signifie que les points sont alignés]
Cette démonstration guidée est formatrice : elle utilise la somme des carrés des écarts, déjà rencontrée avec les moindres carrés. On considère les variables centrées $u_i = x_i - \bar{x}$ et $v_i = y_i - \bar{y}$.

\textbf{Étape 1.} Pour tout réel $t$, considérons la somme des carrés :
\[ S(t) = \sum_{i=1}^{n} (v_i - t\,u_i)^2. \]
C'est une somme de carrés, donc $S(t) \geq 0$ pour tout $t$.

\textbf{Étape 2.} Développons :
\[ S(t) = \sum v_i^2 - 2t \sum u_i v_i + t^2 \sum u_i^2 = n\,V(Y) - 2t\,n\,\mathrm{cov}(X,Y) + t^2\,n\,V(X). \]
C'est un trinôme du second degré en $t$, toujours positif ou nul. Son discriminant est donc négatif ou nul :
\[ \Delta = 4n^2\bigl(\mathrm{cov}(X,Y)\bigr)^2 - 4n^2 V(X)V(Y) \leq 0, \]
ce qui redonne $\bigl(\mathrm{cov}(X,Y)\bigr)^2 \leq V(X)V(Y)$, c'est-à-dire $r^2 \leq 1$, donc $-1 \leq r \leq 1$.

\textbf{Étape 3.} Supposons maintenant $r^2 = 1$. Alors $\Delta = 0$ : le trinôme $S(t)$ admet une racine double $t_0 = \dfrac{\mathrm{cov}(X,Y)}{V(X)}$, et $S(t_0) = 0$. Or une somme de carrés nulle impose que \emph{chaque} carré soit nul :
\[ v_i = t_0\, u_i \quad \text{pour tout } i, \quad \text{c'est-à-dire} \quad y_i - \bar{y} = t_0\,(x_i - \bar{x}). \]
Tous les points $M_i(x_i\,;y_i)$ sont donc alignés sur la droite passant par le point moyen $G$, de coefficient directeur $t_0$. Comme $t_0$ a le signe de $\mathrm{cov}(X,Y)$, on obtient $t_0 > 0$ si $r = 1$ et $t_0 < 0$ si $r = -1$.

\textbf{Étape 4 (réciproque).} Si tous les points sont alignés sur une droite de pente non nulle, alors $y_i = \alpha x_i + \beta$ pour tout $i$, et un calcul direct donne $\mathrm{cov}(X,Y) = \alpha V(X)$ et $V(Y) = \alpha^2 V(X)$, d'où $r = \dfrac{\alpha V(X)}{|\alpha|\,V(X)} = \pm 1$ selon le signe de $\alpha$. \hfill$\blacksquare$
\end{experience}

\courssub{Interprétation du coefficient de corrélation}

Le signe et la valeur absolue de $r$ se lisent directement sur la forme du nuage de points :

\begin{itemize}
\item $r$ \textbf{proche de $1$} : les points sont presque alignés sur une droite « montante ». On dit qu'il y a une \cle{forte corrélation linéaire positive} : quand $X$ augmente, $Y$ augmente.
\item $r$ \textbf{proche de $-1$} : les points sont presque alignés sur une droite « descendante ». \cle{Forte corrélation linéaire négative} : quand $X$ augmente, $Y$ diminue.
\item $r$ \textbf{proche de $0$} : il n'y a \textbf{pas de corrélation linéaire} entre $X$ et $Y$. Attention : cela ne signifie pas qu'il n'y a \emph{aucune} relation (voir le piège plus bas).
\end{itemize}

\begin{popfigure}
\begin{tikzpicture}
\begin{axis}[width=0.48\linewidth,height=0.42\linewidth,title={$r \approx 0{,}98$},title style={font=\small},tick style={draw=none},xtick=\empty,ytick=\empty,axis lines=left,axis line style={popmuted},xmin=0,xmax=10,ymin=0,ymax=11]
\addplot[only marks,mark=*,mark size=1.4pt,popblue] coordinates {(1,1.6)(2,2.4)(3,3.1)(4,4.3)(5,5.2)(6,5.9)(7,7.3)(8,8.1)(9,9.2)};
\addplot[popcurve,domain=0.5:9.5]{0.95*x+0.6};
\end{axis}
\end{tikzpicture}\hfill
\begin{tikzpicture}
\begin{axis}[width=0.48\linewidth,height=0.42\linewidth,title={$r \approx -0{,}9$},title style={font=\small},tick style={draw=none},xtick=\empty,ytick=\empty,axis lines=left,axis line style={popmuted},xmin=0,xmax=10,ymin=0,ymax=11]
\addplot[only marks,mark=*,mark size=1.4pt,poporange] coordinates {(1,9.4)(2,8.2)(3,8.6)(4,6.4)(5,6.8)(6,4.7)(7,4.9)(8,2.6)(9,1.8)};
\addplot[popcurve,poporange,domain=0.5:9.5]{-0.92*x+9.8};
\end{axis}
\end{tikzpicture}

\vspace{0.4cm}
\begin{tikzpicture}
\begin{axis}[width=0.48\linewidth,height=0.42\linewidth,title={$r \approx 0{,}3$},title style={font=\small},tick style={draw=none},xtick=\empty,ytick=\empty,axis lines=left,axis line style={popmuted},xmin=0,xmax=10,ymin=0,ymax=11]
\addplot[only marks,mark=*,mark size=1.4pt,popgreen] coordinates {(1,2.5)(1.5,6.5)(2.5,3.2)(3,7.8)(4,4.1)(4.5,8.9)(5.5,5.3)(6,2.2)(6.5,9.4)(7.5,6.1)(8,3.4)(8.5,8.2)(9.2,5.6)};
\end{axis}
\end{tikzpicture}\hfill
\begin{tikzpicture}
\begin{axis}[width=0.48\linewidth,height=0.42\linewidth,title={$r \approx 0$ (relation non linéaire)},title style={font=\small},tick style={draw=none},xtick=\empty,ytick=\empty,axis lines=left,axis line style={popmuted},xmin=0,xmax=10,ymin=0,ymax=11]
\addplot[only marks,mark=*,mark size=1.4pt,poppurple] coordinates {(1,5.2)(2,2.6)(3,1.2)(4,0.6)(5,0.5)(6,0.8)(7,1.6)(8,3.1)(9,5.3)};
\addplot[poppurple,thick,domain=0.8:9.2,samples=100]{0.3*(x-5)^2+0.5};
\end{axis}
\end{tikzpicture}
\legende{Quatre nuages de points et leurs coefficients de corrélation. Le dernier nuage montre une relation parfaitement organisée... mais non linéaire : $r$ est proche de $0$ malgré un lien évident entre $X$ et $Y$.}
\end{popfigure}

\begin{attention}[Un $r$ proche de $0$ ne veut pas dire « pas de relation »]
Le coefficient $r$ ne mesure \textbf{que} la dépendance \emph{linéaire}. Si les points suivent une parabole, une courbe exponentielle ou toute autre relation non linéaire, on peut obtenir $r \approx 0$ alors que $X$ et $Y$ sont parfaitement liées (voir le quatrième nuage de la galerie ci-dessus). \textbf{Réflexe de survie au Bac :} toujours regarder le nuage de points avant de conclure ! Si le nuage suggère une courbe, un ajustement affine n'est pas adapté, quelle que soit la valeur de $r$.
\end{attention}

\courssub{Lien avec les droites de régression}

Rappelons les coefficients directeurs des deux droites de régression obtenues par les moindres carrés :
\[ a = \frac{\mathrm{cov}(X,Y)}{V(X)} \qquad \text{et} \qquad a' = \frac{\mathrm{cov}(X,Y)}{V(Y)}. \]

\begin{propriete}[$r^2 = a \times a'$]
En notant $a$ le coefficient directeur de la droite de régression de $y$ en $x$ et $a'$ celui de la droite de régression de $x$ en $y$, on a :
\[ a \times a' = \frac{\mathrm{cov}(X,Y)}{V(X)} \times \frac{\mathrm{cov}(X,Y)}{V(Y)} = \frac{\bigl(\mathrm{cov}(X,Y)\bigr)^2}{V(X)\,V(Y)} = r^2. \]
\end{propriete}

Cette formule est très riche en conséquences pratiques :

\begin{itemize}
\item Comme $r^2 = a\,a' \geq 0$, les deux pentes $a$ et $a'$ ont \textbf{toujours le même signe} : les deux droites de régression « montent » ensemble ou « descendent » ensemble.
\item $r$ a le même signe que $\mathrm{cov}(X,Y)$, donc le même signe que $a$ et $a'$. Connaissant $a$ et $a'$, on retrouve $r$ : $r = \pm\sqrt{a\,a'}$, avec le signe de $a$.
\item $r = \pm 1$ équivaut à $a \times a' = 1$ : les deux droites de régression sont alors \textbf{confondues} (les points sont alignés). Plus $|r|$ est proche de $1$, plus les deux droites sont proches l'une de l'autre.
\end{itemize}

\begin{exemplebox}[Retrouver $r$ à partir des pentes]
Pour une série double, on a trouvé les droites de régression $y = 0{,}8x + 2$ et $x = 0{,}45y + 1$. Alors :
\[ r^2 = a \times a' = 0{,}8 \times 0{,}45 = 0{,}36 \quad \text{donc} \quad r = \sqrt{0{,}36} = 0{,}6 \]
(le signe est positif car $a > 0$). La corrélation linéaire est moyenne.
\end{exemplebox}

\courssub{Qualité d'un ajustement affine : le critère de décision}

\begin{methode}[Interpréter $r$ et juger la qualité d'un ajustement affine]
Face à une série statistique double, voici la démarche complète à rédiger au Bac :
\begin{enumerate}
\item \textbf{Observer le nuage de points} : une forme allongée autour d'une droite est un bon signe ; une forme courbe ou éclatée est un mauvais signe.
\item \textbf{Calculer le coefficient de corrélation} $r = \dfrac{\mathrm{cov}(X,Y)}{\sigma_X \sigma_Y}$.
\item \textbf{Appliquer le critère usuel} :
\begin{itemize}
\item si $|r| \geq 0{,}75$ (c'est-à-dire $r^2 \geq 0{,}5625$, en pratique on exige souvent $r^2$ proche de $1$) : la corrélation linéaire est \textbf{forte}, un ajustement affine est \textbf{justifié} ;
\item si $|r| < 0{,}75$ : la corrélation linéaire est \textbf{faible}, un ajustement affine est \textbf{peu pertinent} (ou à proscrire si le nuage est courbe).
\end{itemize}
\item \textbf{Conclure par une phrase} : « Comme $r \approx \ldots$ avec $|r| \geq 0{,}75$, la corrélation linéaire entre $X$ et $Y$ est forte : l'ajustement affine est justifié, et les prévisions issues de cette droite seront fiables. »
\end{enumerate}
\end{methode}

\begin{exoresolu}[Réseau mobile et couverture d'arrondissement]
Une société de téléphonie mobile étudie, dans 6 arrondissements, le lien entre le nombre $x_i$ d'antennes relais installées et le pourcentage $y_i$ de la population couverte par le réseau 4G.

\begin{center}
\begin{tabularx}{\linewidth}{|l|*{6}{>{\centering\arraybackslash}X|}}
\hline
\rowcolor{popblueL} $x_i$ (antennes) & 2 & 4 & 5 & 8 & 10 & 12 \\
\hline
$y_i$ (couverture en \%) & 35 & 48 & 55 & 70 & 79 & 88 \\
\hline
\end{tabularx}
\end{center}

On donne : $\bar{x} = \dfrac{41}{6} \approx 6{,}83$ ; $\bar{y} = 62{,}5$ ; $V(X) = \dfrac{437}{36} \approx 12{,}14$ ; $V(Y) = \dfrac{12009}{36} \approx 333{,}58$ ; $\mathrm{cov}(X,Y) = \dfrac{2283}{36} \approx 63{,}42$.

\textbf{1.} Calculer le coefficient de corrélation linéaire $r$ (arrondir à $10^{-3}$).

\textbf{2.} Un ajustement affine est-il justifié ?

\textbf{3.} Sachant que la droite de régression de $y$ en $x$ a pour coefficient directeur $a = \dfrac{2283}{437} \approx 5{,}22$, retrouver le coefficient directeur $a'$ de la droite de régression de $x$ en $y$.

\tcblower
\textbf{1.} Appliquons la définition :
\[ r = \frac{\mathrm{cov}(X,Y)}{\sqrt{V(X)\,V(Y)}} = \frac{2283/36}{\sqrt{(437/36) \times (12009/36)}} = \frac{2283}{\sqrt{437 \times 12009}} \approx \frac{2283}{2290{,}84} \approx 0{,}997. \]

\textbf{2.} On a $|r| \approx 0{,}997 \geq 0{,}75$ : le coefficient est très proche de $1$. La corrélation linéaire entre le nombre d'antennes et la couverture 4G est \textbf{très forte et positive}. Un ajustement affine est donc parfaitement justifié, et les prévisions faites à partir de la droite de régression seront fiables.

\textbf{3.} On utilise la propriété $r^2 = a \times a'$ :
\[ a' = \frac{r^2}{a} = \frac{(0{,}997)^2}{5{,}22} \approx \frac{0{,}994}{5{,}22} \approx 0{,}190. \]
La droite de régression de $x$ en $y$ a pour coefficient directeur $a' \approx 0{,}190$. Vérification : $a$ et $a'$ sont bien de même signe (positifs), conformément à la théorie.
\end{exoresolu}

\begin{attention}[Pièges fréquents sur le coefficient de corrélation]
\begin{itemize}
\item \textbf{Ne pas confondre $r$ et $r^2$} : le critère porte sur $|r| \geq 0{,}75$. Un élève qui compare $r = 0{,}6$ à $0{,}75$ en oubliant la valeur absolue, ou qui conclut avec $r^2 = 0{,}6$ en le comparant à $0{,}75$, fait une erreur de jugement.
\item \textbf{Le signe de $r$} : en extrayant la racine carrée de $r^2 = a\,a'$, n'oublie pas de choisir le signe de $\mathrm{cov}(X,Y)$ (celui de $a$). Écrire $r = 0{,}6$ quand les pentes sont négatives est faux : c'est $r = -0{,}6$.
\item \textbf{Divisions} : $r$ se calcule avec le produit $\sigma_X \sigma_Y$ au dénominateur, pas avec la somme ni avec les variances non racinées.
\item \textbf{Conclusion obligatoire} : au Bac, calculer $r$ sans phrase d'interprétation fait perdre des points. Toujours conclure sur la \emph{justification de l'ajustement affine}.
\end{itemize}
\end{attention}

\begin{savaistu}[Corrélation n'est pas causalité !]
Une corrélation forte entre deux phénomènes ne prouve \textbf{jamais} que l'un est la cause de l'autre. Exemple classique : sur une année, les ventes de glaces et le nombre de noyades sont fortement corrélés ($r$ proche de $1$). Faut-il interdire les glaces pour sauver des vies ? Bien sûr que non : les deux phénomènes dépendent d'une troisième variable cachée, la \emph{chaleur estivale}, qui pousse à la fois à manger des glaces et à se baigner. On parle de \cle{variable confondante}.

D'autres corrélations fallacieuses célèbres : le nombre de cigognes et le taux de natalité dans certaines régions d'Europe, ou encore, chez nous, on pourrait observer une corrélation entre les ventes d'ombrelles et les dégâts de pluie sur les routes de Yaoundé — ce n'est évidemment pas l'ombrelle qui abîme la route ! En statistiques, $r$ mesure une \emph{association}, jamais un \emph{lien de cause à effet}. Seule une étude approfondie (expérience contrôlée, analyse du contexte) peut établir une causalité. Retiens la formule des statisticiens : \emph{« Correlation is not causation »}.
\end{savaistu}

\begin{exercice}[À toi de jouer]
Le tableau suivant donne, pour 7 élèves de Terminale C, le nombre d'heures $x_i$ de révision hebdomadaire et la moyenne $y_i$ obtenue en mathématiques.

\begin{center}
\begin{tabularx}{\linewidth}{|l|*{7}{>{\centering\arraybackslash}X|}}
\hline
\rowcolor{popblueL} $x_i$ (heures) & 2 & 4 & 5 & 7 & 9 & 11 & 12 \\
\hline
$y_i$ (moyenne) & 8 & 9 & 11 & 12 & 14 & 15 & 16 \\
\hline
\end{tabularx}
\end{center}

On donne : $V(X) = \dfrac{580}{49} \approx 11{,}84$ ; $V(Y) = \dfrac{384}{49} \approx 7{,}84$ ; $\mathrm{cov}(X,Y) = \dfrac{468}{49} \approx 9{,}55$.
\begin{enumerate}
\item Calculer le coefficient de corrélation linéaire $r$, arrondi à $10^{-3}$.
\item Un ajustement affine du nuage est-il justifié ? Justifier soigneusement.
\item La droite de régression de $y$ en $x$ a pour équation $y = 0{,}76x + 6{,}67$. Déterminer le coefficient directeur de la droite de régression de $x$ en $y$.
\end{enumerate}
\end{exercice}

\begin{corrige}[Exercice : révisions et moyennes]
\textbf{1.} $r = \dfrac{468/49}{\sqrt{(580/49) \times (384/49)}} = \dfrac{468}{\sqrt{580 \times 384}} = \dfrac{468}{\sqrt{222720}} \approx \dfrac{468}{471{,}93} \approx 0{,}992$.

\textbf{2.} $|r| \approx 0{,}992 \geq 0{,}75$ : la corrélation linéaire est très forte et positive. L'ajustement affine est justifié.

\textbf{3.} $a' = \dfrac{r^2}{a} = \dfrac{0{,}992^2}{0{,}76} \approx \dfrac{0{,}984}{0{,}76} \approx 1{,}29$.
\end{corrige}

\begin{aretenir}[L'essentiel sur le coefficient de corrélation linéaire]
\begin{itemize}
\item $r = \dfrac{\mathrm{cov}(X,Y)}{\sigma_X\,\sigma_Y}$, et l'on a toujours $-1 \leq r \leq 1$.
\item $r = \pm 1$ si et seulement si les points du nuage sont parfaitement alignés (pente positive pour $r = 1$, négative pour $r = -1$).
\item $r^2 = a \times a'$ (produit des coefficients directeurs des deux droites de régression) ; $r$ a le signe de $\mathrm{cov}(X,Y)$.
\item Critère de qualité : $|r| \geq 0{,}75$ $\Rightarrow$ corrélation linéaire forte, ajustement affine justifié.
\item $r \approx 0$ n'exclut pas une relation \emph{non linéaire} : toujours examiner le nuage.
\item Une corrélation, même forte, ne prouve jamais une relation de cause à effet.
\end{itemize}
\end{aretenir}

\coursec{Prévisions et synthèse méthodologique}

Dans la section précédente, nous avons appris à mesurer la \cle{qualité} d'un ajustement affine grâce au \cle{coefficient de corrélation linéaire} $r$. Nous savons donc désormais construire un nuage de points, placer le point moyen $G$, tracer une droite d'ajustement (Mayer ou moindres carrés) et juger si cet ajustement est pertinent. Il reste à répondre à la question qui motive toute cette leçon : \emph{à quoi sert une droite d'ajustement ?} Réponse : à \cle{prévoir}. Le gérant d'une station-service de Yaoundé qui connaît la relation entre le trafic routier et ses ventes de carburant veut anticiper ses commandes ; l'INS veut projeter la population du Cameroun en 2035 ; un médecin veut estimer la dose d'un médicament adaptée à la masse d'un patient. Dans tous ces cas, on \emph{substitue} une valeur dans l'équation de la droite. Cette dernière section te donne la méthode complète, ses limites, et un complément précieux pour le Bac : les ajustements non affines par changement de variable.

\courssub{Principe de la prévision}

\textbf{Intuition.} Une droite d'ajustement $y = ax + b$ résume la tendance globale du nuage par une relation \emph{fonctionnelle} : à chaque valeur de $x$ correspond une valeur estimée de $y$. Prévoir, c'est simplement \textbf{lire cette relation} pour une valeur de $x$ qui n'a pas (encore) été observée.

\begin{definition}[Prévision à l'aide d'un ajustement affine]
Soit une série statistique à deux variables $(x\,;y)$ ajustée par une droite. Effectuer une \cle{prévision} consiste à :
\begin{itemize}
\item estimer $y$ pour une valeur $x_0$ donnée en calculant $\widehat{y_0} = a x_0 + b$ à l'aide de la \cle{droite de régression de $y$ en $x$} ;
\item estimer $x$ pour une valeur $y_0$ donnée en calculant $\widehat{x_0} = a' y_0 + b'$ à l'aide de la \cle{droite de régression de $x$ en $y$}.
\end{itemize}
\end{definition}

\begin{attention}[La bonne droite pour la bonne question]
Les deux droites de régression ne sont pas interchangeables ! Elles minimisent des écarts différents (verticaux pour $y$ en $x$, horizontaux pour $x$ en $y$) et ne coïncident que si $|r| = 1$. Retiens :
\begin{itemize}
\item pour estimer $y$ \textbf{connaissant} $x$ : droite de régression de $y$ en $x$, $y = ax+b$ ;
\item pour estimer $x$ \textbf{connaissant} $y$ : droite de régression de $x$ en $y$, $x = a'y+b'$.
\end{itemize}
Erreur classique au Bac : isoler $x$ dans l'équation $y = ax+b$ pour estimer $x$ connaissant $y$. C'est mathématiquement possible, mais ce n'est \emph{pas} la droite de régression de $x$ en $y$, et l'examinateur le sanctionne.
\end{attention}

\courssub{Interpolation, extrapolation et condition de fiabilité}

\begin{definition}[Interpolation et extrapolation]
Soit $[x_{\min}\,;\,x_{\max}]$ la plage des valeurs observées de $x$.
\begin{itemize}
\item Une prévision est une \cle{interpolation} si $x_0 \in [x_{\min}\,;\,x_{\max}]$ : on estime à l'intérieur de la zone où le modèle a été « calibré ».
\item C'est une \cle{extrapolation} si $x_0 \notin [x_{\min}\,;\,x_{\max}]$ : on prolonge la tendance au-delà des observations.
\end{itemize}
\end{definition}

\begin{propriete}[Conditions de fiabilité d'une prévision]
Une prévision par ajustement affine n'est raisonnable que si :
\begin{enumerate}
\item la corrélation est forte, en pratique $|r| \geq \num{0,87}$ (au minimum $\num{0,75}$ pour une estimation grossière) ;
\item la prévision est une interpolation, ou une extrapolation \emph{proche} de la plage observée ;
\item le phénomène étudié n'a pas subi de changement de contexte (crise, nouvelle technologie, saturation\dots).
\end{enumerate}
\end{propriete}

\begin{attention}[Le piège de l'extrapolation lointaine]
Une extrapolation lointaine des données est \textbf{toujours risquée, même avec un excellent $r$}. Le coefficient $r$ mesure la linéarité du nuage \emph{observé}, il ne dit rien du comportement du phénomène en dehors. Exemple : la taille d'un enfant entre 4 et 10 ans croît presque linéairement ($r \approx \num{0,99}$) ; extrapoler cette droite prédirait un adulte de \SI{3}{m} à 30 ans ! Toute croissance réelle finit par saturer. Au Bac, signale toujours la prudence quand la valeur demandée sort de la plage observée : c'est un réflexe qui rapporte des points.
\end{attention}

\begin{popfigure}
\begin{tikzpicture}
\begin{axis}[
  width=\linewidth, height=7cm,
  xlabel={Trafic $x$ (centaines de véhicules)},
  ylabel={Ventes $y$ (centaines de litres)},
  xmin=0, xmax=16, ymin=0, ymax=14,
  grid=both, grid style={popline!40},
  legend style={at={(0.02,0.98)}, anchor=north west, font=\small}
]
\fill[popgreen!25] (axis cs:2,0) rectangle (axis cs:10,14);
\fill[poporange!25] (axis cs:10,0) rectangle (axis cs:16,14);
\addplot[only marks, mark=*, mark size=2pt, popblue] coordinates {(2,3.1) (3,4.2) (4,4.6) (5,5.9) (6,6.4) (7,7.6) (8,8.1) (9,9.3) (10,9.8)};
\addplot[popcurve, domain=1.5:10, samples=50]{0.847*x + 1.48};
\addlegendentry{Droite des moindres carrés (zone observée)}
\addplot[poporange, very thick, dashed, domain=10:13, samples=20]{0.847*x + 1.48};
\addlegendentry{Prolongement (extrapolation)}
\addplot[only marks, mark=star, mark size=4pt, popink] coordinates {(12,11.64)};
\node[popgreen!60!black, font=\small\bfseries] at (axis cs:6,13) {Interpolation};
\node[poporange!80!black, font=\small\bfseries] at (axis cs:13,13) {Extrapolation};
\node[popink, font=\small, anchor=west] at (axis cs:12.2,11.5) {Prévision pour $x_0 = 12$};
\end{axis}
\end{tikzpicture}
\legende{Zone verte : interpolation à l'intérieur des données observées. Zone orange : extrapolation, où la droite prolongée en pointillés n'est qu'une hypothèse. La droite tracée correspond à la régression des moindres carrés recalculée exactement depuis le tableau.}
\end{popfigure}

\courssub{Méthode complète et exemple chiffré}

\begin{methode}[Effectuer une prévision à l'aide d'un ajustement affine]
\begin{enumerate}
\item \textbf{Identifier la variable à estimer} : si l'on cherche $y$ connaissant $x_0$, on utilise la droite de régression de $y$ en $x$ ; si l'on cherche $x$ connaissant $y_0$, on utilise la droite de régression de $x$ en $y$.
\item \textbf{Vérifier la qualité de l'ajustement} : calculer $r$ et s'assurer que $|r|$ est proche de $1$ (au moins $\num{0,87}$).
\item \textbf{Situer la valeur} : comparer $x_0$ à la plage $[x_{\min}\,;\,x_{\max}]$ pour distinguer interpolation et extrapolation.
\item \textbf{Substituer} : calculer $\widehat{y_0} = a x_0 + b$ (ou $\widehat{x_0} = a' y_0 + b'$).
\item \textbf{Conclure en rédigeant} : donner la valeur estimée, son unité, et commenter sa fiabilité (interpolation/extrapolation, valeur de $r$).
\end{enumerate}
\end{methode}

\begin{exoresolu}[Prévision des ventes de carburant]
Une station-service de la nationale N1 observe, pendant 9 jours, le trafic $x$ (en centaines de véhicules) et les ventes $y$ (en centaines de litres) :
\begin{center}
\begin{tabular}{|c|*{9}{c|}}
\hline
\rowcolor{popblueL}
$x$ & 2 & 3 & 4 & 5 & 6 & 7 & 8 & 9 & 10 \\
\hline
$y$ & 3,1 & 4,2 & 4,6 & 5,9 & 6,4 & 7,6 & 8,1 & 9,3 & 9,8 \\
\hline
\end{tabular}
\end{center}
On calcule : $\bar{x} = 6$, $\bar{y} = \frac{59}{9} \approx \num{6,56}$, $V(x) = \frac{60}{9} \approx \num{6,67}$, $\mathrm{cov}(x,y) = \frac{50,8}{9} \approx \num{5,64}$, $r \approx \num{0,994}$.
\begin{enumerate}
\item Déterminer la droite de régression de $y$ en $x$ (moindres carrés) et la droite de Mayer. Estimer les ventes pour un trafic de \num{1200} véhicules avec les deux droites. Comparer.
\item Discuter la fiabilité de cette prévision.
\end{enumerate}
\tcblower
\textbf{1. Droite des moindres carrés (régression de $y$ en $x$).}
\[ a = \frac{\mathrm{cov}(x,y)}{V(x)} = \frac{50,8/9}{60/9} = \frac{50,8}{60} \approx \num{0,847}, \qquad b = \bar{y} - a\bar{x} \approx \num{6,56} - \num{0,847}\times 6 \approx \num{1,48}. \]
Droite : $y = \num{0,847}x + \num{1,48}$.

\textbf{Droite de Mayer.} On partage la série en deux groupes de même effectif (ou presque) autour de la médiane de $x$ (ici $x=6$).
Groupe 1 ($x \le 6$) : 5 points. $G_1(\bar{x}_1\,;\bar{y}_1) = (4\,;\, \frac{3,1+4,2+4,6+5,9+6,4}{5}) = (4\,;\, 4,84)$.
Groupe 2 ($x > 6$) : 4 points. $G_2(\bar{x}_2\,;\bar{y}_2) = (8,5\,;\, \frac{7,6+8,1+9,3+9,8}{4}) = (8,5\,;\, 8,7)$.
\[ a_M = \frac{8,7 - 4,84}{8,5 - 4} = \frac{3,86}{4,5} \approx \num{0,858}, \qquad b_M = 4,84 - \num{0,858}\times 4 \approx \num{1,41}. \]
Droite : $y = \num{0,858}x + \num{1,41}$.

\textbf{Prévision pour $x_0 = 12$ (soit \num{1200} véhicules).}
\begin{align*}
\text{Moindres carrés : } & \widehat{y_0} = \num{0,847}\times 12 + \num{1,48} \approx \num{11,64} \quad (\text{soit } \num{1164} \text{ litres}), \\
\text{Mayer : } & \widehat{y_0} = \num{0,858}\times 12 + \num{1,41} \approx \num{11,71} \quad (\text{soit } \num{1171} \text{ litres}).
\end{align*}
Les deux méthodes donnent des prévisions très proches (écart $< 1\,\%$), ce qui conforte la robustesse de l'estimation.

\textbf{2.} Fiabilité : $r \approx \num{0,994}$, corrélation quasi parfaite, donc l'ajustement affine est excellent sur la zone observée. En revanche $x_0 = 12$ est \emph{hors} de la plage observée $[2\,;\,10]$ : c'est une \textbf{extrapolation}. Elle reste proche des données (2 centaines au-delà), donc acceptable, mais le gérant commandera avec une marge de sécurité.
\end{exoresolu}

\courssub{Complément Bac : ajustements non affines par changement de variable}

Certains nuages ont une allure visiblement \emph{courbe} (croissance exponentielle d'une population, décroissance radioactive, saturation\dots). L'idée, classique au Bac, est de \cle{linéariser} : poser une nouvelle variable qui rend le nuage affine, ajuster, puis revenir à la variable initiale.

\begin{methode}[Linéarisation par changement de variable]
\begin{enumerate}
\item Observer l'allure du nuage : croissance de plus en plus rapide $\Rightarrow$ modèle exponentiel $y = k\,e^{cx}$, on pose $z = \ln y$ ; nuage en « racine » $\Rightarrow$ $z = \sqrt{y}$ ; hyperbole $\Rightarrow$ $z = \dfrac{1}{x}$.
\item Construire la série $(x_i\,; z_i)$ et vérifier que son nuage est affine (calculer le nouveau $r$).
\item Ajuster $z$ en $x$ par les moindres carrés : $z = ax + b$.
\item Revenir à $y$ : par exemple, si $z = \ln y$, alors $\ln y = ax + b$ donne $y = e^{b}\,e^{ax}$.
\item Utiliser ce modèle pour la prévision demandée.
\end{enumerate}
\end{methode}

\begin{exemplebox}[Ajustement exponentiel d'une population bactérienne]
Une culture compte $y$ bactéries (en milliers) après $x$ heures : $(0\,; 5)$, $(1\,; 8)$, $(2\,; 13)$, $(3\,; 21)$, $(4\,; 34)$. Le nuage n'est pas affine. On pose $z = \ln y$ : la série $(x\,; z)$ donne $z \approx \num{0,48} x + \num{1,60}$ avec $r \approx \num{0,999}$. Le modèle est donc $y = e^{\num{1,60}}\,e^{\num{0,48}x} \approx \num{4,95}\,e^{\num{0,48}x}$. Prévision à 6 heures : $y \approx \num{4,95}\,e^{\num{2,88}} \approx 88$ milliers de bactéries.
\end{exemplebox}

\courssub{Synthèse : l'étude complète d'une série à deux variables}

\begin{aretenir}[Fiche-méthode : du tableau à la prévision]
\begin{enumerate}
\item \textbf{Nuage de points} : représenter les points $M_i(x_i\,; y_i)$, observer l'allure (affine ? croissante ?).
\item \textbf{Point moyen} : calculer $G(\bar{x}\,; \bar{y})$ et le placer.
\item \textbf{Covariance} : $\mathrm{cov}(x,y) = \dfrac{1}{n}\sum x_i y_i - \bar{x}\,\bar{y}$ ; son signe donne le sens de la relation.
\item \textbf{Droite d'ajustement} : moindres carrés $a = \dfrac{\mathrm{cov}(x,y)}{V(x)}$, $b = \bar{y} - a\bar{x}$ (ou Mayer si l'énoncé l'impose) ; vérifier que $G$ appartient à la droite.
\item \textbf{Coefficient de corrélation} : $r = \dfrac{\mathrm{cov}(x,y)}{\sigma_x\,\sigma_y}$ ; valider l'ajustement si $|r| \geq \num{0,87}$.
\item \textbf{Prévision} : substituer dans la \emph{bonne} droite, puis commenter (interpolation ou extrapolation, fiabilité).
\end{enumerate}
\end{aretenir}

\begin{popfigure}
\begin{tikzpicture}[
  node distance=0.55cm,
  etape/.style={rectangle, rounded corners=3pt, draw=popblue, fill=popblueL, thick, text width=4.6cm, align=center, font=\small, minimum height=0.9cm},
  decision/.style={diamond, draw=poporange, fill=poporangeL, thick, aspect=2.2, align=center, font=\small, inner sep=1pt},
  fleche/.style={-{Stealth[length=2.5mm]}, thick, popdark}
]
\node[etape, align=center] (n1) {1. Nuage de points\\ (allure du phénomène)};
\node[etape, below=of n1] (n2) {2. Point moyen $G(\bar{x}\,;\bar{y})$};
\node[etape, below=of n2, align=center] (n3) {3. Covariance et droite d'ajustement\\ (moindres carrés ou Mayer)};
\node[etape, below=of n3] (n4) {4. Coefficient de corrélation $r$};
\node[decision, below=of n4] (n5) {$|r| \geq \num{0,87}$ ?};
\node[etape, below=1.1cm of n5, fill=popgreenL, draw=popgreen, align=center] (n6) {5. Prévision par substitution\\ + commentaire de fiabilité};
\node[etape, right=2.2cm of n5, fill=poppinkL, draw=poppink, align=center] (n7) {Ajustement affine non pertinent\\ (changement de variable ?)};
\draw[fleche] (n1) -- (n2);
\draw[fleche] (n2) -- (n3);
\draw[fleche] (n3) -- (n4);
\draw[fleche] (n4) -- (n5);
\draw[fleche] (n5) -- node[right, font=\small]{oui} (n6);
\draw[fleche] (n5) -- node[above, font=\small]{non} (n7);
\end{tikzpicture}
\legende{Organigramme de synthèse : les étapes de l'étude complète d'une série statistique à deux variables.}
\end{popfigure}

\begin{savaistu}[Des régressions qui façonnent le pays]
Les prévisions météorologiques qui guident les semis des agriculteurs de l'Adamaoua reposent sur des modèles de régression alimentés par des décennies d'observations. De même, l'Institut National de la Statistique (INS) du Cameroun produit des projections de population — essentielles pour planifier écoles, hôpitaux et électrification rurale — en ajustant des modèles de croissance sur les recensements passés. La droite des moindres carrés, publiée par Legendre en 1805, reste l'un des outils mathématiques les plus utilisés au monde : économie, médecine, sport, télécommunications\dots{} Tu viens d'en acquérir la maîtrise complète !
\end{savaistu}

\begin{exercice}[Pour s'entraîner]
Une entreprise de distribution d'eau relève la consommation $y$ (en m$^3$) d'un quartier en fonction du nombre $x$ de foyers raccordés, pour $x$ variant de 20 à 100. On obtient $\bar{x} = 60$, $\bar{y} = 480$, $\mathrm{cov}(x,y) = \num{3900}$, $V(x) = 500$, $\sigma_y = 200$.
\begin{enumerate}
\item Déterminer la droite de régression de $y$ en $x$.
\item Calculer $r$ et juger la qualité de l'ajustement.
\item Estimer la consommation pour 130 foyers raccordés. Cette prévision est-elle fiable ?
\end{enumerate}
\end{exercice}

\begin{corrige}[Exercice 1]
\textbf{1.} $a = \dfrac{\num{3900}}{500} = \num{7,8}$ et $b = 480 - \num{7,8}\times 60 = 12$. La droite est $y = \num{7,8}x + 12$.
\textbf{2.} $\sigma_x = \sqrt{500} \approx \num{22,36}$, d'où $r = \dfrac{\num{3900}}{\num{22,36}\times 200} = \dfrac{\num{3900}}{\num{4472}} \approx \num{0,872}$. L'ajustement est acceptable ($|r| \geq \num{0,87}$).
\textbf{3.} $\widehat{y} = \num{7,8}\times 130 + 12 = \num{1026}$ m$^3$. Mais $x = 130$ est hors de la plage $[20\,;\,100]$ : extrapolation, à annoncer avec prudence.
\end{corrige}

\newpage

\coursec{Activité d'intégration}

\courssub{Objectifs de l'activité}

À l'issue de cette activité d'intégration, tu seras capable de mobiliser, dans une situation authentique et non guidée pas à pas, l'ensemble de la chaîne statistique étudiée dans cette leçon :

\begin{itemize}
\item représenter une série statistique à deux variables par un \cle{nuage de points} et y placer le \cle{point moyen} ;
\item déterminer une droite d'ajustement par la \cle{méthode de Mayer} ;
\item calculer une \cle{covariance}, des \cle{variances} et la \cle{droite de régression} de $y$ en $x$ par la méthode des moindres carrés ;
\item calculer et interpréter le \cle{coefficient de corrélation linéaire} ;
\item utiliser un ajustement affine pour effectuer une \cle{prévision}, comparer deux méthodes et rédiger une conclusion argumentée à destination d'un décideur.
\end{itemize}

Au-delà des calculs, cette activité développe des savoir-être essentiels : la rigueur dans la tenue d'un tableau de calcul, l'esprit critique face à une prévision chiffrée, et la capacité à communiquer un résultat mathématique en langage courant.

\courssub{Situation}

La CAMWATER, société en charge de la production et de la distribution de l'eau potable au Cameroun, doit anticiper la demande du quartier Tougang, à Bafoussam (région de l'Ouest). L'expérience montre que la consommation d'eau augmente avec la chaleur : les gestionnaires souhaitent donc exploiter les relevés de l'année écoulée pour prévoir la consommation en fonction de la température.

Le chef de centre CAMWATER de Bafoussam a relevé, pendant dix mois consécutifs, la température moyenne mensuelle $x_i$ (en degrés Celsius) et la consommation totale d'eau du quartier $y_i$ (en milliers de m$^3$) :

\begin{center}
\begin{tabularx}{\linewidth}{|l|*{10}{>{\centering\arraybackslash}X|}}
\hline
\rowcolor{popblueL}
Mois & 1 & 2 & 3 & 4 & 5 & 6 & 7 & 8 & 9 & 10 \\
\hline
$x_i$ (°C) & 24 & 25 & 26 & 27 & 28 & 26 & 25 & 27 & 28 & 29 \\
\hline
$y_i$ (milliers de m$^3$) & 18 & 19 & 20 & 22 & 23 & 21 & 19 & 22 & 24 & 25 \\
\hline
\end{tabularx}
\end{center}

Les services météorologiques annoncent pour le mois prochain une température moyenne de $29$ °C. Le chef de centre se pose la question suivante : \emph{peut-on prévoir raisonnablement la consommation du quartier le mois prochain, et cette prévision est-elle assez fiable pour dimensionner le nouveau château d'eau ?}

Par groupes de quatre élèves, vous êtes les experts-statisticiens missionnés pour répondre à cette question et rédiger une note de synthèse.

\courssub{Consignes}

\begin{enumerate}
\item \textbf{Observation graphique.} Construire le nuage de points $M_i(x_i\,;\,y_i)$ dans un repère orthogonal bien choisi (unités adaptées aux données). Calculer les coordonnées du point moyen $G$ et le placer sur le graphique. Le nuage suggère-t-il un ajustement affine ? Justifier.
\item \textbf{Méthode de Mayer.} Partager le nuage en deux sous-nuages de cinq points (les cinq premiers mois, puis les cinq derniers). Déterminer les points moyens $G_1$ et $G_2$, puis l'équation réduite de la droite de Mayer $(G_1G_2)$.
\item \textbf{Méthode des moindres carrés.} À l'aide d'un tableau à cinq colonnes $(x_i-\bar{x})$, $(y_i-\bar{y})$, $(x_i-\bar{x})(y_i-\bar{y})$, $(x_i-\bar{x})^2$, $(y_i-\bar{y})^2$, calculer $\mathrm{cov}(X,Y)$, $V(X)$ et $V(Y)$. En déduire l'équation de la droite de régression de $y$ en $x$. Vérifier que cette droite passe par $G$.
\item \textbf{Qualité de l'ajustement.} Calculer le coefficient de corrélation linéaire $r$ (arrondi à $10^{-3}$). L'ajustement affine est-il justifié ?
\item \textbf{Prévision.} En utilisant chacune des deux droites obtenues, prévoir la consommation d'eau pour une température moyenne de $29$ °C. Comparer les deux prévisions et commenter l'écart observé.
\item \textbf{Note de synthèse.} Rédiger en cinq à huit lignes une note à l'intention du chef de centre CAMWATER : donner la prévision retenue, préciser le degré de confiance accordé au modèle, et indiquer s'il est raisonnable de dimensionner le château d'eau sur cette seule prévision (penser aux limites du modèle : marge de sécurité, croissance démographique du quartier, fuites du réseau).
\end{enumerate}

\courssub{Ressources à mobiliser}

\begin{aretenir}[Boîte à outils de la leçon]
\begin{itemize}
\item \textbf{Point moyen :} $G\left(\bar{x}\,;\,\bar{y}\right)$ avec $\bar{x}=\dfrac{1}{n}\sum x_i$ et $\bar{y}=\dfrac{1}{n}\sum y_i$.
\item \textbf{Méthode de Mayer :} partager le nuage en deux sous-nuages de même effectif (ou presque), calculer $G_1$ et $G_2$ ; la droite $(G_1G_2)$ est la droite d'ajustement.
\item \textbf{Covariance et variance :} $\mathrm{cov}(X,Y)=\dfrac{1}{n}\sum (x_i-\bar{x})(y_i-\bar{y})$ et $V(X)=\dfrac{1}{n}\sum (x_i-\bar{x})^2$.
\item \textbf{Droite de régression de $y$ en $x$ :} $y=ax+b$ avec $a=\dfrac{\mathrm{cov}(X,Y)}{V(X)}$ et $b=\bar{y}-a\bar{x}$. Elle passe toujours par $G$.
\item \textbf{Coefficient de corrélation :} $r=\dfrac{\mathrm{cov}(X,Y)}{\sqrt{V(X)\,V(Y)}}$, avec $-1\leq r\leq 1$. L'ajustement affine est jugé pertinent lorsque $|r|$ est proche de $1$ (en pratique $|r|\geq 0{,}87$ environ).
\item \textbf{Prévision :} on remplace $x$ par la valeur annoncée dans l'équation de la droite d'ajustement, en restant vigilant si cette valeur sort de la plage des données (extrapolation risquée).
\end{itemize}
\end{aretenir}

\courssub{Démarche et corrigé commenté}

\begin{exoresolu}[Corrigé complet de l'activité]
Prévoir la consommation d'eau du quartier Tougang à partir des dix relevés CAMWATER, par la méthode de Mayer puis par les moindres carrés, juger la qualité de l'ajustement et prévoir la consommation pour $29$ °C.
\tcblower
\textbf{Étape 1 : nuage de points et point moyen.}

On calcule les moyennes :
\[
\bar{x}=\frac{24+25+26+27+28+26+25+27+28+29}{10}=\frac{265}{10}=26{,}5
\]
\[
\bar{y}=\frac{18+19+20+22+23+21+19+22+24+25}{10}=\frac{213}{10}=21{,}3
\]
Le point moyen est $G(26{,}5\,;\,21{,}3)$. Sur le graphique (abscisses de $23$ à $30$, ordonnées de $17$ à $26$), les points sont sensiblement alignés le long d'une direction montante : un ajustement affine semble pertinent, ce que confirmera le calcul de $r$.

\textbf{Étape 2 : droite de Mayer.}

Sous-nuage 1 (mois 1 à 5) :
\[
\bar{x}_1=\frac{24+25+26+27+28}{5}=26 \quad ; \quad \bar{y}_1=\frac{18+19+20+22+23}{5}=20{,}4
\]
donc $G_1(26\,;\,20{,}4)$. Sous-nuage 2 (mois 6 à 10) :
\[
\bar{x}_2=\frac{26+25+27+28+29}{5}=27 \quad ; \quad \bar{y}_2=\frac{21+19+22+24+25}{5}=22{,}2
\]
donc $G_2(27\,;\,22{,}2)$. Le coefficient directeur de $(G_1G_2)$ est :
\[
a_M=\frac{22{,}2-20{,}4}{27-26}=1{,}8
\]
L'équation s'écrit $y-20{,}4=1{,}8(x-26)$, soit :
\[
\boxed{y=1{,}8x-26{,}4}
\]

\textbf{Étape 3 : moindres carrés.}

On complète le tableau à cinq colonnes (les écarts sont calculés par rapport à $\bar{x}=26{,}5$ et $\bar{y}=21{,}3$) :

\begin{center}
\begin{tabular}{|c|c|c|c|c|}
\hline
\rowcolor{popblueL}
$x_i-\bar{x}$ & $y_i-\bar{y}$ & $(x_i-\bar{x})(y_i-\bar{y})$ & $(x_i-\bar{x})^2$ & $(y_i-\bar{y})^2$ \\
\hline
$-2{,}5$ & $-3{,}3$ & $8{,}25$ & $6{,}25$ & $10{,}89$ \\
$-1{,}5$ & $-2{,}3$ & $3{,}45$ & $2{,}25$ & $5{,}29$ \\
$-0{,}5$ & $-1{,}3$ & $0{,}65$ & $0{,}25$ & $1{,}69$ \\
$0{,}5$ & $0{,}7$ & $0{,}35$ & $0{,}25$ & $0{,}49$ \\
$1{,}5$ & $1{,}7$ & $2{,}55$ & $2{,}25$ & $2{,}89$ \\
$-0{,}5$ & $-0{,}3$ & $0{,}15$ & $0{,}25$ & $0{,}09$ \\
$-1{,}5$ & $-2{,}3$ & $3{,}45$ & $2{,}25$ & $5{,}29$ \\
$0{,}5$ & $0{,}7$ & $0{,}35$ & $0{,}25$ & $0{,}49$ \\
$1{,}5$ & $2{,}7$ & $4{,}05$ & $2{,}25$ & $7{,}29$ \\
$2{,}5$ & $3{,}7$ & $9{,}25$ & $6{,}25$ & $13{,}69$ \\
\hline
\rowcolor{popgreenL}
\multicolumn{2}{|c|}{Sommes} & $32{,}5$ & $22{,}5$ & $48{,}1$ \\
\hline
\end{tabular}
\end{center}

On en déduit :
\[
\mathrm{cov}(X,Y)=\frac{32{,}5}{10}=3{,}25 \quad ; \quad V(X)=\frac{22{,}5}{10}=2{,}25 \quad ; \quad V(Y)=\frac{48{,}1}{10}=4{,}81
\]
Coefficient directeur de la droite de régression de $y$ en $x$ :
\[
a=\frac{\mathrm{cov}(X,Y)}{V(X)}=\frac{3{,}25}{2{,}25}=\frac{13}{9}\approx 1{,}444
\]
puis $b=\bar{y}-a\bar{x}=21{,}3-\dfrac{13}{9}\times 26{,}5\approx 21{,}3-38{,}278=-16{,}978$. D'où :
\[
\boxed{y\approx 1{,}444x-16{,}978}
\]
Vérification : pour $x=\bar{x}=26{,}5$, on obtient $y\approx 1{,}444\times 26{,}5-16{,}978\approx 21{,}3=\bar{y}$ : la droite passe bien par $G$.

\textbf{Étape 4 : coefficient de corrélation.}
\[
r=\frac{\mathrm{cov}(X,Y)}{\sqrt{V(X)V(Y)}}=\frac{3{,}25}{\sqrt{2{,}25\times 4{,}81}}=\frac{3{,}25}{\sqrt{10{,}8225}}\approx \frac{3{,}25}{3{,}290}\approx 0{,}988
\]
$r\approx 0{,}988$ est très proche de $1$ : la corrélation linéaire est \textbf{excellente}, l'ajustement affine est pleinement justifié, et la relation est croissante ($r>0$) : plus il fait chaud, plus le quartier consomme d'eau.

\textbf{Étape 5 : prévisions pour $x=29$ °C.}

Avec la droite de Mayer : $y=1{,}8\times 29-26{,}4=25{,}8$ milliers de m$^3$.

Avec la droite des moindres carrés : $y\approx 1{,}444\times 29-16{,}978\approx 24{,}91$ milliers de m$^3$.

Les deux prévisions ($\num{25,8}$ et $\num{24,91}$ milliers de m$^3$) sont proches, à environ $0{,}9$ millier de m$^3$ près, ce qui renforce la confiance dans l'ordre de grandeur : environ \textbf{25 000 m$^3$}. On retient la prévision des moindres carrés, méthode qui minimise globalement les écarts sur l'ensemble des points.

\textbf{Étape 6 : note au chef de centre (exemple de rédaction).}

\emph{« Monsieur le Chef de centre, l'étude statistique des dix derniers mois montre une très forte corrélation linéaire ($r\approx 0{,}99$) entre la température moyenne et la consommation d'eau du quartier Tougang. Pour une température annoncée de $29$ °C, notre modèle prévoit une consommation d'environ 25 000 m$^3$ (entre 24 900 et 25 800 m$^3$ selon la méthode). Cette prévision est fiable car $29$ °C reste dans la plage des températures observées. Toutefois, nous recommandons de dimensionner le château d'eau avec une marge de sécurité d'au moins 10 à 15 \%, afin de tenir compte de la croissance démographique du quartier, des fuites du réseau et des écarts naturels autour de la droite d'ajustement. »}
\end{exoresolu}

\courssub{Bilan de l'activité}

Cette activité a permis de mettre en évidence plusieurs idées essentielles de la leçon :

\begin{itemize}
\item \textbf{La démarche statistique est une chaîne complète} : on ne calcule pas une droite de régression « à l'aveugle » ; on commence toujours par observer le nuage de points, et on termine par juger la qualité de l'ajustement grâce à $r$ avant toute prévision.
\item \textbf{Les deux méthodes d'ajustement se complètent} : la méthode de Mayer est rapide et géométrique, celle des moindres carrés est plus fine car elle utilise toutes les données ; quand leurs prévisions sont proches, la conclusion est renforcée.
\item \textbf{Le coefficient de corrélation est un juge de paix} : $r\approx 0{,}988$ autorise la prévision, alors qu'un $|r|$ faible l'aurait interdite, même si l'on peut toujours calculer mécaniquement une droite.
\item \textbf{Une prévision reste une estimation} : le statisticien honnête communique un ordre de grandeur et une marge de prudence, jamais une certitude. C'est cette honnêteté intellectuelle, autant que les formules, qui fait la valeur d'une étude statistique pour un décideur comme le chef de centre CAMWATER.
\end{itemize}

\begin{attention}[Ce qu'il fallait éviter]
Oublier de vérifier que la droite des moindres carrés passe par $G$ (excellent réflexe de contrôle) ; confondre $\mathrm{cov}(X,Y)$ et $V(X)$ dans le calcul de $a$ ; conclure à une « cause » entre température et consommation alors que la statistique ne met en évidence qu'une \emph{corrélation} ; enfin, extrapoler sans précaution à des températures très éloignées de la plage observée.
\end{attention}

\newpage

\coursec{Méthodes et exercices résolus}

\coursec{Statistiques à deux variables}

\courssub{Séries statistiques à deux caractères : tableau à double entrée et séries marginales}

\begin{definition}[Série statistique à deux caractères]
On appelle \cle{série statistique à deux caractères} (ou série double) une série où chaque individu de la population est caractérisé par deux variables $X$ et $Y$. On note $(x_i\,;\,y_i)_{1\le i\le n}$ les $n$ couples de valeurs observées.
\end{definition}

\begin{definition}[Tableau à double entrée et séries marginales]
Lorsque les effectifs sont importants, on regroupe les données dans un \cle{tableau à double entrée} (ou tableau de contingence pour des caractères qualitatifs, tableau de effectifs pour des caractères quantitatifs groupés en classes).
Les \cle{séries marginales} sont les séries statistiques de chaque variable prise séparément : on obtient leurs effectifs (ou fréquences) en sommant les lignes et les colonnes du tableau.
\end{definition}

\begin{methode}[Construire les séries marginales à partir d'un tableau à double entrée]
\begin{enumerate}
\item Repérer les modalités (ou classes) de $X$ en en-tête de colonnes et celles de $Y$ en en-tête de lignes.
\item Les effectifs \emph{marginaux} de $X$ (dernière ligne du tableau) s'obtiennent en sommant les effectifs de chaque colonne.
\item Les effectifs \emph{marginaux} de $Y$ (dernière colonne du tableau) s'obtiennent en sommant les effectifs de chaque ligne.
\item La somme des effectifs marginaux de $X$ (ou de $Y$) donne l'effectif total $n$.
\end{enumerate}
\end{methode}

\begin{exemplebox}[Ventes de téléphones (MTN/Orange) à Yaoundé]
Un magasin a relevé les ventes hebdomadaires selon le modèle (lignes) et la marque (colonnes) :
\begin{center}
\begin{tabular}{|c|ccc|c|}
\hline
\rowcolor{popblueL} \diagbox{$Y$ (Modèle)}{$X$ (Marque)} & \textbf{MTN} & \textbf{Orange} & \textbf{Autres} & \textbf{Total $Y$} \\
\hline
\textbf{Entrée de gamme} & 12 & 18 & 5 & 35 \\
\textbf{Milieu de gamme} & 20 & 25 & 10 & 55 \\
\textbf{Haut de gamme} & 8 & 12 & 5 & 25 \\
\hline
\textbf{Total $X$} & \textbf{40} & \textbf{55} & \textbf{20} & \textbf{115} \\
\hline
\end{tabular}
\end{center}
Les séries marginales sont :
\begin{itemize}
\item Pour la marque $X$ : $(40\,;\,55\,;\,20)$ effectif total $115$.
\item Pour le modèle $Y$ : $(35\,;\,55\,;\,25)$ effectif total $115$.
\end{itemize}
\end{exemplebox}

\courssub{Nuage de points et point moyen}

\begin{methode}[Nuage de points et point moyen]
Pour représenter une série statistique double $(x_i\,;\,y_i)$ et situer son point moyen :
\begin{enumerate}
\item Choisir deux axes gradués adaptés aux valeurs de $x$ et de $y$ (on peut commencer les graduations ailleurs qu'en $0$ si les valeurs sont concentrées ; indiquer alors une \cle{rupture d'échelle}).
\item Placer chaque point $M_i(x_i\,;\,y_i)$ : l'ensemble forme le \cle{nuage de points}.
\item Calculer les moyennes marginales :
\[ \bar{x}=\frac{1}{n}\sum_{i=1}^{n}x_i \qquad \text{et} \qquad \bar{y}=\frac{1}{n}\sum_{i=1}^{n}y_i. \]
\item Le \cle{point moyen} est $G(\bar{x}\,;\,\bar{y})$ : le placer sur le graphique (souvent d'une couleur différente).
\item Observer la forme du nuage : s'il est allongé autour d'une direction, un \cle{ajustement affine} est envisageable.
\end{enumerate}
\textbf{Réflexe Bac :} vérifier que $G$ est bien \og au centre \fg{} du nuage ; une erreur de calcul de $\bar{x}$ ou $\bar{y}$ se voit immédiatement sur le graphique.
\end{methode}

\begin{exoresolu}[Récolte de maïs à Bafoussam]
Un agronome a relevé, pour cinq exploitations de la région de l'Ouest, la surface cultivée $x_i$ (en hectares) et la récolte de maïs $y_i$ (en tonnes) :
\begin{center}
\begin{tabular}{|l|*{5}{c|}}
\hline
\rowcolor{popblueL} Surface $x_i$ (ha) & 2 & 4 & 6 & 8 & 10 \\
\hline
Récolte $y_i$ (t) & 3 & 5 & 7 & 8 & 10 \\
\hline
\end{tabular}
\end{center}
\begin{enumerate}
\item Représenter le nuage de points associé à cette série.
\item Déterminer les coordonnées du point moyen $G$ et le placer sur le graphique.
\end{enumerate}
\tcblower
\textbf{1. Nuage de points.} On place les points $M_1(2\,;\,3)$, $M_2(4\,;\,5)$, $M_3(6\,;\,7)$, $M_4(8\,;\,8)$, $M_5(10\,;\,10)$.
\begin{popfigure}
\begin{center}
\begin{tikzpicture}[scale=0.62]
\draw[->,thick] (0,0) -- (11.5,0) node[below right]{$x$ (ha)};
\draw[->,thick] (0,0) -- (0,11.5) node[above left]{$y$ (t)};
\foreach \x in {2,4,6,8,10} \draw (\x,0.08) -- (\x,-0.08) node[below]{\small \x};
\foreach \y in {2,4,6,8,10} \draw (0.08,\y) -- (-0.08,\y) node[left]{\small \y};
\foreach \x/\y in {2/3,4/5,6/7,8/8,10/10} \fill[popblue] (\x,\y) circle (0.12);
\fill[poporange] (6,6.6) circle (0.16) node[above right,popdark]{\small $G(6\,;\,6{,}6)$};
\end{tikzpicture}
\end{center}
\legende{Nuage de points et point moyen $G$}
\end{popfigure}

\textbf{2. Point moyen.} On calcule les moyennes marginales :
\[ \bar{x}=\frac{2+4+6+8+10}{5}=\frac{30}{5}=6 \qquad ; \qquad \bar{y}=\frac{3+5+7+8+10}{5}=\frac{33}{5}=6{,}6. \]
Le point moyen est donc $G(6\,;\,6{,}6)$. Il est bien situé au centre du nuage, qui apparaît allongé selon une direction montante : un ajustement affine est justifié visuellement.
\end{exoresolu}

\courssub{Ajustement affine par la méthode de Mayer}

\begin{methode}[Méthode de Mayer (double moyenne)]
\begin{enumerate}
\item \textbf{Trier} les points par abscisses croissantes (indispensable !).
\item \textbf{Partager} le nuage en deux sous-nuages de effectifs égaux (ou différant de $1$ si $n$ est impair).
\item Calculer les points moyens $G_1(x_1\,;\,y_1)$ et $G_2(x_2\,;\,y_2)$ de chaque sous-nuage.
\item La \cle{droite de Mayer} est la droite $(G_1G_2)$. Son équation s'obtient par :
\[ a=\frac{y_2-y_1}{x_2-x_1}, \qquad \text{puis} \qquad y=ax+b \ \text{avec}\ b=y_1-a\,x_1. \]
\item Vérifier que la droite passe \og au milieu \fg{} du nuage sur le graphique.
\end{enumerate}
\textbf{Réflexe :} la droite de Mayer passe aussi par le point moyen global $G$ quand les deux sous-nuages ont le même effectif : c'est un bon contrôle.
\end{methode}

\begin{exoresolu}[Droite de Mayer pour la récolte de maïs]
On reprend la série de l'exercice précédent :
\begin{center}
\begin{tabular}{|l|*{5}{c|}}
\hline
\rowcolor{popblueL} $x_i$ & 2 & 4 & 6 & 8 & 10 \\
\hline
$y_i$ & 3 & 5 & 7 & 8 & 10 \\
\hline
\end{tabular}
\end{center}
Déterminer, par la méthode de Mayer, une équation de la droite d'ajustement de $y$ en $x$.
\tcblower
\textbf{Étape 1 et 2.} Les points sont déjà rangés par abscisses croissantes. Avec $n=5$ (impair), on partage en un premier groupe de $2$ points et un second de $3$ points :
\[ \text{Groupe 1 : } (2\,;\,3),\ (4\,;\,5) \qquad ; \qquad \text{Groupe 2 : } (6\,;\,7),\ (8\,;\,8),\ (10\,;\,10). \]

\textbf{Étape 3.} Points moyens des sous-nuages :
\[ G_1\left(\frac{2+4}{2}\,;\,\frac{3+5}{2}\right)=G_1(3\,;\,4), \]
\[ G_2\left(\frac{6+8+10}{3}\,;\,\frac{7+8+10}{3}\right)=G_2\left(8\,;\,\frac{25}{3}\right). \]

\textbf{Étape 4.} Coefficient directeur de $(G_1G_2)$ :
\[ a=\frac{\dfrac{25}{3}-4}{8-3}=\frac{\dfrac{25-12}{3}}{5}=\frac{13}{15}\approx 0{,}87. \]
Ordonnée à l'origine, en utilisant $G_1$ :
\[ b=y_1-a x_1=4-\frac{13}{15}\times 3=4-\frac{13}{5}=\frac{7}{5}=1{,}4. \]

\textbf{Conclusion :} la droite de Mayer a pour équation
\[ y=\frac{13}{15}x+\frac{7}{5} \qquad \text{soit environ} \qquad y=0{,}87\,x+1{,}4. \]
\textbf{Vérification :} pour $G(6\,;\,6{,}6)$ : $\frac{13}{15}\times 6+\frac{7}{5}=\frac{26}{5}+\frac{7}{5}=\frac{33}{5}=6{,}6$. La droite passe bien par $G$ : le calcul est cohérent.
\end{exoresolu}

\courssub{Covariance de deux variables statistiques}

\begin{methode}[Covariance de deux variables statistiques]
\begin{enumerate}
\item Calculer les moyennes $\bar{x}$ et $\bar{y}$.
\item Appliquer l'une des deux formules équivalentes :
\[ \operatorname{cov}(x,y)=\frac{1}{n}\sum_{i=1}^{n}(x_i-\bar{x})(y_i-\bar{y}) \qquad \text{ou} \qquad \operatorname{cov}(x,y)=\frac{1}{n}\sum_{i=1}^{n}x_i y_i-\bar{x}\,\bar{y}. \]
La seconde (dite \emph{formule de König-Huygens}) est souvent plus rapide avec un tableau.
\item Construire un tableau à colonnes $x_i$, $y_i$, $x_i y_i$ (et $x_i^2$, $y_i^2$ si l'on vise ensuite les variances).
\item \textbf{Interprétation du signe :} $\operatorname{cov}(x,y)>0$ : $x$ et $y$ varient globalement dans le même sens ; $\operatorname{cov}(x,y)<0$ : en sens contraires ; $\operatorname{cov}(x,y)=0$ : pas de liaison affine (mais possible liaison non-linéaire).
\end{enumerate}
\textbf{Réflexe :} $\operatorname{cov}(x,x)=V(x)$ : la variance est un cas particulier de covariance.
\end{methode}

\begin{exoresolu}[Covariance surface/récolte]
On considère à nouveau la série :
\begin{center}
\begin{tabular}{|l|*{5}{c|}}
\hline
\rowcolor{popblueL} $x_i$ & 2 & 4 & 6 & 8 & 10 \\
\hline
$y_i$ & 3 & 5 & 7 & 8 & 10 \\
\hline
\end{tabular}
\end{center}
Calculer $\operatorname{cov}(x,y)$ par les deux méthodes et interpréter le signe du résultat.
\tcblower
On sait déjà que $\bar{x}=6$ et $\bar{y}=6{,}6$.

\textbf{Méthode 1 : formule de König-Huygens.} On complète le tableau :
\begin{center}
\begin{tabular}{|c|c|c|c|}
\hline
\rowcolor{popblueL} $x_i$ & $y_i$ & $x_i y_i$ & $x_i^2$ \\
\hline
2 & 3 & 6 & 4 \\
\hline
4 & 5 & 20 & 16 \\
\hline
6 & 7 & 42 & 36 \\
\hline
8 & 8 & 64 & 64 \\
\hline
10 & 10 & 100 & 100 \\
\hline
\rowcolor{poporangeL} \textbf{Sommes} & & \textbf{232} & \textbf{220} \\
\hline
\end{tabular}
\end{center}
\[ \operatorname{cov}(x,y)=\frac{232}{5}-6\times 6{,}6=46{,}4-39{,}6=6{,}8. \]

\textbf{Méthode 2 : écarts aux moyennes.}
\begin{align*}
\operatorname{cov}(x,y)&=\frac{1}{5}\big[(2-6)(3-6{,}6)+(4-6)(5-6{,}6)+(6-6)(7-6{,}6)\\
&\quad +(8-6)(8-6{,}6)+(10-6)(10-6{,}6)\big]\\
&=\frac{1}{5}\big[(-4)(-3{,}6)+(-2)(-1{,}6)+0\times 0{,}4+2\times 1{,}4+4\times 3{,}4\big]\\
&=\frac{1}{5}\big[14{,}4+3{,}2+0+2{,}8+13{,}6\big]=\frac{34}{5}=6{,}8.
\end{align*}

\textbf{Conclusion :} $\operatorname{cov}(x,y)=6{,}8>0$ : surface cultivée et récolte varient globalement dans le \textbf{même sens}, ce qui confirme l'observation du nuage.
\end{exoresolu}

\courssub{Droites de régression par la méthode des moindres carrés}

\begin{methode}[Droites de régression de $y$ en $x$ et de $x$ en $y$]
\begin{enumerate}
\item Calculer $\bar{x}$, $\bar{y}$, $\operatorname{cov}(x,y)$, $V(x)$ et $V(y)$ avec :
\[ V(x)=\frac{1}{n}\sum x_i^2-\bar{x}^2 \qquad ; \qquad V(y)=\frac{1}{n}\sum y_i^2-\bar{y}^2. \]
\item \textbf{Droite de régression de $y$ en $x$} (on prédit $y$ à partir de $x$) : $y=ax+b$ avec
\[ a=\frac{\operatorname{cov}(x,y)}{V(x)} \qquad \text{et} \qquad b=\bar{y}-a\bar{x}. \]
\item \textbf{Droite de régression de $x$ en $y$} (on prédit $x$ à partir de $y$) : $x=a'y+b'$ avec
\[ a'=\frac{\operatorname{cov}(x,y)}{V(y)} \qquad \text{et} \qquad b'=\bar{x}-a'\bar{y}. \]
\item \textbf{Propriété de contrôle :} les deux droites passent par le point moyen $G(\bar{x}\,;\,\bar{y})$.
\end{enumerate}
\begin{attention}[Ne pas confondre les deux droites]
La droite de $x$ en $y$ s'écrit $x=a'y+b'$ : pour la comparer à $y=ax+b$, il faut la réécrire $y=\dfrac{x-b'}{a'}$. Ce ne sont \textbf{pas} les mêmes droites, sauf corrélation parfaite ($|r|=1$).
\end{attention}
\end{methode}

\begin{exoresolu}[Droites des moindres carrés]
Toujours avec la série $(x_i\,;\,y_i)$ : $(2;3)$, $(4;5)$, $(6;7)$, $(8;8)$, $(10;10)$.
\begin{enumerate}
\item Déterminer l'équation de la droite de régression de $y$ en $x$.
\item Déterminer l'équation de la droite de régression de $x$ en $y$, et vérifier que les deux droites passent par $G$.
\end{enumerate}
\tcblower
On dispose déjà de : $\bar{x}=6$, $\bar{y}=6{,}6$, $\operatorname{cov}(x,y)=6{,}8$, $\sum x_i^2=220$. Il reste à calculer $\sum y_i^2$ :
\[ \sum y_i^2=9+25+49+64+100=247. \]

\textbf{Variances :}
\[ V(x)=\frac{220}{5}-6^2=44-36=8 \qquad ; \qquad V(y)=\frac{247}{5}-6{,}6^2=49{,}4-43{,}56=5{,}84. \]

\textbf{1. Droite de régression de $y$ en $x$.}
\[ a=\frac{\operatorname{cov}(x,y)}{V(x)}=\frac{6{,}8}{8}=0{,}85 \qquad ; \qquad b=\bar{y}-a\bar{x}=6{,}6-0{,}85\times 6=6{,}6-5{,}1=1{,}5. \]
\[ \boxed{y=0{,}85\,x+1{,}5} \]

\textbf{2. Droite de régression de $x$ en $y$.}
\[ a'=\frac{\operatorname{cov}(x,y)}{V(y)}=\frac{6{,}8}{5{,}84}\approx 1{,}164 \qquad ; \qquad b'=\bar{x}-a'\bar{y}=6-1{,}164\times 6{,}6\approx 6-7{,}685=-1{,}685. \]
\[ \boxed{x\approx 1{,}164\,y-1{,}685} \]

\textbf{Vérification du passage par $G(6\,;\,6{,}6)$ :}
\begin{itemize}
\item $0{,}85\times 6+1{,}5=5{,}1+1{,}5=6{,}6=\bar{y}$ \checkmark
\item $1{,}164\times 6{,}6-1{,}685\approx 7{,}685-1{,}685=6=\bar{x}$ \checkmark
\end{itemize}
Les deux droites concourent bien au point moyen $G$.
\end{exoresolu}

\courssub{Coefficient de corrélation linéaire et qualité d'un ajustement}

\begin{methode}[Coefficient de corrélation linéaire $r$]
\begin{enumerate}
\item Calculer $\operatorname{cov}(x,y)$, $\sigma_x=\sqrt{V(x)}$ et $\sigma_y=\sqrt{V(y)}$.
\item Appliquer la formule :
\[ r=\frac{\operatorname{cov}(x,y)}{\sigma_x\,\sigma_y}=\frac{\operatorname{cov}(x,y)}{\sqrt{V(x)\,V(y)}}. \]
\item \textbf{Interpréter} (on a toujours $-1\le r\le 1$) :
\begin{itemize}
\item $|r|$ proche de $1$ (en pratique $|r|\ge 0{,}87$, c'est-à-dire $r^2\ge 0{,}75$) : forte corrélation linéaire, l'ajustement affine est \textbf{pertinent} ;
\item $r$ proche de $0$ : pas de relation affine (l'ajustement affine est à rejeter) ;
\item le signe de $r$ est celui de la covariance : $r>0$ relation croissante, $r<0$ décroissante.
\end{itemize}
\item \textbf{Lien avec les droites :} $a\times a'=r^2$ (produit des coefficients directeurs des deux droites de régression).
\end{enumerate}
\end{methode}

\begin{exoresolu}[Qualité de l'ajustement surface/récolte]
Avec les résultats précédents : $\operatorname{cov}(x,y)=6{,}8$, $V(x)=8$, $V(y)=5{,}84$.
\begin{enumerate}
\item Calculer le coefficient de corrélation linéaire $r$ (arrondir à $10^{-3}$).
\item Un ajustement affine de cette série est-il justifié ?
\item Retrouver $r^2$ à partir des coefficients $a=0{,}85$ et $a'\approx 1{,}164$.
\end{enumerate}
\tcblower
\textbf{1. Calcul de $r$.}
\[ r=\frac{6{,}8}{\sqrt{8\times 5{,}84}}=\frac{6{,}8}{\sqrt{46{,}72}}\approx \frac{6{,}8}{6{,}835}\approx 0{,}995. \]

\textbf{2. Interprétation.} On a $r\approx 0{,}995$, donc $r^2\approx 0{,}99\ge 0{,}75$ (et même $|r|\ge 0{,}87$). La corrélation linéaire entre la surface cultivée et la récolte est \textbf{très forte et positive} : l'ajustement affine est pleinement justifié.

\textbf{3. Vérification par le produit $a\times a'$.}
\[ a\times a'=0{,}85\times 1{,}164\approx 0{,}990\approx r^2. \]
Le contrôle est concluant : les calculs des deux droites de régression sont cohérents avec le coefficient de corrélation.
\end{exoresolu}

\courssub{Utiliser un ajustement affine pour effectuer une prévision}

\begin{methode}[Prévision à partir d'une droite d'ajustement]
\begin{enumerate}
\item \textbf{Justifier d'abord} que la prévision est légitime : coefficient de corrélation satisfaisant ($|r|$ proche de $1$) et valeur de la variable dans (ou proche de) la plage des données observées.
\item Pour prévoir $y$ à partir d'une valeur $x_0$ : utiliser la droite de régression \textbf{de $y$ en $x$} : $y_0=a x_0+b$.
\item Pour prévoir $x$ à partir d'une valeur $y_0$ : utiliser la droite de régression \textbf{de $x$ en $y$} : $x_0=a' y_0+b'$.
\item Conclure par une phrase \textbf{en contexte}, avec l'unité, et mentionner qu'il s'agit d'une estimation.
\end{enumerate}
\textbf{Réflexe :} se méfier des extrapolations lointaines (prévoir très loin en dehors des données observées est hasardeux, même avec un bon $r$).
\end{methode}

\begin{exoresolu}[Prévision de récolte et de surface]
On admet que la droite de régression de $y$ en $x$ est $y=0{,}85x+1{,}5$ et celle de $x$ en $y$ est $x=1{,}164y-1{,}685$ (avec $r\approx 0{,}995$).
\begin{enumerate}
\item Estimer la récolte attendue pour une exploitation de $12$ hectares.
\item Un exploitant vise une récolte de $15$ tonnes. Quelle surface devrait-il cultiver ?
\end{enumerate}
\tcblower
\textbf{Justification préalable :} $r\approx 0{,}995$, très proche de $1$ : l'ajustement affine est de bonne qualité, les prévisions sont donc pertinentes.

\textbf{1. Prévision de la récolte pour $x_0=12$ ha.} On utilise la droite de régression \textbf{de $y$ en $x$} :
\[ y_0=0{,}85\times 12+1{,}5=10{,}2+1{,}5=11{,}7. \]
Pour une exploitation de $12$ hectares, on peut estimer la récolte à environ $\mathbf{11{,}7}$ \textbf{tonnes} de maïs. (La valeur $12$ est proche des données observées, l'estimation est raisonnable.)

\textbf{2. Surface nécessaire pour $y_0=15$ tonnes.} On utilise la droite de régression \textbf{de $x$ en $y$} :
\[ x_0=1{,}164\times 15-1{,}685=17{,}46-1{,}685=15{,}775\approx 15{,}8. \]
Pour espérer une récolte de $15$ tonnes, l'exploitant devrait cultiver environ $\mathbf{15{,}8}$ \textbf{hectares}. Il s'agit d'une extrapolation au-delà des données ($y$ maximal observé : $10$ t) : le résultat est à prendre avec prudence.
\end{exoresolu}

\begin{attention}[Les 5 erreurs qui coûtent des points au Bac]
\begin{enumerate}
\item \textbf{Oublier de trier les points pour Mayer.} La méthode de Mayer exige des points rangés par \emph{abscisses croissantes} avant le partage en deux sous-nuages. Un partage dans l'ordre du tableau non trié donne une droite fausse.
\item \textbf{Confondre les deux droites de régression.} $y=ax+b$ sert à prévoir $y$ connaissant $x$ ; $x=a'y+b'$ sert à prévoir $x$ connaissant $y$. Utiliser la mauvaise droite pour une prévision est une erreur éliminatoire. Et on ne passe pas de l'une à l'autre par une simple inversion !
\item \textbf{Se tromper de dénominateur dans les coefficients.} $a=\dfrac{\operatorname{cov}(x,y)}{V(x)}$ mais $a'=\dfrac{\operatorname{cov}(x,y)}{V(y)}$. Mélanger $V(x)$ et $V(y)$ fausse tout le reste. Contrôle possible : $a\times a'=r^2$ et les deux droites passent par $G(\bar{x}\,;\,\bar{y})$.
\item \textbf{Interpréter $r$ sans le justifier, ou mal.} Il ne suffit pas d'écrire \og $r=0{,}99$ \fg{} : il faut conclure explicitement (\og $|r|$ proche de $1$, donc forte corrélation linéaire, l'ajustement affine est pertinent \fg). Rappel : $r$ proche de $0$ ne signifie pas \og pas de relation \fg, mais \og pas de relation \emph{affine} \fg.
\item \textbf{Erreurs de calcul sur les moyennes et la covariance.} Oublier de diviser par $n$, confondre $\overline{xy}$ et $\bar{x}\bar{y}$ dans la formule $\operatorname{cov}(x,y)=\overline{xy}-\bar{x}\bar{y}$, ou perdre des signes dans les écarts $(x_i-\bar{x})(y_i-\bar{y})$. Rédiger un tableau à colonnes $x_i$, $y_i$, $x_iy_i$, $x_i^2$, $y_i^2$ sécurise tous les calculs.
\end{enumerate}
\end{attention}

\begin{savaistu}[Le coefficient de corrélation et la causalité]
Un coefficient de corrélation $r$ proche de $1$ ou $-1$ prouve une \emph{liaison affine forte}, mais \textbf{pas} une relation de cause à effet. Exemple : au Cameroun, le nombre d'abonnés MTN et la production de cacao ont tous deux augmenté ces 20 dernières années ($r\approx 0,95$). Cela ne veut pas dire que téléphoner fait pousser le cacao ! Les deux sont liés à une troisième variable : le développement économique et démographique. Au Bac, on écrit : \og forte corrélation linéaire \fg, jamais \og $x$ cause $y$ \fg.
\end{savaistu}

\begin{exercice}[Entraînement Bac -- Consommation d'électricité à Douala]
Le tableau ci-dessous donne la consommation mensuelle d'électricité $y$ (en kWh) en fonction du nombre d'appareils électroménagers $x$ dans 10 foyers doualais.
\begin{center}
\begin{tabular}{|c|*{10}{c|}}
\hline
\rowcolor{popblueL} $x_i$ & 3 & 4 & 5 & 5 & 6 & 7 & 7 & 8 & 9 & 10 \\
\hline
$y_i$ & 120 & 150 & 180 & 190 & 210 & 240 & 250 & 280 & 300 & 330 \\
\hline
\end{tabular}
\end{center}
\begin{enumerate}
\item Représenter le nuage de points et le point moyen $G$.
\item Déterminer la droite de Mayer.
\item Calculer $\operatorname{cov}(x,y)$, $V(x)$, $V(y)$ et les deux droites de régression.
\item Calculer $r$ et conclure sur la pertinence d'un ajustement affine.
\item Prévoir la consommation pour un foyer possédant 11 appareils. Est-ce une interpolation ou une extrapolation ?
\end{enumerate}
\end{exercice}

\begin{corrige}[Exercice Entraînement Bac]
\textbf{1.} $\bar{x}=\frac{64}{10}=6{,}4$ ; $\bar{y}=\frac{2250}{10}=225$. $G(6{,}4\,;\,225)$. Nuage allongé croissant.
\textbf{2.} $n=10$ pair. Groupe 1 (5 premiers) : $G_1(4{,}4\,;\,170)$. Groupe 2 (5 derniers) : $G_2(8{,}2\,;\,280)$. $a=\frac{110}{3,8}\approx 28,95$. $b=170-28,95\times 4,4\approx 42,6$. Droite de Mayer : $y\approx 28,95x+42,6$.
\textbf{3.} $\sum x_i^2=430$, $\sum y_i^2=537\,500$, $\sum x_iy_i=15\,380$.
$\operatorname{cov}=\frac{15380}{10}-6,4\times 225=1538-1440=98$.
$V(x)=\frac{430}{10}-6,4^2=43-40,96=2,04$.
$V(y)=\frac{537500}{10}-225^2=53750-50625=3125$.
$a=\frac{98}{2,04}\approx 48,04$ ; $b=225-48,04\times 6,4\approx -82,46$. Droite $y$ en $x$ : $y\approx 48,04x-82,46$.
$a'=\frac{98}{3125}=0,03136$ ; $b'=6,4-0,03136\times 225\approx -0,656$. Droite $x$ en $y$ : $x\approx 0,03136y-0,656$.
\textbf{4.} $r=\frac{98}{\sqrt{2,04\times 3125}}=\frac{98}{\sqrt{6375}}\approx \frac{98}{79,84}\approx 0,977$. $r^2\approx 0,955\ge 0,75$. Ajustement très pertinent.
\textbf{5.} $y_0=48,04\times 11-82,46\approx 446$ kWh. Extrapolation (11 > 10 max observé), prudence.
\end{corrige}

\newpage

\coursec{Exercices}

\coursec{Statistiques à deux variables}

\courssub{Séries statistiques à deux caractères}

\begin{definition}[Série statistique à deux caractères]
On appelle \cle{série statistique à deux caractères} (ou série double) une série où chaque individu est caractérisé par deux variables $X$ et $Y$. Elle se présente sous la forme d'un tableau à double entrée ou d'une liste de couples $(x_i\,;\,y_i)$ avec des effectifs $n_i$ (souvent $n_i=1$).
\end{definition}

\begin{definition}[Séries marginales]
Soit un tableau à double entrée croisant les modalités (ou classes) de $X$ et $Y$.
\begin{itemize}
\item La \cle{série marginale de $X$} est la série obtenue en sommant les effectifs de chaque ligne (ou colonne selon l'orientation) : elle donne la distribution de $X$ seule.
\item La \cle{série marginale de $Y$} est la série obtenue en sommant les effectifs de chaque colonne (ou ligne) : elle donne la distribution de $Y$ seule.
\end{itemize}
Les effectifs marginaux sont notés $n_{i\cdot}$ pour $X$ et $n_{\cdot j}$ pour $Y$. L'effectif total est $n = \sum n_{i\cdot} = \sum n_{\cdot j}$.
\end{definition}

\begin{propriete}[Moyennes et variances marginales]
Les moyennes $\bar{x}$, $\bar{y}$ et les variances $V(X)$, $V(Y)$ se calculent \emph{uniquement} à partir des séries marginales, exactement comme pour une série à une variable.
\[
\bar{x} = \frac{1}{n}\sum n_{i\cdot}x_i,\quad V(X) = \frac{1}{n}\sum n_{i\cdot}(x_i-\bar{x})^2 = \overline{x^2}-\bar{x}^2
\]
Idem pour $Y$.
\end{propriete}

\begin{exemplebox}[Tableau à double entrée : notes de maths et physique]
On reprend l'exercice Bac (p. \pageref{exo:bac-notes}). Le tableau croise les notes de maths $X$ (lignes) et physique $Y$ (colonnes). La série marginale de $X$ donne le nombre d'élèves par classe de note en maths, celle de $Y$ le nombre par classe en physique.
\end{exemplebox}

\courssub{Nuage de points et point moyen}

\begin{definition}[Nuage de points]
À chaque couple $(x_i\,;\,y_i)$ de la série double, on associe un point $M_i(x_i\,;\,y_i)$ dans un repère orthogonal $(O\,;\,\vec{\imath}\,,\,\vec{\jmath})$. L'ensemble de ces points forme le \cle{nuage de points}.
\end{definition}

\begin{definition}[Point moyen]
Le \cle{point moyen} $G$ du nuage a pour coordonnées $(\bar{x}\,;\,\bar{y})$, où $\bar{x}$ et $\bar{y}$ sont les moyennes des séries marginales.
\end{definition}

\begin{propriete}[Propriété fondamentale du point moyen]
Le point moyen $G$ est le centre de gravité du nuage de points (affecté des effectifs $n_i$). Toute droite d'ajustement affine \emph{raisonnable} (Mayer, moindres carrés) passe par $G$.
\end{propriete}

\begin{methode}[Construire un nuage de points et placer $G$]
\begin{enumerate}
\item Choisir des unités graphiques adaptées à l'étendue des valeurs (ex: \SI{1}{cm} pour 1 unité en $x$, \SI{0,5}{cm} pour 10 unités en $y$).
\item Gradué les axes en commençant si possible à l'origine ou à une valeur proche du minimum.
\item Reporter chaque point $M_i(x_i\,;\,y_i)$.
\item Calculer $\bar{x}$ et $\bar{y}$, placer $G(\bar{x}\,;\,\bar{y})$ (souvent en rouge ou avec un symbole distinctif).
\end{enumerate}
\end{methode}

\courssub{Ajustement affine par la méthode de Mayer}

\begin{definition}[Principe de la méthode de Mayer]
Lorsque le nuage de points suggère une tendance affine, on peut tracer une droite d'ajustement « à l'œil » ou par la \cle{méthode de Mayer} :
\begin{enumerate}
\item Ordonner les points par $x$ croissant.
\item Partager le nuage en \textbf{deux sous-nuages} de même effectif (ou au plus 1 d'écart). Si $n$ est impair, le point central est souvent exclu ou mis dans l'un des deux.
\item Calculer les points moyens $G_1(\bar{x}_1\,;\,\bar{y}_1)$ et $G_2(\bar{x}_2\,;\,\bar{y}_2)$ de chaque sous-nuage.
\item La \cle{droite de Mayer} est la droite $(G_1G_2)$.
\end{enumerate}
\end{definition}

\begin{propriete}[Équation de la droite de Mayer]
Si $\bar{x}_1 \neq \bar{x}_2$, l'équation est $y = ax + b$ avec :
\[
a = \frac{\bar{y}_2 - \bar{y}_1}{\bar{x}_2 - \bar{x}_1},\qquad b = \bar{y}_1 - a\bar{x}_1 \quad(\text{ou } \bar{y}_2 - a\bar{x}_2).
\]
Si $\bar{x}_1 = \bar{x}_2$, la droite est verticale (pas de fonction $y=f(x)$).
\end{propriete}

\begin{attention}[Limites de la méthode de Mayer]
\begin{itemize}
\item Le partage en deux sous-nuages est \textbf{arbitraire} (on pourrait couper autrement).
\item Elle ne minimise aucun critère d'erreur global.
\item Elle est sensible aux valeurs extrêmes (points aberrants) qui tirent $G_1$ ou $G_2$.
\item \textbf{Au Bac :} on l'utilise seulement si l'énoncé le demande explicitement. Sinon, on privilégie les moindres carrés.
\end{itemize}
\end{attention}

\courssub{Covariance et moindres carrés}

\begin{definition}[Covariance]
La \cle{covariance} de $X$ et $Y$ mesure leur liaison linéaire moyenne :
\[
\operatorname{cov}(X,Y) = \frac{1}{n}\sum_{i=1}^{n} n_i (x_i-\bar{x})(y_i-\bar{y}).
\]
\end{definition}

\begin{propriete}[Formule de König-Huygens pour la covariance]
\[
\operatorname{cov}(X,Y) = \overline{xy} - \bar{x}\bar{y} = \frac{1}{n}\sum n_i x_i y_i - \bar{x}\bar{y}.
\]
C'est la formule la plus pratique pour le calcul numérique.
\end{propriete}

\begin{propriete}[Propriétés de la covariance]
\begin{itemize}
\item $\operatorname{cov}(X,X) = V(X)$.
\item $\operatorname{cov}(X,Y) = \operatorname{cov}(Y,X)$ (symétrie).
\item Bilinéarité : $\operatorname{cov}(aX+b,\,cY+d) = ac\,\operatorname{cov}(X,Y)$.
\item Si $X$ et $Y$ sont indépendantes $\Rightarrow \operatorname{cov}(X,Y)=0$. Le réciproque est \textbf{faux} (liaison non linéaire possible).
\end{itemize}
\end{propriete}

\begin{experience}[Découverte de la méthode des moindres carrés]
On cherche une droite $y=ax+b$ qui « colle » au nuage. L'erreur verticale pour $M_i$ est $e_i = y_i - (ax_i+b)$. On minimise la somme des carrés des erreurs $S(a,b) = \sum (y_i - ax_i - b)^2$.
\begin{enumerate}
\item Exprimer $S(a,b)$ en fonction de $a,b,\bar{x},\bar{y},V(X),\operatorname{cov}(X,Y)$.
\item Annuler les dérivées partielles $\frac{\partial S}{\partial a}=0$ et $\frac{\partial S}{\partial b}=0$.
\item Retrouver les formules de $a$ et $b$.
\end{enumerate}
\end{experience}

\begin{propriete}[Droites de régression par les moindres carrés — \textbf{Démonstration exigible au Bac}]
\begin{itemize}
\item \textbf{Droite de régression de $y$ en $x$ (notée $D_{y/x}$)} : prédit $y$ connaissant $x$.
\[
a_{y/x} = \frac{\operatorname{cov}(X,Y)}{V(X)},\qquad b_{y/x} = \bar{y} - a_{y/x}\bar{x}.
\]
\item \textbf{Droite de régression de $x$ en $y$ (notée $D_{x/y}$)} : prédit $x$ connaissant $y$.
\[
a_{x/y} = \frac{\operatorname{cov}(X,Y)}{V(Y)},\qquad b_{x/y} = \bar{x} - a_{x/y}\bar{y}.
\]
\end{itemize}
Les deux droites passent par le point moyen $G(\bar{x}\,;\,\bar{y})$.
\end{propriete}

\begin{aretenir}[Choix de la droite de régression]
\begin{itemize}
\item Pour \textbf{prévoir $y$ à partir de $x$} $\rightarrow$ utiliser $D_{y/x}$ (minimise erreurs verticales).
\item Pour \textbf{prévoir $x$ à partir de $y$} $\rightarrow$ utiliser $D_{x/y}$ (minimise erreurs horizontales).
\item Les deux droites sont \textbf{différentes} (sauf si $|r|=1$). Elles se coupent en $G$.
\end{itemize}
\end{aretenir}

\courssub{Coefficient de corrélation linéaire}

\begin{definition}[Coefficient de corrélation linéaire (de Pearson)]
\[
r = \frac{\operatorname{cov}(X,Y)}{\sqrt{V(X)\,V(Y)}} = \frac{\operatorname{cov}(X,Y)}{\sigma_X\,\sigma_Y}.
\]
\end{definition}

\begin{propriete}[Propriétés fondamentales de $r$]
\begin{enumerate}
\item $-1 \le r \le 1$ (inégalité de Cauchy-Schwarz).
\item $r = 1 \Leftrightarrow$ tous les points sont alignés sur une droite de pente \textbf{positive} ($y=ax+b, a>0$).
\item $r = -1 \Leftrightarrow$ tous les points sont alignés sur une droite de pente \textbf{négative} ($a<0$).
\item $r = 0 \Leftrightarrow \operatorname{cov}(X,Y)=0$ (aucune liaison \textbf{linéaire}). \textbf{Attention :} $r=0$ n'implique pas l'indépendance (ex: $y=x^2$ symétrique).
\item $r$ est \textbf{adimensionnel} (sans unité) et invariant par changement d'origine et d'échelle (affine) sur $X$ et/ou $Y$.
\end{enumerate}
\end{propriete}

\begin{aretenir}[Interprétation de $|r|$ — Qualité de l'ajustement affine]
\begin{center}
\begin{tabular}{|c|c|}
\hline
\rowcolor{popblueL} Valeur de $|r|$ & Interprétation \\ \hline
$|r| < 0{,}5$ & Liaison linéaire faible ou nulle — ajustement \textbf{peu pertinent} \\ \hline
$0{,}5 \le |r| < 0{,}8$ & Liaison linéaire modérée — ajustement \textbf{acceptable} \\ \hline
$0{,}8 \le |r| < 0{,}95$ & Liaison linéaire forte — ajustement \textbf{bon} \\ \hline
$|r| \ge 0{,}95$ & Liaison linéaire très forte — ajustement \textbf{excellent} \\ \hline
\end{tabular}
\end{center}
\textbf{Au Bac :} on considère souvent qu'un ajustement affine est justifié si $|r| \ge 0{,}8$ (seuil usuel). On commente toujours : « $r \approx 0,92$, proche de 1, donc liaison linéaire forte, l'ajustement affine est pertinent. »
\end{aretenir}

\begin{propriete}[Coefficient de détermination $r^2$]
$r^2$ représente la part de la variance de $Y$ « expliquée » par la liaison affine avec $X$.
\[
r^2 = \frac{[\operatorname{cov}(X,Y)]^2}{V(X)V(Y)} = 1 - \frac{\text{Variance résiduelle}}{V(Y)}.
\]
Si $r^2 = 0,85$, 85\% de la dispersion de $Y$ est expliquée par la droite de régression.
\end{propriete}

\courssub{Prévisions et synthèse méthodologique}

\begin{attention}[Interpolation vs Extrapolation]
\begin{itemize}
\item \textbf{Interpolation :} prévoir $y$ pour un $x$ \emph{à l'intérieur} de l'intervalle des $x_i$ observés. Fiable si $|r|$ fort.
\item \textbf{Extrapolation :} prévoir $y$ pour un $x$ \emph{à l'extérieur} de l'intervalle observé. \textbf{Très risqué} : la relation affine peut ne plus tenir (saturation, seuil, courbure). Au Bac, on le fait si demandé, mais on \emph{doit} signaler le risque.
\end{itemize}
\end{attention}

\begin{methode}[Synthèse : Étapes d'une étude de liaison affine]
\begin{enumerate}
\item \textbf{Représenter} le nuage de points et placer $G(\bar{x}\,;\,\bar{y})$.
\item \textbf{Calculer} $\bar{x},\bar{y}, V(X), V(Y), \operatorname{cov}(X,Y)$ (formule de König-Huygens).
\item \textbf{Calculer} $r$ et \textbf{interpréter} sa valeur (signe, intensité, pertinence de l'ajustement).
\item Si $|r|$ est assez grand ($\ge 0,8$) : déterminer l'équation de la droite de régression adaptée ($D_{y/x}$ ou $D_{x/y}$).
\item \textbf{Tracer} la droite de régression (deux points : $G$ et un autre $x$ simple).
\item \textbf{Prévoir} si demandé (interpolation/extrapolation) et \textbf{commenter} la fiabilité.
\end{enumerate}
\end{methode}

\begin{savaistu}[Linéralisation : quand le nuage n'est pas affine]
Si le nuage de points $(x_i\,;\,y_i)$ a une allure \textbf{courbe} (ex: croissance exponentielle, loi de puissance), on peut parfois transformer les variables pour retrouver une liaison affine :
\begin{itemize}
\item Modèle exponentiel $y = A e^{kx}$ $\rightarrow$ poser $z = \ln y$. Le nuage $(x_i\,;\,z_i)$ devient affine : $z = kx + \ln A$.
\item Modèle puissance $y = A x^k$ $\rightarrow$ poser $u = \ln x$, $v = \ln y$. Le nuage $(u_i\,;\,v_i)$ devient affine : $v = k u + \ln A$.
\end{itemize}
On applique alors toute la machinerie (régression, $r$, prévisions) sur les variables transformées, puis on revient au modèle initial. C'est un classique du Bac (voir exercice « Croissance d'une population »).
\end{savaistu}

\courssub{Je vérifie mes connaissances}

\begin{exercice}[QCM]
Pour chaque question, une seule réponse est exacte. Recopie la lettre correspondante sur ta copie.

\begin{enumerate}
\item On considère une série statistique double $(x_i\,;\,y_i)$ de point moyen $G(\bar{x}\,;\,\bar{y})$. La droite de régression de $y$ en $x$ :
\begin{enumerate}
\item passe toujours par l'origine du repère ;
\item passe toujours par le point moyen $G$ ;
\item passe toujours par au moins deux points du nuage.
\end{enumerate}
\item La covariance de deux variables $X$ et $Y$ est définie par :
\begin{enumerate}
\item $\operatorname{cov}(X,Y)=\dfrac{1}{n}\displaystyle\sum_{i=1}^{n}x_i\,y_i$ ;
\item $\operatorname{cov}(X,Y)=\dfrac{1}{n}\displaystyle\sum_{i=1}^{n}(x_i-\bar{x})(y_i-\bar{y})$ ;
\item $\operatorname{cov}(X,Y)=\sqrt{V(X)\,V(Y)}$.
\end{enumerate}
\item Si le coefficient de corrélation linéaire vaut $r=-0{,}98$, alors :
\begin{enumerate}
\item il n'y a aucune relation entre $X$ et $Y$ ;
\item un ajustement affine est pertinent et $Y$ est fonction décroissante de $X$ ;
\item un ajustement affine est pertinent et $Y$ est fonction croissante de $X$.
\end{enumerate}
\item La droite de régression de $y$ en $x$ a pour équation $y=ax+b$ avec :
\begin{enumerate}
\item $a=\dfrac{\operatorname{cov}(X,Y)}{V(X)}$ et $b=\bar{y}-a\bar{x}$ ;
\item $a=\dfrac{V(X)}{\operatorname{cov}(X,Y)}$ et $b=\bar{x}-a\bar{y}$ ;
\item $a=\dfrac{\operatorname{cov}(X,Y)}{V(Y)}$ et $b=\bar{y}-a\bar{x}$.
\end{enumerate}
\end{enumerate}
\end{exercice}

\begin{exercice}[Vrai ou faux, en justifiant]
Réponds par vrai ou faux en justifiant chaque réponse par un argument ou un contre-exemple.

\begin{enumerate}
\item Deux droites de régression (de $y$ en $x$ et de $x$ en $y$) passent toujours par le point moyen du nuage.
\item Si $r=0$, alors les variables $X$ et $Y$ sont indépendantes.
\item Si $|r|=1$, alors tous les points du nuage sont alignés.
\item La méthode de Mayer et la méthode des moindres carrés donnent toujours la même droite d'ajustement.
\item La covariance peut être un nombre négatif.
\end{enumerate}
\end{exercice}

\begin{exercice}[Définitions et vocabulaire]
\begin{enumerate}
\item Définis la \cle{covariance} de deux variables statistiques $X$ et $Y$, puis donne la formule développée de König-Huygens permettant de la calculer.
\item Définis le \cle{coefficient de corrélation linéaire} $r$ de deux variables $X$ et $Y$. Quel encadrement vérifie-t-il toujours ?
\item Qu'appelle-t-on \cle{séries marginales} d'un tableau statistique à double entrée ?
\item À quelle condition sur $r$ juge-t-on, en pratique, qu'un ajustement affine est de bonne qualité ?
\end{enumerate}
\end{exercice}

\begin{exercice}[Calculs directs]
On considère la série statistique double suivante :

\begin{center}
\begin{tabular}{|c|cccc|}
\hline
\rowcolor{popblueL} $x_i$ & 1 & 2 & 3 & 4 \\
\hline
$y_i$ & 3 & 5 & 6 & 10 \\
\hline
\end{tabular}
\end{center}

\begin{enumerate}
\item Calcule les moyennes $\bar{x}$ et $\bar{y}$, puis place le point moyen $G$.
\item Calcule la variance $V(X)$ et la covariance $\operatorname{cov}(X,Y)$.
\item Détermine l'équation de la droite de régression de $y$ en $x$.
\item Calcule le coefficient de corrélation linéaire $r$ (on donne $V(Y)=7{,}25$). Un ajustement affine est-il justifié ?
\end{enumerate}
\end{exercice}

\courssub{Je m'entraîne}

\begin{exercice}[Nuage, point moyen et méthode de Mayer]
Une coopérative agricole de l'Ouest Cameroun relève la quantité d'engrais utilisée $x_i$ (en sacs) et le rendement du maïs $y_i$ (en quintaux par hectare) sur six parcelles :

\begin{center}
\begin{tabular}{|c|cccccc|}
\hline
\rowcolor{popblueL} $x_i$ (sacs) & 1 & 2 & 3 & 4 & 5 & 6 \\
\hline
$y_i$ (q/ha) & 8 & 11 & 13 & 15 & 18 & 21 \\
\hline
\end{tabular}
\end{center}

\begin{enumerate}
\item Représente le nuage de points dans un repère orthogonal (unités : \SI{2}{cm} pour 1 sac en abscisse, \SI{0,5}{cm} pour 1 quintal en ordonnée).
\item Calcule les coordonnées du point moyen $G$ et place-le sur le graphique.
\item Partage le nuage en deux sous-nuages de trois points chacun ; détermine les points moyens $G_1$ et $G_2$.
\item Détermine l'équation de la droite de Mayer $(G_1G_2)$ et trace-la.
\item À l'aide de cette droite, estime le rendement attendu pour 8 sacs d'engrais.
\end{enumerate}
\end{exercice}

\begin{exercice}[Droites de régression par les moindres carrés]
Le tableau ci-dessous donne la dépense mensuelle en données internet $x_i$ (en milliers de FCFA) et le nombre d'heures de connexion $y_i$ de huit abonnés d'un quartier de Douala :

\begin{center}
\begin{tabular}{|c|cccccccc|}
\hline
\rowcolor{popblueL} $x_i$ & 2 & 3 & 5 & 5 & 7 & 8 & 10 & 12 \\
\hline
$y_i$ & 10 & 14 & 18 & 22 & 26 & 30 & 36 & 40 \\
\hline
\end{tabular}
\end{center}

On donne : $\displaystyle\sum x_i=52$ ; $\displaystyle\sum y_i=196$ ; $\displaystyle\sum x_i^2=424$ ; $\displaystyle\sum y_i^2=5616$ ; $\displaystyle\sum x_iy_i=1554$.

\begin{enumerate}
\item Calcule $\bar{x}$, $\bar{y}$, $V(X)$ et $\operatorname{cov}(X,Y)$.
\item Détermine l'équation de la droite de régression $D_{y/x}$ de $y$ en $x$.
\item Détermine l'équation de la droite de régression $D_{x/y}$ de $x$ en $y$ (on donne $V(Y)=102$).
\item Vérifie que les deux droites passent par le point moyen $G$.
\end{enumerate}
\end{exercice}

\begin{exercice}[Comparaison de deux ajustements]
On considère la série double suivante :

\begin{center}
\begin{tabular}{|c|cccccc|}
\hline
\rowcolor{popblueL} $x_i$ & 1 & 2 & 3 & 4 & 5 & 6 \\
\hline
$y_i$ & 2{,}1 & 3{,}9 & 6{,}2 & 7{,}8 & 10{,}1 & 11{,}9 \\
\hline
\end{tabular}
\end{center}

\begin{enumerate}
\item Détermine, par la méthode de Mayer (deux sous-nuages de trois points), une équation de la droite d'ajustement $D_1$.
\item Détermine, par la méthode des moindres carrés, l'équation de la droite de régression $D_2$ de $y$ en $x$.
\item Pour $x=10$, calcule la prévision donnée par chacune des deux droites, puis l'écart absolu entre ces deux prévisions.
\item Calcule le coefficient de corrélation linéaire $r$. Laquelle des deux méthodes te semble la plus fiable pour prévoir ? Justifie.
\end{enumerate}
\end{exercice}

\begin{exercice}[Exercice inverse]
Pour une série statistique double $(X\,;\,Y)$ de 50 observations, on sait que :
\begin{itemize}
\item le point moyen est $G(12\,;\,\bar{y})$ avec $\bar{y}$ inconnu ;
\item $V(X)=16$ et $\operatorname{cov}(X,Y)=24$ ;
\item la droite de régression de $y$ en $x$ a pour équation $y=1{,}5x+4$.
\end{itemize}

\begin{enumerate}
\item Retrouve la valeur de $\bar{y}$ en utilisant le fait que la droite de régression passe par $G$.
\item Vérifie la cohérence du coefficient directeur $1{,}5$ avec les données de l'énoncé.
\item On donne de plus $V(Y)=40$. Calcule le coefficient de corrélation linéaire $r$.
\item Détermine l'équation de la droite de régression de $x$ en $y$ sous la forme $x=a'y+b'$.
\end{enumerate}
\end{exercice}

\courssub{Je me prépare au Bac}

\begin{exercice}[Type Bac — Chiffre d'affaires d'une entreprise de téléphonie]\label{exo:bac-telecom}
Une entreprise de vente de téléphones et de services mobiles, installée à Yaoundé, a relevé son chiffre d'affaires annuel $y_i$ (en millions de FCFA) en fonction du rang $x_i$ de l'année :

\begin{center}
\begin{tabular}{|c|cccccccc|}
\hline
\rowcolor{popblueL} Rang $x_i$ & 1 & 2 & 3 & 4 & 5 & 6 & 7 & 8 \\
\hline
CA $y_i$ (millions FCFA) & 42 & 50 & 57 & 65 & 74 & 81 & 90 & 97 \\
\hline
\end{tabular}
\end{center}

On donne les sommes \textbf{corrigées} (calculées à partir du tableau) :
$\displaystyle\sum_{i=1}^{8}x_i^2=204$ ; $\displaystyle\sum_{i=1}^{8}x_iy_i=2835$ ; $\displaystyle\sum_{i=1}^{8}y_i^2=41284$.

\begin{enumerate}
\item Représente le nuage de points $M_i(x_i\,;\,y_i)$ dans un repère orthogonal (unités : \SI{1}{cm} pour 1 an en abscisse ; \SI{1}{cm} pour 10 millions de FCFA en ordonnée, en commençant la graduation à 40).
\item Calcule les coordonnées du point moyen $G$ et place-le dans le repère.
\item \begin{enumerate}
\item Calcule la variance $V(X)$ et la covariance $\operatorname{cov}(X,Y)$.
\item Détermine une équation de la droite de régression $D$ de $y$ en $x$ (arrondir les coefficients à $10^{-2}$ près).
\item Trace $D$ dans le repère.
\end{enumerate}
\item \begin{enumerate}
\item Calcule le coefficient de corrélation linéaire $r$ (arrondir à $10^{-3}$ près).
\item Un ajustement affine de $y$ en $x$ est-il justifié ? Justifie ta réponse.
\end{enumerate}
\item En admettant que l'évolution se poursuive ainsi, estime le chiffre d'affaires de l'entreprise pour l'année de rang 9, puis pour l'année de rang 12.
\item À partir de quel rang d'année le chiffre d'affaires dépassera-t-il, selon ce modèle, 150 millions de FCFA ?
\end{enumerate}
\end{exercice}

\begin{exercice}[Type Bac — Tableau à double entrée : notes de maths et de physique]\label{exo:bac-notes}
Dans une classe de Terminale C de 40 élèves, on a relevé les notes de mathématiques $X$ et de physique $Y$ (sur 20) au premier devoir harmonisé. Les résultats sont regroupés dans le tableau à double entrée suivant :

\begin{center}
\begin{tabular}{|c|c|c|c|c|}
\hline
\rowcolor{popblueL} $X$ \textbackslash{} $Y$ & $[0\,;\,5[$ & $[5\,;\,10[$ & $[10\,;\,15[$ & $[15\,;\,20]$ \\
\hline
$[0\,;\,5[$ & 2 & 1 & 0 & 0 \\
\hline
$[5\,;\,10[$ & 1 & 8 & 4 & 0 \\
\hline
$[10\,;\,15[$ & 0 & 3 & 10 & 3 \\
\hline
$[15\,;\,20]$ & 0 & 0 & 2 & 6 \\
\hline
\end{tabular}
\end{center}

On prendra les centres des classes comme valeurs des caractères : 2,5 ; 7,5 ; 12,5 ; 17,5.

\begin{enumerate}
\item Détermine la série marginale de la variable $X$ (effectifs par classe), puis celle de la variable $Y$.
\item Calcule les moyennes $\bar{x}$ et $\bar{y}$ des deux séries marginales.
\item \begin{enumerate}
\item Calcule les variances $V(X)$ et $V(Y)$.
\item Calcule la covariance $\operatorname{cov}(X,Y)$ (on pourra utiliser $\operatorname{cov}(X,Y)=\overline{xy}-\bar{x}\,\bar{y}$ en dressant le tableau des produits $x_iy_j$ pondérés).
\end{enumerate}
\item Calcule le coefficient de corrélation linéaire $r$ entre les notes de mathématiques et de physique (arrondir à $10^{-2}$ près).
\item Le professeur de physique affirme : « Les élèves forts en mathématiques sont en général forts en physique. » Ce résultat statistique confirme-t-il son affirmation ? Rédige un commentaire de trois à quatre lignes.
\item Détermine l'équation de la droite de régression de $y$ en $x$, puis estime la note de physique d'un élève ayant 14 en mathématiques.
\end{enumerate}
\end{exercice}

\begin{exercice}[Type Bac — Changement de variable : croissance d'une population]
La population d'une ville du Nord-Cameroun, bien desservie par un nouveau programme d'électrification rurale, a été relevée tous les ans pendant six ans. On note $x_i$ le rang de l'année et $y_i$ la population en milliers d'habitants :

\begin{center}
\begin{tabular}{|c|cccccc|}
\hline
\rowcolor{popblueL} Rang $x_i$ & 1 & 2 & 3 & 4 & 5 & 6 \\
\hline
Population $y_i$ (milliers) & 20{,}0 & 24{,}2 & 29{,}3 & 35{,}5 & 43{,}0 & 52{,}1 \\
\hline
\end{tabular}
\end{center}

Le nuage de points $M_i(x_i\,;\,y_i)$ présente une allure nettement courbe : un ajustement affine de $y$ en $x$ ne semble pas approprié. On pose alors $z_i=\ln(y_i)$.

\begin{enumerate}
\item Recopie et complète le tableau suivant (arrondir les $z_i$ à $10^{-3}$ près) :

\begin{center}
\begin{tabular}{|c|cccccc|}
\hline
\rowcolor{popblueL} $x_i$ & 1 & 2 & 3 & 4 & 5 & 6 \\
\hline
$z_i=\ln(y_i)$ & 2,996 & 3,186 & 3,378 & 3,570 & 3,761 & 3,953 \\
\hline
\end{tabular}
\end{center}

\item Représente le nuage de points $N_i(x_i\,;\,z_i)$. Que constates-tu quant à son allure ?
\item On admet que $z$ et $x$ sont liés par une relation affine $z=ax+b$.
\begin{enumerate}
\item Calcule le coefficient de corrélation linéaire $r'$ entre $x$ et $z$ (on donne : $\bar{z}\approx 3{,}474$ ; $\displaystyle\sum x_iz_i\approx 76{,}304$ ; $\displaystyle\sum z_i^2\approx 73{,}05$ ; arrondir $r'$ à $10^{-3}$ près). Conclus sur la pertinence de l'ajustement affine de $z$ en $x$.
\item Détermine l'équation de la droite de régression de $z$ en $x$ (arrondir $a$ à $10^{-3}$ et $b$ à $10^{-2}$ près).
\end{enumerate}
\item Déduis-en une expression de $y$ en fonction de $x$ de la forme $y=A\,\mathrm{e}^{ax}$, où $A$ est un réel à déterminer (arrondir à l'unité).
\item Estime la population de la ville à l'année de rang 11 (c'est-à-dire dans 5 ans). Arrondir au millier d'habitants.
\item Selon ce modèle, à partir de quelle année la population dépassera-t-elle 100 000 habitants ? (On pourra résoudre une inéquation avec la fonction $\ln$.)
\end{enumerate}
\end{exercice}



\newpage

\coursec{Corrigés des exercices}

\begin{corrige}[Exercice 1 — QCM]
\begin{enumerate}
\item \textbf{Réponse b.} La droite de régression de $y$ en $x$ a pour équation $y=ax+b$ avec $b=\bar{y}-a\bar{x}$ ; en $x=\bar{x}$, on obtient $y=a\bar{x}+(\bar{y}-a\bar{x})=\bar{y}$ : elle passe donc toujours par $G(\bar{x}\,;\,\bar{y})$. Elle ne passe ni par l'origine en général, ni nécessairement par des points du nuage.
\item \textbf{Réponse b.} Par définition, $\operatorname{cov}(X,Y)=\dfrac{1}{n}\displaystyle\sum_{i=1}^{n}(x_i-\bar{x})(y_i-\bar{y})$. La réponse a oublie le centrage (c'est $\overline{xy}$, pas la covariance) ; la réponse c est le produit des écarts types.
\item \textbf{Réponse b.} $r=-0{,}98$ : $|r|$ est très proche de 1, donc la liaison linéaire est très forte et l'ajustement affine est pertinent ; le signe négatif de $r$ (et donc de $\operatorname{cov}(X,Y)$, donc de la pente $a$) traduit une fonction \textbf{décroissante}.
\item \textbf{Réponse a.} $a=\dfrac{\operatorname{cov}(X,Y)}{V(X)}$ et $b=\bar{y}-a\bar{x}$. Attention : la réponse c est le coefficient de la droite de régression de $x$ en $y$.
\end{enumerate}
\end{corrige}

\begin{corrige}[Exercice 2 — Vrai ou faux, en justifiant]
\begin{enumerate}
\item \textbf{Vrai.} Pour $D_{y/x}$ : $b=\bar{y}-a\bar{x}$, donc $y(\bar{x})=a\bar{x}+b=\bar{y}$. Pour $D_{x/y}$ : $b'=\bar{x}-a'\bar{y}$, donc $x(\bar{y})=a'\bar{y}+b'=\bar{x}$. Les deux droites passent par $G(\bar{x}\,;\,\bar{y})$.
\item \textbf{Faux.} $r=0$ signifie seulement l'absence de liaison \emph{linéaire}. Contre-exemple : les points $(-2\,;\,4)$, $(-1\,;\,1)$, $(0\,;\,0)$, $(1\,;\,1)$, $(2\,;\,4)$ vérifient $y=x^2$ (liaison parfaite, donc dépendance totale) et pourtant $\operatorname{cov}(X,Y)=0$ par symétrie, donc $r=0$.
\item \textbf{Vrai.} $|r|=1$ si et seulement si tous les points sont alignés sur une droite non verticale (pente positive si $r=1$, négative si $r=-1$). C'est une conséquence du cas d'égalité dans l'inégalité de Cauchy-Schwarz.
\item \textbf{Faux.} La méthode de Mayer dépend du partage (arbitraire) en deux sous-nuages, alors que les moindres carrés minimisent un critère global unique. Voir l'exercice « Comparaison de deux ajustements » : on obtient $y\approx1{,}956x+0{,}155$ (Mayer) et $y\approx1{,}977x+0{,}081$ (moindres carrés) : deux droites distinctes.
\item \textbf{Vrai.} La covariance donne le sens de la liaison : elle est négative dès que $Y$ tend à décroître quand $X$ croît. Exemple : $(1\,;\,10)$, $(2\,;\,6)$, $(3\,;\,2)$ : $\bar{x}=2$, $\bar{y}=6$,
\[
\operatorname{cov}(X,Y)=\frac{1\times10+2\times6+3\times2}{3}-2\times6=\frac{28}{3}-12=-\frac{8}{3}<0.
\]
\end{enumerate}
\end{corrige}

\begin{corrige}[Exercice 3 — Définitions et vocabulaire]
\begin{enumerate}
\item La \cle{covariance} de $X$ et $Y$ est la moyenne des produits des écarts à la moyenne :
\[
\operatorname{cov}(X,Y)=\frac{1}{n}\sum_{i=1}^{p}n_i(x_i-\bar{x})(y_i-\bar{y}).
\]
Formule de König-Huygens (forme développée, pratique pour les calculs) :
\[
\operatorname{cov}(X,Y)=\overline{xy}-\bar{x}\,\bar{y}=\frac{1}{n}\sum_{i=1}^{p}n_i x_i y_i-\bar{x}\,\bar{y}.
\]
\item Le \cle{coefficient de corrélation linéaire} est
\[
r=\frac{\operatorname{cov}(X,Y)}{\sqrt{V(X)\,V(Y)}}=\frac{\operatorname{cov}(X,Y)}{\sigma_X\,\sigma_Y}.
\]
Il vérifie toujours l'encadrement $-1\le r\le 1$ (conséquence de l'inégalité de Cauchy-Schwarz).
\item Les \cle{séries marginales} d'un tableau à double entrée sont les deux séries à une variable obtenues en sommant les effectifs par ligne (série marginale de $X$, effectifs $n_{i\cdot}$) et par colonne (série marginale de $Y$, effectifs $n_{\cdot j}$). Elles décrivent la distribution de chaque caractère pris isolément.
\item En pratique, on juge l'ajustement affine de bonne qualité lorsque $|r|\ge 0{,}8$ (liaison linéaire forte), et excellent lorsque $|r|\ge 0{,}95$. En dessous de $0{,}5$, l'ajustement affine n'a guère de sens.
\end{enumerate}
\end{corrige}

\begin{corrige}[Exercice 4 — Calculs directs]
\begin{enumerate}
\item $n=4$. $\bar{x}=\dfrac{1+2+3+4}{4}=\dfrac{10}{4}=2{,}5$ ; $\bar{y}=\dfrac{3+5+6+10}{4}=\dfrac{24}{4}=6$. Donc $G(2{,}5\,;\,6)$.
\item Variance de $X$ (König-Huygens) :
\[
V(X)=\frac{1^2+2^2+3^2+4^2}{4}-(2{,}5)^2=\frac{30}{4}-6{,}25=7{,}5-6{,}25=1{,}25.
\]
Moyenne des produits : $\overline{xy}=\dfrac{1\times3+2\times5+3\times6+4\times10}{4}=\dfrac{3+10+18+40}{4}=\dfrac{71}{4}=17{,}75$.
\[
\operatorname{cov}(X,Y)=\overline{xy}-\bar{x}\bar{y}=17{,}75-2{,}5\times6=17{,}75-15=2{,}75.
\]
\item $a=\dfrac{\operatorname{cov}(X,Y)}{V(X)}=\dfrac{2{,}75}{1{,}25}=2{,}2$ ; $b=\bar{y}-a\bar{x}=6-2{,}2\times2{,}5=6-5{,}5=0{,}5$.
\[
D_{y/x}\ :\ y=2{,}2x+0{,}5.
\]
\item Variance de $Y$ (König-Huygens) :
\[
V(Y)=\frac{3^2+5^2+6^2+10^2}{4}-6^2=\frac{9+25+36+100}{4}-36=\frac{170}{4}-36=42{,}5-36=6{,}5.
\]
$\sigma_X=\sqrt{1{,}25}\approx1{,}118$ ; $\sigma_Y=\sqrt{6{,}5}\approx2{,}550$.
\[
r=\frac{2{,}75}{1{,}118\times2{,}550}\approx\frac{2{,}75}{2{,}851}\approx0{,}965.
\]
Comme $|r|\approx0{,}97>0{,}95$, la liaison linéaire est très forte : l'ajustement affine est \textbf{pleinement justifié}.
\end{enumerate}
\end{corrige}

\begin{corrige}[Exercice 5 — Nuage, point moyen et méthode de Mayer]
\begin{enumerate}
\item On place les points $M_1(1\,;\,8)$, $M_2(2\,;\,11)$, $M_3(3\,;\,13)$, $M_4(4\,;\,15)$, $M_5(5\,;\,18)$, $M_6(6\,;\,21)$ dans le repère avec les unités indiquées. Le nuage a une allure nettement rectiligne ascendante.
\item $\bar{x}=\dfrac{1+2+3+4+5+6}{6}=\dfrac{21}{6}=3{,}5$ ; $\bar{y}=\dfrac{8+11+13+15+18+21}{6}=\dfrac{86}{6}=\dfrac{43}{3}\approx14{,}33$.
Donc $G\left(3{,}5\,;\,\dfrac{43}{3}\right)\approx(3{,}5\,;\,14{,}33)$, placé sur le graphique.
\item Sous-nuage 1 (points 1 à 3) : $\bar{x}_1=\dfrac{1+2+3}{3}=2$, $\bar{y}_1=\dfrac{8+11+13}{3}=\dfrac{32}{3}\approx10{,}67$, donc $G_1\left(2\,;\,\dfrac{32}{3}\right)$.
Sous-nuage 2 (points 4 à 6) : $\bar{x}_2=\dfrac{4+5+6}{3}=5$, $\bar{y}_2=\dfrac{15+18+21}{3}=18$, donc $G_2(5\,;\,18)$.
\item Pente : $a=\dfrac{18-\frac{32}{3}}{5-2}=\dfrac{\frac{22}{3}}{3}=\dfrac{22}{9}\approx2{,}44$.
Ordonnée à l'origine : $b=\dfrac{32}{3}-\dfrac{22}{9}\times2=\dfrac{96-44}{9}=\dfrac{52}{9}\approx5{,}78$.
\[
(G_1G_2)\ :\ y=\frac{22}{9}x+\frac{52}{9}\quad\text{(soit environ } y=2{,}44x+5{,}78\text{)}.
\]
On la trace en reliant $G_1$ et $G_2$ ; on vérifie qu'elle passe par $G$ : $\frac{22}{9}\times3{,}5+\frac{52}{9}=\frac{77+52}{9}=\frac{129}{9}=\frac{43}{3}=\bar{y}$.
\item Pour $x=8$ : $y=\dfrac{22}{9}\times8+\dfrac{52}{9}=\dfrac{176+52}{9}=\dfrac{228}{9}\approx25{,}3$ q/ha.
\textbf{Commentaire :} il s'agit d'une \emph{extrapolation} ($x=8$ est hors de l'intervalle $[1\,;\,6]$ observé) : l'estimation de 25,3 q/ha est à prendre avec prudence (saturation possible du rendement).
\end{enumerate}
\end{corrige}

\begin{corrige}[Exercice 6 — Droites de régression par les moindres carrés]
\begin{enumerate}
\item $n=8$. $\bar{x}=\dfrac{52}{8}=6{,}5$ ; $\bar{y}=\dfrac{196}{8}=24{,}5$.
\[
V(X)=\frac{424}{8}-(6{,}5)^2=53-42{,}25=10{,}75.
\]
$\overline{xy}=\dfrac{1554}{8}=194{,}25$, donc
\[
\operatorname{cov}(X,Y)=194{,}25-6{,}5\times24{,}5=194{,}25-159{,}25=35.
\]
\item $D_{y/x}$ : $a=\dfrac{\operatorname{cov}(X,Y)}{V(X)}=\dfrac{35}{10{,}75}=\dfrac{140}{43}\approx3{,}256$.
\[
b=\bar{y}-a\bar{x}=24{,}5-\frac{140}{43}\times6{,}5=\frac{1053{,}5-910}{43}=\frac{143{,}5}{43}\approx3{,}337.
\]
\[
D_{y/x}\ :\ y=\frac{140}{43}x+\frac{143{,}5}{43}\quad\text{(soit } y\approx3{,}26x+3{,}34\text{)}.
\]
\item $D_{x/y}$ : $a'=\dfrac{\operatorname{cov}(X,Y)}{V(Y)}=\dfrac{35}{102}\approx0{,}343$.
\[
b'=\bar{x}-a'\bar{y}=6{,}5-\frac{35}{102}\times24{,}5=\frac{663-857{,}5}{102}=-\frac{194{,}5}{102}\approx-1{,}907.
\]
\[
D_{x/y}\ :\ x=\frac{35}{102}y-\frac{194{,}5}{102}\quad\text{(soit } x\approx0{,}343y-1{,}91\text{)}.
\]
\item Pour $D_{y/x}$ : $a\bar{x}+b=\dfrac{140}{43}\times6{,}5+\dfrac{143{,}5}{43}=\dfrac{910+143{,}5}{43}=\dfrac{1053{,}5}{43}=24{,}5=\bar{y}$.
Pour $D_{x/y}$ : $a'\bar{y}+b'=\dfrac{35}{102}\times24{,}5-\dfrac{194{,}5}{102}=\dfrac{857{,}5-194{,}5}{102}=\dfrac{663}{102}=6{,}5=\bar{x}$.
Les deux droites passent bien par $G(6{,}5\,;\,24{,}5)$.
\end{enumerate}
\end{corrige}

\begin{corrige}[Exercice 7 — Comparaison de deux ajustements]
\begin{enumerate}
\item \textbf{Mayer.} Sous-nuage 1 : $\bar{x}_1=2$, $\bar{y}_1=\dfrac{2{,}1+3{,}9+6{,}2}{3}=\dfrac{12{,}2}{3}\approx4{,}067$, donc $G_1(2\,;\,4{,}067)$.
Sous-nuage 2 : $\bar{x}_2=5$, $\bar{y}_2=\dfrac{7{,}8+10{,}1+11{,}9}{3}=\dfrac{29{,}8}{3}\approx9{,}933$, donc $G_2(5\,;\,9{,}933)$.
\[
a_1=\frac{9{,}933-4{,}067}{5-2}=\frac{5{,}866}{3}\approx1{,}956\,;\qquad b_1=4{,}067-1{,}956\times2\approx0{,}155.
\]
\[
D_1\ :\ y\approx1{,}956x+0{,}155.
\]
\item \textbf{Moindres carrés.} $\bar{x}=\dfrac{21}{6}=3{,}5$ ; $\bar{y}=\dfrac{42}{6}=7$.
$V(X)=\dfrac{1+4+9+16+25+36}{6}-12{,}25=\dfrac{91}{6}-12{,}25=\dfrac{35}{12}\approx2{,}917$.
$\overline{xy}=\dfrac{2{,}1+7{,}8+18{,}6+31{,}2+50{,}5+71{,}4}{6}=\dfrac{181{,}6}{6}\approx30{,}267$.
$\operatorname{cov}(X,Y)=30{,}267-3{,}5\times7=30{,}267-24{,}5=5{,}767$.
\[
a_2=\frac{5{,}767}{2{,}917}\approx1{,}977\,;\qquad b_2=7-1{,}977\times3{,}5\approx0{,}081.
\]
\[
D_2\ :\ y\approx1{,}977x+0{,}081.
\]
\item Pour $x=10$ : $D_1$ donne $1{,}956\times10+0{,}155\approx19{,}71$ ; $D_2$ donne $1{,}977\times10+0{,}081\approx19{,}85$.
Écart absolu : $|19{,}85-19{,}71|\approx0{,}14$.
\item $V(Y)=\dfrac{2{,}1^2+3{,}9^2+6{,}2^2+7{,}8^2+10{,}1^2+11{,}9^2}{6}-49=\dfrac{362{,}52}{6}-49=60{,}42-49=11{,}42$.
\[
r=\frac{5{,}767}{\sqrt{2{,}917\times11{,}42}}=\frac{5{,}767}{\sqrt{33{,}31}}\approx\frac{5{,}767}{5{,}772}\approx0{,}999.
\]
$r$ est extrêmement proche de 1 : la liaison linéaire est quasi parfaite. La méthode des \textbf{moindres carrés} est la plus fiable : elle minimise globalement la somme des carrés des erreurs et ne dépend d'aucun choix arbitraire, contrairement au partage de Mayer.
\end{enumerate}
\end{corrige}

\begin{corrige}[Exercice 8 — Exercice inverse]
\begin{enumerate}
\item La droite de régression de $y$ en $x$ passe par $G(12\,;\,\bar{y})$, donc :
\[
\bar{y}=1{,}5\times12+4=18+4=22.
\]
\item Le coefficient directeur théorique est $a=\dfrac{\operatorname{cov}(X,Y)}{V(X)}=\dfrac{24}{16}=1{,}5$. Il coïncide avec celui de l'énoncé : les données sont \textbf{cohérentes}.
\item \[
r=\frac{\operatorname{cov}(X,Y)}{\sqrt{V(X)V(Y)}}=\frac{24}{\sqrt{16\times40}}=\frac{24}{\sqrt{640}}=\frac{24}{8\sqrt{10}}=\frac{3}{\sqrt{10}}\approx0{,}949.
\]
Liaison linéaire forte ($|r|\approx0{,}95$).
\item $D_{x/y}$ : $a'=\dfrac{\operatorname{cov}(X,Y)}{V(Y)}=\dfrac{24}{40}=0{,}6$ ; $b'=\bar{x}-a'\bar{y}=12-0{,}6\times22=12-13{,}2=-1{,}2$.
\[
D_{x/y}\ :\ x=0{,}6y-1{,}2.
\]
\textbf{Vérification :} le produit des pentes $a\times a'=1{,}5\times0{,}6=0{,}9=r^2$ (car $r^2=\frac{9}{10}$) : cohérent.
\end{enumerate}
\end{corrige}

\begin{corrige}[Exercice 9 — Type Bac : chiffre d'affaires d'une entreprise de téléphonie]
\begin{enumerate}
\item On place les points $M_i(x_i\,;\,y_i)$ avec les unités indiquées (graduation en ordonnée à partir de 40). Le nuage est quasiment rectiligne, croissant.
\item $n=8$ ; $\displaystyle\sum x_i=36$, $\displaystyle\sum y_i=556$.
\[
\bar{x}=\frac{36}{8}=4{,}5\,;\qquad \bar{y}=\frac{556}{8}=69{,}5.\qquad G(4{,}5\,;\,69{,}5).
\]
\item \begin{enumerate}
\item $V(X)=\dfrac{204}{8}-(4{,}5)^2=25{,}5-20{,}25=5{,}25$.
$\overline{xy}=\dfrac{2835}{8}=354{,}375$, donc
\[
\operatorname{cov}(X,Y)=354{,}375-4{,}5\times69{,}5=354{,}375-312{,}75=41{,}625.
\]
\item $a=\dfrac{41{,}625}{5{,}25}\approx7{,}93$ ; $b=69{,}5-7{,}93\times4{,}5\approx69{,}5-35{,}68\approx33{,}82$.
\[
D\ :\ y=7{,}93x+33{,}82.
\]
\item On trace $D$ par deux points : $G(4{,}5\,;\,69{,}5)$ et par exemple $A(0\,;\,33{,}82)$ ou $B(8\,;\,97{,}26)$.
\end{enumerate}
\item \begin{enumerate}
\item $V(Y)=\dfrac{41284}{8}-(69{,}5)^2=5160{,}5-4830{,}25=330{,}25$.
\[
r=\frac{41{,}625}{\sqrt{5{,}25\times330{,}25}}=\frac{41{,}625}{\sqrt{1733{,}8125}}\approx\frac{41{,}625}{41{,}639}\approx0{,}9997.
\]
\item $r\approx0{,}9997$ : la liaison linéaire est quasi parfaite ($|r|\ge0{,}95$). L'ajustement affine de $y$ en $x$ est \textbf{pleinement justifié}.
\end{enumerate}
\item Rang 9 : $y=7{,}93\times9+33{,}82\approx105{,}2$ millions de FCFA.
Rang 12 : $y=7{,}93\times12+33{,}82\approx129{,}0$ millions de FCFA.
Ce sont des \emph{extrapolations} : elles supposent que la tendance se maintienne, ce qui doit être signalé.
\item $7{,}93x+33{,}82>150\iff7{,}93x>116{,}18\iff x>14{,}65$.
Le chiffre d'affaires dépassera 150 millions de FCFA à partir du \textbf{rang 15}.
\end{enumerate}
\textbf{Barème indicatif (sur 10) :} nuage 1 pt ; $G$ 1 pt ; $V(X)$ et covariance 2 pts ; équation de $D$ 1,5 pt ; tracé 0,5 pt ; $r$ et justification 1,5 pt ; prévisions 1,5 pt ; inéquation 1 pt.

\textbf{Points de vigilance :} citer la formule de König-Huygens avant tout calcul de variance ou covariance ; justifier la pertinence de l'ajustement \emph{uniquement} à partir de $r$ (pas « à l'œil ») ; mentionner le caractère extrapolé des prévisions ; bien arrondir comme demandé ($10^{-2}$ pour $a,b$ ; $10^{-3}$ pour $r$).
\end{corrige}

\begin{corrige}[Exercice 10 — Type Bac : tableau à double entrée, notes de maths et de physique]
\begin{enumerate}
\item En sommant par lignes, série marginale de $X$ : $3\,;\,13\,;\,16\,;\,8$ (total 40).
En sommant par colonnes, série marginale de $Y$ : $3\,;\,12\,;\,16\,;\,9$ (total 40).
\item Avec les centres $2{,}5\,;\,7{,}5\,;\,12{,}5\,;\,17{,}5$ :
\[
\bar{x}=\frac{3\times2{,}5+13\times7{,}5+16\times12{,}5+8\times17{,}5}{40}=\frac{7{,}5+97{,}5+200+140}{40}=\frac{445}{40}=11{,}125.
\]
\[
\bar{y}=\frac{3\times2{,}5+12\times7{,}5+16\times12{,}5+9\times17{,}5}{40}=\frac{7{,}5+90+200+157{,}5}{40}=\frac{455}{40}=11{,}375.
\]
\item \begin{enumerate}
\item $\overline{x^2}=\dfrac{3\times6{,}25+13\times56{,}25+16\times156{,}25+8\times306{,}25}{40}=\dfrac{5700}{40}=142{,}5$.
\[
V(X)=142{,}5-(11{,}125)^2=142{,}5-123{,}766\approx18{,}73.
\]
$\overline{y^2}=\dfrac{3\times6{,}25+12\times56{,}25+16\times156{,}25+9\times306{,}25}{40}=\dfrac{5950}{40}=148{,}75$.
\[
V(Y)=148{,}75-(11{,}375)^2=148{,}75-129{,}391\approx19{,}36.
\]
\item Tableau des produits pondérés $n_{ij}x_iy_j$ (cases non nulles) :
$2\times6{,}25=12{,}5$ ; $1\times18{,}75=18{,}75$ ; $1\times18{,}75=18{,}75$ ; $8\times56{,}25=450$ ; $4\times93{,}75=375$ ; $3\times93{,}75=281{,}25$ ; $10\times156{,}25=1562{,}5$ ; $3\times218{,}75=656{,}25$ ; $2\times218{,}75=437{,}5$ ; $6\times306{,}25=1837{,}5$.
Somme : $5650$, donc $\overline{xy}=\dfrac{5650}{40}=141{,}25$.
\[
\operatorname{cov}(X,Y)=141{,}25-11{,}125\times11{,}375=141{,}25-126{,}547\approx14{,}70.
\]
\end{enumerate}
\item \[
r=\frac{14{,}70}{\sqrt{18{,}73\times19{,}36}}=\frac{14{,}70}{\sqrt{362{,}7}}\approx\frac{14{,}70}{19{,}04}\approx0{,}77.
\]
\item $r\approx0{,}77$ est \textbf{positif} et proche de $0{,}8$ : il existe une liaison linéaire modérée à forte entre les deux notes, qui varient globalement dans le même sens. L'affirmation du professeur est donc \textbf{plutôt confirmée} : un élève fort en maths a de bonnes chances d'être fort en physique. Toutefois, comme $|r|<0{,}8$, la dispersion reste notable : la note de maths ne suffit pas à prédire précisément celle de physique.
\item $a=\dfrac{14{,}70}{18{,}73}\approx0{,}785$ ; $b=11{,}375-0{,}785\times11{,}125\approx2{,}64$.
\[
D_{y/x}\ :\ y=0{,}785x+2{,}64.
\]
Pour $x=14$ : $y\approx0{,}785\times14+2{,}64\approx13{,}6$. Note de physique estimée : \textbf{environ 13,6/20}.
\end{enumerate}
\textbf{Barème indicatif (sur 10) :} séries marginales 1,5 pt ; moyennes 1,5 pt ; variances 1,5 pt ; covariance 2 pts ; $r$ 1 pt ; commentaire 1 pt ; droite et estimation 1,5 pt.

\textbf{Points de vigilance :} utiliser les \emph{centres de classes} (et non les bornes) ; vérifier que les deux totaux marginaux valent 40 ; dans le commentaire, citer la valeur de $r$, son signe et le seuil de comparaison ; l'estimation en 6. est une interpolation ($14\in[2{,}5\,;\,17{,}5]$), acceptable mais à nuancer vu $|r|<0{,}8$.
\end{corrige}

\begin{corrige}[Exercice 11 — Type Bac : changement de variable, croissance d'une population]
\begin{enumerate}
\item Les valeurs $z_i=\ln(y_i)$ arrondies à $10^{-3}$ sont : $2{,}996$ ; $3{,}186$ ; $3{,}378$ ; $3{,}570$ ; $3{,}761$ ; $3{,}953$.
\item Le nuage $N_i(x_i\,;\,z_i)$ est \textbf{quasi rectiligne}, croissant : la transformation $z=\ln y$ a linéarisé la relation. Un ajustement affine de $z$ en $x$ s'impose.
\item \begin{enumerate}
\item $\bar{x}=\dfrac{21}{6}=3{,}5$ ; $\bar{z}=\dfrac{2{,}996+3{,}186+3{,}378+3{,}570+3{,}761+3{,}953}{6}=\dfrac{20{,}844}{6}=3{,}474$.
$V(X)=\dfrac{91}{6}-12{,}25=\dfrac{35}{12}\approx2{,}917$.
$\overline{xz}=\dfrac{76{,}305}{6}\approx12{,}718$ ; $\operatorname{cov}(X,Z)=12{,}718-3{,}5\times3{,}474=12{,}718-12{,}159\approx0{,}559$.
$V(Z)=\dfrac{73{,}054}{6}-(3{,}474)^2\approx12{,}176-12{,}069\approx0{,}107$.
\[
r'=\frac{0{,}559}{\sqrt{2{,}917\times0{,}107}}=\frac{0{,}559}{\sqrt{0{,}312}}\approx\frac{0{,}559}{0{,}559}\approx0{,}999.
\]
$r'\approx1$ : l'ajustement affine de $z$ en $x$ est \textbf{excellent}, donc pertinent.
\item $a=\dfrac{0{,}559}{2{,}917}\approx0{,}191$ ; $b=3{,}474-0{,}191\times3{,}5\approx3{,}474-0{,}669\approx2{,}81$.
\[
z=0{,}191x+2{,}81.
\]
\end{enumerate}
\item $\ln y=0{,}191x+2{,}81\iff y=\mathrm{e}^{2{,}81}\,\mathrm{e}^{0{,}191x}$. Or $\mathrm{e}^{2{,}81}\approx16{,}6\approx17$.
\[
y=17\,\mathrm{e}^{0{,}191x}\quad\text{(en milliers d'habitants).}
\]
\item Rang 11 : $y=17\,\mathrm{e}^{0{,}191\times11}=17\,\mathrm{e}^{2{,}101}\approx17\times8{,}17\approx139$ milliers.
Population estimée : \textbf{environ 139 000 habitants} (extrapolation : résultat à manier avec prudence).
\item $17\,\mathrm{e}^{0{,}191x}>100\iff \mathrm{e}^{0{,}191x}>\dfrac{100}{17}\approx5{,}882\iff0{,}191x>\ln(5{,}882)\approx1{,}772\iff x>\dfrac{1{,}772}{0{,}191}\approx9{,}3$.
La population dépassera 100 000 habitants à partir du \textbf{rang 10}, c'est-à-dire dans 4 ans.
\end{enumerate}
\textbf{Barème indicatif (sur 10) :} tableau 0,5 pt ; nuage et constatation 1 pt ; covariance, variances et $r'$ 2,5 pts ; conclusion 0,5 pt ; équation de la droite 1,5 pt ; modèle exponentiel 1,5 pt ; estimation 1 pt ; inéquation 1,5 pt.

\textbf{Points de vigilance :} travailler sur le couple $(x\,;\,z)$ et non $(x\,;\,y)$ dans toute la question 3 ; justifier la pertinence par $r'$ (et non par l'allure seule) ; pour passer de $z$ à $y$, composer par l'exponentielle ($A=\mathrm{e}^{b}$, et non $A=b$) ; dans l'inéquation, la fonction $\ln$ étant croissante, le sens de l'inégalité est conservé ; arrondir la population au millier comme demandé.
\end{corrige}

\newpage

\coursec{Fiche bilan}

\coursec{Séries statistiques à deux caractères}

\begin{definition}[Série statistique à deux caractères]
Soit une population de taille $n$ sur laquelle on observe simultanément deux caractères quantitatifs $x$ et $y$. On obtient une \cle{série statistique à deux caractères} (ou série bidimensionnelle) représentée par les couples $(x_i\,;\,y_i)$ pour $i=1,\dots,n$, avec des effectifs $n_i$ (souvent $n_i=1$ pour des données brutes).
\end{definition}

\begin{definition}[Tableau à double entrée et séries marginales]
Lorsque les effectifs sont importants, on regroupe les valeurs de $x$ en $p$ classes (ou valeurs distinctes) et celles de $y$ en $q$ classes. Le \cle{tableau à double entrée} croise ces classes : la case $(i,j)$ contient l'effectif $n_{ij}$ des individus ayant $x$ dans la classe $i$ et $y$ dans la classe $j$.

Les \cle{séries marginales} sont les séries de $x$ seul et de $y$ seul, obtenues en sommant les effectifs par lignes et par colonnes :
\[
n_{i\cdot} = \sum_{j=1}^q n_{ij} \quad\text{(effectif marginal de la ligne $i$)}\qquad
n_{\cdot j} = \sum_{i=1}^p n_{ij} \quad\text{(effectif marginal de la colonne $j$)}
\]
On a bien $\sum_i n_{i\cdot} = \sum_j n_{\cdot j} = n$.
\end{definition}

\begin{exemplebox}[Consommation d'électricité et revenu — Douala]
Une enquête auprès de $20$ foyers à Douala donne le revenu mensuel $x$ (en milliers de FCFA) et la consommation d'électricité $y$ (en kWh). Le tableau à double entrée ci-dessous regroupe les données en classes.

\begin{center}
\begin{tabularx}{\linewidth}{|c|*{5}{C|}|C|}
\hline
\rowcolor{popblueL} \textbf{Revenu $x$ (kFCFA) $\backslash$ Conso $y$ (kWh)} & \textbf{[100;200[} & \textbf{[200;300[} & \textbf{[300;400[} & \textbf{[400;500[} & \textbf{[500;600[} & \textbf{Total $n_{i\cdot}$} \\
\hline
\textbf{[50;100[} & 1 & 2 & 1 & 0 & 0 & 4 \\
\hline
\textbf{[100;150[} & 0 & 1 & 3 & 2 & 0 & 6 \\
\hline
\textbf{[150;200[} & 0 & 0 & 2 & 3 & 1 & 6 \\
\hline
\textbf{[200;250[} & 0 & 0 & 0 & 2 & 2 & 4 \\
\hline
\rowcolor{popblueL} \textbf{Total $n_{\cdot j}$} & 1 & 3 & 6 & 7 & 3 & \textbf{20} \\
\hline
\end{tabularx}
\end{center}
\emph{Série marginale de $x$ :} effectifs $4, 6, 6, 4$. \emph{Série marginale de $y$ :} effectifs $1, 3, 6, 7, 3$.
\end{exemplebox}

\begin{methode}[Lire un tableau à double entrée]
\begin{enumerate}
  \item Identifier les classes (ou valeurs) des deux caractères en lignes et en colonnes.
  \item L'effectif $n_{ij}$ à l'intersection correspond aux individus vérifiant \emph{les deux} conditions simultanément.
  \item Les totaux des lignes (série marginale de $x$) et des colonnes (série marginale de $y$) permettent de retrouver les distributions de chaque caractère pris séparément.
  \item Pour calculer une moyenne marginale, utiliser les centres des classes (ou les valeurs exactes) et les effectifs marginaux.
\end{enumerate}
\end{methode}

\coursec{Nuage de points et point moyen}

\begin{definition}[Nuage de points]
Dans un repère orthonormé $(O;\vec{\imath},\vec{\jmath})$, on représente chaque individu par le point $M_i(x_i\,;\,y_i)$. L'ensemble de ces $n$ points forme le \cle{nuage de points}. Le choix des échelles sur les axes doit permettre de visualiser l'ensemble du nuage sans écrasement excessif.
\end{definition}

\begin{propriete}[Point moyen (barycentre)]
Le \cle{point moyen} $G$ du nuage a pour coordonnées les moyennes arithmétiques des deux caractères :
\[
G(\bar{x}\,;\,\bar{y}) \quad\text{avec}\quad \bar{x}=\frac{1}{n}\sum_{i=1}^n x_i,\quad \bar{y}=\frac{1}{n}\sum_{i=1}^n y_i.
\]
Toute droite d'ajustement affine (Mayer ou moindres carrés) passe par $G$. C'est un point de contrôle indispensable.
\end{propriete}

\begin{exemplebox}[Nuage et point moyen pour l'exemple Douala]
On utilise les centres des classes : $x \in \{75; 125; 175; 225\}$, $y \in \{150; 250; 350; 450; 550\}$.
Calcul des moyennes marginales (pondérées par les effectifs marginaux) :
\[
\bar{x} = \frac{75\times4 + 125\times6 + 175\times6 + 225\times4}{20} = \frac{300+750+1050+900}{20} = \frac{3000}{20} = 150
\]
\[
\bar{y} = \frac{150\times1 + 250\times3 + 350\times6 + 450\times7 + 550\times3}{20} = \frac{150+750+2100+3150+1650}{20} = \frac{7800}{20} = 390
\]
Le point moyen est $G(150\,;\,390)$. Sur le graphique, il se place au « centre de gravité » du nuage.
\end{exemplebox}

\begin{popfigure}
\begin{tikzpicture}[scale=0.7]
\begin{axis}[
  width=\linewidth, height=0.6\linewidth,
  xlabel={Revenu $x$ (kFCFA)}, ylabel={Conso $y$ (kWh)},
  xmin=50, xmax=250, ymin=100, ymax=600,
  grid=major, grid style={popline},
  scatter/classes={a={mark=*,popblue}, b={mark=*,poporange}, c={mark=*,popgreen}},
  legend style={at={(0.02,0.98)}, anchor=north west, fill=popcream}
]
\addplot[scatter,only marks,scatter src=explicit symbolic] coordinates {
(75,150) [a] (75,250) [a] (75,250) [a] (75,350) [a]
(125,250) [b] (125,350) [b] (125,350) [b] (125,350) [b] (125,450) [b] (125,450) [b]
(175,350) [c] (175,350) [c] (175,450) [c] (175,450) [c] (175,450) [c] (175,550) [c]
(225,450) [a] (225,450) [a] (225,550) [a] (225,550) [a]
};
\addplot[mark=*,mark size=3,popred] coordinates {(150,390)};
\addlegendentry{Nuage de points}
\addlegendentry{Point moyen $G(150;390)$}
\end{axis}
\end{tikzpicture}
\legende{Nuage de points de l'enquête Douala et position du point moyen $G$.}
\end{popfigure}

\coursec{Ajustement affine par la méthode de Mayer}

\begin{definition}[Principe de l'ajustement affine]
Un \cle{ajustement affine} consiste à remplacer le nuage de points par une droite d'équation $y = ax + b$ qui « suit » au mieux la tendance globale. Cette droite permet de résumer la relation entre $x$ et $y$ et de faire des prévisions.
\end{definition}

\begin{methode}[Méthode de Mayer (ou des deux points moyens)]
\begin{enumerate}
  \item Ordonner les $n$ points du nuage par valeurs croissantes de $x$ (en cas d'égalité, ordonner par $y$ croissant).
  \item Partager la liste ordonnée en \textbf{deux sous-nuages} d'effectifs égaux (ou différant au maximum de $1$ si $n$ est impair). Noter $n_1$ et $n_2$ leurs effectifs ($n_1+n_2=n$).
  \item Calculer les points moyens de chaque sous-nuage :
  \[
  G_1(\bar{x}_1\,;\,\bar{y}_1),\quad \bar{x}_1=\frac{1}{n_1}\sum_{\text{sous-nuage 1}} x_i,\quad \bar{y}_1=\frac{1}{n_1}\sum_{\text{sous-nuage 1}} y_i
  \]
  et de même pour $G_2(\bar{x}_2\,;\,\bar{y}_2)$.
  \item La \cle{droite de Mayer} est la droite $(G_1G_2)$. Son équation s'obtient par :
  \[
  a = \frac{\bar{y}_2-\bar{y}_1}{\bar{x}_2-\bar{x}_1} \quad (\text{si } \bar{x}_1\neq\bar{x}_2),\qquad b = \bar{y}_1 - a\bar{x}_1.
  \]
  \item \textbf{Vérification :} La droite passe par $G(\bar{x}\,;\,\bar{y})$ car $G$ est le barycentre de $G_1$ (poids $n_1$) et $G_2$ (poids $n_2$).
\end{enumerate}
\end{methode}

\begin{attention}[Pièges fréquents — Méthode de Mayer]
\begin{itemize}
  \item Ne pas partager « au milieu de l'axe des $x$ » (par la médiane des valeurs de $x$), mais \textbf{par l'effectif} : les deux sous-nuages doivent contenir le même nombre de points.
  \item Si $n$ est impair, le point central (médiane) est \textbf{exclu} des deux sous-nuages (ou attribué à l'un des deux de façon à équilibrer les effectifs, la méthode standard l'exclut).
  \item La méthode de Mayer est sensible aux valeurs extrêmes (points aberrants) car elle ne utilise que $G_1$ et $G_2$.
\end{itemize}
\end{attention}

\begin{exoresolu}[Droite de Mayer — Enquête Douala]
\emph{Énoncé :} Déterminer l'équation de la droite de Mayer pour les $20$ foyers (données brutes triées par $x$ croissant).
\tcblower
\emph{Solution :} $n=20$ pair $\Rightarrow$ deux sous-nuages de $10$ points chacun.
On calcule les moyennes sur les $10$ premiers et les $10$ derniers points (détail omis ici, supposons $G_1(110; 300)$ et $G_2(190; 480)$).
\[
a = \frac{480-300}{190-110} = \frac{180}{80} = 2,25,\quad b = 300 - 2,25\times 110 = 300 - 247,5 = 52,5
\]
Droite de Mayer : $y = 2,25x + 52,5$.
Vérification : pour $x=150$, $y = 2,25\times150 + 52,5 = 337,5 + 52,5 = 390 = \bar{y}$. $G$ appartient bien à la droite.
\end{exoresolu}

\coursec{Covariance et moindres carrés}

\begin{definition}[Covariance]
La \cle{covariance} des deux caractères $x$ et $y$ mesure leur variation conjointe :
\[
\mathrm{cov}(x,y) = \frac{1}{n}\sum_{i=1}^n (x_i-\bar{x})(y_i-\bar{y}) = \overline{xy} - \bar{x}\bar{y}
\]
où $\overline{xy} = \frac{1}{n}\sum x_i y_i$ (moyenne des produits).
\end{definition}

\begin{propriete}[Propriétés de la covariance]
\begin{itemize}
  \item $\mathrm{cov}(x,x) = V(x)$ (variance de $x$).
  \item $\mathrm{cov}(x,y) = \mathrm{cov}(y,x)$ (symétrie).
  \item Bilinéarité : $\mathrm{cov}(ax+b,\,cy+d) = ac\,\mathrm{cov}(x,y)$ pour $a,b,c,d\in\mathbb{R}$.
  \item \textbf{Interprétation du signe :} $\mathrm{cov}(x,y) > 0$ signifie que $x$ et $y$ ont tendance à varier dans le même sens (corrélation positive) ; $<0$ dans le sens inverse ; $=0$ absence de liaison linéaire (mais pas nécessairement indépendance).
\end{itemize}
\end{propriete}

\begin{experience}[Démonstration guidée — Méthode des moindres carrés (Droite de régression $y$ en $x$)]
On cherche la droite $y=ax+b$ qui minimise la somme des carrés des écarts verticaux $S(a,b) = \sum_{i=1}^n (y_i - (ax_i+b))^2$.
\begin{enumerate}
  \item Exprimer $S(a,b)$ en fonction de $a,b$ et des sommes $\sum x_i, \sum y_i, \sum x_i^2, \sum x_i y_i$.
  \item Annuler les dérivées partielles $\frac{\partial S}{\partial a}=0$ et $\frac{\partial S}{\partial b}=0$ (système normal).
  \item Montrer que le système donne :
  \[
  \begin{cases}
  a\sum x_i^2 + b\sum x_i = \sum x_i y_i \\
  a\sum x_i + nb = \sum y_i
  \end{cases}
  \iff
  \begin{cases}
  a(\overline{x^2}-\bar{x}^2) = \overline{xy}-\bar{x}\bar{y} \\
  b = \bar{y} - a\bar{x}
  \end{cases}
  \]
  \item Conclure : $a = \dfrac{\mathrm{cov}(x,y)}{V(x)}$, $b = \bar{y} - a\bar{x}$.
\end{enumerate}
\end{experience}

\begin{propriete}[Droites de régression par les moindres carrés]
\begin{itemize}
  \item \cle{Droite de régression de $y$ en $x$} (notée $D_{y/x}$) : $y = ax + b$ avec
  \[
  a = \frac{\mathrm{cov}(x,y)}{V(x)},\qquad b = \bar{y} - a\bar{x}.
  \]
  Elle minimise la somme des carrés des distances \textbf{verticales} aux points. On l'utilise pour \textbf{prévoir $y$ connaissant $x$}.
  \item \cle{Droite de régression de $x$ en $y$} (notée $D_{x/y}$) : $x = a'y + b'$ avec
  \[
  a' = \frac{\mathrm{cov}(x,y)}{V(y)},\qquad b' = \bar{x} - a'\bar{y}.
  \]
  Elle minimise la somme des carrés des distances \textbf{horizontales}. On l'utilise pour \textbf{prévoir $x$ connaissant $y$}.
\end{itemize}
Les deux droites passent par $G(\bar{x}\,;\,\bar{y})$ et s'y coupent. Elles sont distinctes sauf si $|r|=1$ (points alignés).
\end{propriete}

\begin{exemplebox}[Calcul des droites de régression — Enquête Douala]
On dispose des sommes (calculées à partir du tableau à double entrée avec centres de classes) :
$\sum x_i = 3000$, $\sum y_i = 7800$, $\sum x_i^2 = 487\,500$, $\sum y_i^2 = 3\,318\,000$, $\sum x_i y_i = 1\,252\,500$, $n=20$.
Moyennes : $\bar{x}=150$, $\bar{y}=390$.
$\overline{xy} = 62\,625$, $\overline{x^2} = 24\,375$, $\overline{y^2} = 165\,900$.
$V(x) = 24\,375 - 150^2 = 24\,375 - 22\,500 = 1\,875$.
$V(y) = 165\,900 - 390^2 = 165\,900 - 152\,100 = 13\,800$.
$\mathrm{cov}(x,y) = 62\,625 - 150\times390 = 62\,625 - 58\,500 = 4\,125$.

\emph{Droite $D_{y/x}$ :} $a = \frac{4\,125}{1\,875} = 2,2$ ; $b = 390 - 2,2\times150 = 390 - 330 = 60$.
$\Rightarrow D_{y/x} : y = 2,2x + 60$.

\emph{Droite $D_{x/y}$ :} $a' = \frac{4\,125}{13\,800} \approx 0,2989$ ; $b' = 150 - 0,2989\times390 \approx 150 - 116,57 = 33,43$.
$\Rightarrow D_{x/y} : x = 0,2989y + 33,43$ (ou $y \approx 3,345x - 111,8$ en tirant $y$ en fonction de $x$, mais ce n'est plus la régression de $y$ en $x$).
\end{exemplebox}

\coursec{Coefficient de corrélation linéaire}

\begin{definition}[Coefficient de corrélation linéaire (de Pearson)]
Le \cle{coefficient de corrélation linéaire} $r$ (ou $\rho$) est défini par :
\[
r = \frac{\mathrm{cov}(x,y)}{\sigma_x\,\sigma_y} = \frac{\mathrm{cov}(x,y)}{\sqrt{V(x)V(y)}}
\]
où $\sigma_x = \sqrt{V(x)}$, $\sigma_y = \sqrt{V(y)}$ sont les écarts-types.
\end{definition}

\begin{propriete}[Propriétés fondamentales de $r$]
\begin{itemize}
  \item $-1 \le r \le 1$ (inégalité de Cauchy-Schwarz).
  \item $r = 1 \iff$ tous les points sont alignés sur une droite de pente \textbf{positive} ($y=ax+b, a>0$).
  \item $r = -1 \iff$ tous les points sont alignés sur une droite de pente \textbf{négative} ($a<0$).
  \item $r = 0 \iff \mathrm{cov}(x,y)=0$ (absence de corrélation \textbf{linéaire} ; il peut exister une relation non-linéaire).
  \item $r^2 = a \times a'$ (produit des pentes des deux droites de régression). $r^2$ est le \cle{coefficient de détermination} : proportion de la variance de $y$ expliquée par l'ajustement affine sur $x$.
\end{itemize}
\end{propriete}

\begin{aretenir}[Interprétation de $|r|$ — Seuil de pertinence au Bac]
\begin{center}
\begin{tabularx}{\linewidth}{|c|X|}
\hline
\rowcolor{popblueL} \textbf{Valeur de $|r|$} & \textbf{Interprétation} \\
\hline
$|r| \ge 0,87$ & Corrélation linéaire \textbf{forte} — L'ajustement affine est \textbf{pertinent} (on peut prévoir). \\
\hline
$0,5 \le |r| < 0,87$ & Corrélation linéaire \textbf{modérée} — Ajustement possible mais prévisions peu fiables. \\
\hline
$|r| < 0,5$ & Corrélation linéaire \textbf{faible ou nulle} — Ajustement affine \textbf{non pertinent} (ne pas prévoir). \\
\hline
\end{tabularx}
\end{center}
\textbf{Attention :} Le seuil $0,87$ (correspondant à $r^2 \approx 0,75$) est la convention usuelle dans l'enseignement secondaire camerounais. Vérifiez toujours l'énoncé du sujet.
\end{aretenir}

\begin{exemplebox}[Calcul et interprétation de $r$ — Enquête Douala]
$\sigma_x = \sqrt{1\,875} \approx 43,30$ ; $\sigma_y = \sqrt{13\,800} \approx 117,47$.
$r = \frac{4\,125}{43,30 \times 117,47} \approx \frac{4\,125}{5\,086} \approx 0,811$.
$r^2 \approx 0,658$.
\emph{Conclusion :} $r > 0$ (corrélation positive : plus le revenu augmente, plus la consommation augmente). $|r| = 0,811 < 0,87$ : la corrélation linéaire est modérée, l'ajustement affine est \textbf{peu pertinent} pour des prévisions précises. Environ $66\%$ de la variabilité de la consommation est expliquée par le revenu.
\end{exemplebox}

\coursec{Prévisions et synthèse méthodologique}

\begin{methode}[Réaliser une prévision par ajustement affine]
\begin{enumerate}
  \item \textbf{Vérifier la pertinence :} Calculer $r$. Si $|r| < 0,87$ (seuil Bac), écrire : « L'ajustement affine n'est pas pertinent, la prévision n'est pas fiable » et s'arrêter (ou préciser « à titre indicatif » si le sujet l'impose).
  \item \textbf{Choisir la bonne droite :}
  \begin{itemize}
    \item Pour prévoir $y$ à partir d'une valeur $x_0$ : utiliser $D_{y/x}$ ($y_0 = ax_0 + b$).
    \item Pour prévoir $x$ à partir d'une valeur $y_0$ : utiliser $D_{x/y}$ ($x_0 = a'y_0 + b'$).
  \end{itemize}
  \item \textbf{Calculer la prévision} en remplaçant dans l'équation de la droite choisie.
  \item \textbf{Discuter le résultat :} Interpolation ($x_0$ dans l'intervalle des $x$ observés) $\rightarrow$ prévision assez fiable. Extrapolation ($x_0$ hors intervalle) $\rightarrow$ prévision très incertaine, à formuler avec prudence.
\end{enumerate}
\end{methode}

\begin{attention}[Erreurs classiques en prévision]
\begin{itemize}
  \item Utiliser $D_{y/x}$ pour prévoir $x$ (ou l'inverse). Les deux droites sont \textbf{différentes} (sauf $|r|=1$).
  \item Oublier de vérifier $|r| \ge 0,87$ avant de conclure « la prévision est fiable ».
  \item Faire une extrapolation lointaine sans avertir sur le risque (ex: prévoir la consommation pour un revenu de 1\,000 kFCFA alors que l'échantillon s'arrête à 250).
  \item Arrondir les coefficients $a,b$ trop tôt (garder 3 à 4 décimales pour les calculs intermédiaires, arrondir le résultat final selon le contexte).
\end{itemize}
\end{attention}

\begin{savaistu}[Le seuil 0,87 et le coefficient de détermination]
Le seuil $|r|\ge 0,87$ correspond à $r^2 \ge 0,7569 \approx 75\%$. Cela signifie que la droite de régression explique au moins $75\%$ de la variance du caractère prédit. En économie ou en biologie (ex: rendement maïs vs engrais à Dschang), on accepte parfois des seuils plus bas ($0,7$ ou $0,8$) selon la variabilité naturelle du phénomène. Au Bac, \textbf{appliquez le seuil indiqué dans l'énoncé} (souvent $0,87$ ou $0,9$).
\end{savaistu}

\begin{exoresolu}[Prévision — Enquête Douala]
\emph{Énoncé :} Un foyer a un revenu de $180$ kFCFA. Estimer sa consommation. Un autre foyer consomme $300$ kWh. Estimer son revenu.
\tcblower
\emph{Solution :}
$r \approx 0,811 < 0,87$ $\Rightarrow$ \textbf{L'ajustement n'est pas pertinent au seuil $0,87$.} Les prévisions ci-dessous sont données à titre indicatif seulement.

1. Prévoir $y$ pour $x_0=180$ : utiliser $D_{y/x} : y = 2,2x + 60$.
$y_0 = 2,2 \times 180 + 60 = 396 + 60 = 456$ kWh.
(Interpolation : $180 \in [50;250]$, prévision acceptable mais incertaine car $r<0,87$).

2. Prévoir $x$ pour $y_0=300$ : utiliser $D_{x/y} : x = 0,2989y + 33,43$.
$x_0 = 0,2989 \times 300 + 33,43 \approx 89,67 + 33,43 = 123,1$ kFCFA.
(Interpolation : $300 \in [100;600]$, même réserve).
\end{exoresolu}

\begin{aretenir}[Synthèse : La « boîte à outils » pour le Bac]
\begin{itemize}
  \item \textbf{Tableau double entrée} $\rightarrow$ Séries marginales (totaux lignes/colonnes).
  \item \textbf{Nuage de points} $\rightarrow$ Repère adapté, point moyen $G(\bar{x},\bar{y})$.
  \item \textbf{Mayer} : Trier par $x$, couper en 2 (effectifs), $G_1,G_2$, droite $(G_1G_2)$. Vérifier $G \in (G_1G_2)$.
  \item \textbf{Moindres carrés} : $\mathrm{cov}(x,y)=\overline{xy}-\bar{x}\bar{y}$, $V(x)=\overline{x^2}-\bar{x}^2$.
  $D_{y/x}: a=\frac{\mathrm{cov}}{V(x)},\ b=\bar{y}-a\bar{x}$.
  $D_{x/y}: a'=\frac{\mathrm{cov}}{V(y)},\ b'=\bar{x}-a'\bar{y}$.
  \item \textbf{Corrélation} : $r=\frac{\mathrm{cov}}{\sigma_x\sigma_y}$. Vérifier $-1\le r\le 1$. Interpréter signe et $|r|$ vs seuil.
  \item \textbf{Prévision} : Choisir la bonne droite ($D_{y/x}$ pour $y$, $D_{x/y}$ pour $x$). Dire si interpolation/extrapolation. Conclure avec prudence.
\end{itemize}
\end{aretenir}

\begin{exercice}[Entraînement — Ventes de mangues et pluviométrie]
Un producteur de mangues à Njombé note sur 12 mois la pluviométrie $x$ (mm) et la production $y$ (tonnes) :
$(120; 4,2),\ (150; 5,1),\ (180; 5,8),\ (200; 6,5),\ (220; 7,0),\ (250; 7,8),\ (280; 8,5),\ (300; 9,0),\ (320; 9,5),\ (350; 10,2),\ (380; 10,8),\ (400; 11,0)$.
\begin{enumerate}
  \item Représenter le nuage de points et le point moyen $G$.
  \item Déterminer la droite de Mayer.
  \item Calculer la covariance, les variances, le coefficient $r$.
  \item Déterminer les deux droites de régression.
  \item Prévoir la production pour une pluviométrie de $330$ mm. Commenter.
\end{enumerate}
\end{exercice}

\begin{corrige}[Exercice — Ventes de mangues]
\begin{enumerate}
  \item $n=12$. $\sum x = 3150 \Rightarrow \bar{x} = 262,5$. $\sum y = 95,4 \Rightarrow \bar{y} = 7,95$. $G(262,5; 7,95)$.
  \item Mayer : $n$ pair, 2 sous-nuages de 6 points.
  Sous-nuage 1 (6 premiers) : $\bar{x}_1 = 187,5$, $\bar{y}_1 = 6,05 \Rightarrow G_1(187,5; 6,05)$.
  Sous-nuage 2 (6 derniers) : $\bar{x}_2 = 337,5$, $\bar{y}_2 = 9,85 \Rightarrow G_2(337,5; 9,85)$.
  $a = \frac{9,85-6,05}{337,5-187,5} = \frac{3,8}{150} \approx 0,02533$. $b = 6,05 - 0,02533\times187,5 \approx 1,30$.
  Mayer : $y = 0,0253x + 1,30$. (Vérif : $0,0253\times262,5+1,30 \approx 7,95 = \bar{y}$).
  \item $\sum x^2 = 917\,700 \Rightarrow \overline{x^2}=76\,475 \Rightarrow V(x)=76\,475-262,5^2=7\,593,75$.
  $\sum y^2 = 798,48 \Rightarrow \overline{y^2}=66,54 \Rightarrow V(y)=66,54-7,95^2=3,3375$.
  $\sum xy = 27\,399 \Rightarrow \overline{xy}=2\,283,25 \Rightarrow \mathrm{cov}=2\,283,25 - 262,5\times7,95 = 200,5625$.
  $r = \frac{200,5625}{\sqrt{7\,593,75 \times 3,3375}} \approx \frac{200,56}{159,1} \approx 0,996$.
  \item $D_{y/x}$ : $a = \frac{200,56}{7\,593,75} \approx 0,02641$, $b = 7,95 - 0,02641\times262,5 \approx 1,02$. $y = 0,0264x + 1,02$.
  $D_{x/y}$ : $a' = \frac{200,56}{3,3375} \approx 60,10$, $b' = 262,5 - 60,10\times7,95 \approx -215,3$. $x = 60,10y - 215,3$.
  \item $r=0,996 > 0,87$ : ajustement \textbf{très pertinent}.
  Prévoir $y$ pour $x_0=330$ (interpolation) : utiliser $D_{y/x}$.
  $y_0 = 0,0264\times330 + 1,02 \approx 9,73$ tonnes.
\end{enumerate}
\end{corrige}

\findecours

\end{document}
